%%% Copyright (C) 2017-2022 Vincent Goulet
%%%
%%% Ce fichier et tous les fichiers .tex ou .Rnw dont la racine est
%%% mentionnée dans les commandes \include ci-dessous font partie du
%%% projet «Programmer avec R».
%%% https://gitlab.com/vigou3/programmer-avec-r
%%%
%%% Cette création est mise à disposition sous licence
%%% Attribution-Partage dans les mêmes conditions 4.0
%%% International de Creative Commons.
%%% https://creativecommons.org/licenses/by-sa/4.0/

\documentclass[letterpaper,11pt,x11names,english,french]{memoir}
  \usepackage{natbib,url}
  \usepackage{babel}
  \usepackage[autolanguage]{numprint}
  \usepackage{amsmath,amsthm}
  \usepackage[noae]{Sweave}
  \usepackage{graphicx}
  \usepackage{pict2e}                  % dessin de la clôture
  \usepackage{calc}                    % \widthof
  \usepackage{currfile}                % noms des fichiers de script
  \usepackage{actuarialangle}          % \angl et al.
  \usepackage{framed}                  % env. snugshade*, oframed
  \usepackage[absolute]{textpos}       % éléments des pages de titre
  \usepackage[shortlabels]{enumitem}   % configuration listes
  \usepackage{multicol}                % environnement multicols
  \usepackage{relsize}                 % \smaller et al.
  \usepackage{manfnt}                  % \mantriangleright (puce)
  \usepackage{metalogo}                % \XeLaTeX logo
  \usepackage{applekeys}               % touches Mac
  \usepackage{dirtree}                 % arborescence de fichiers
  \usepackage{fontawesome5}            % icônes \fa*
  \usepackage{awesomebox}              % boites info, important, etc.
  \usepackage{ullrcorners}             % expressions régulières
  \usepackage{soul}                    % surbrillance (expressions régulières)
  \usepackage{answers}                 % exercices et solutions
  \usepackage{listings}                % code informatique
  \usepackage{refcount}                % numéros de ligne en surbrillance
  \usepackage{imakeidx}                % configuration de l'index

  %%% =============================
  %%%  Informations de publication
  %%% =============================
  \title{Programmer avec R}
  \author{Vincent Goulet}
  \renewcommand{\year}{2022}
  \renewcommand{\month}{08}
  \newcommand{\reposurl}{https://gitlab.com/vigou3/programmer-avec-r}

  %%% ===================
  %%%  Style du document
  %%% ===================

  %% Polices de caractères
  \usepackage{fontspec}
  \usepackage{unicode-math}
  \defaultfontfeatures
  {
    Scale = 0.92
  }
  \setmainfont{Lucida Bright OT}
  [
    Ligatures = TeX,
    Numbers = OldStyle
  ]
  \setmathfont{Lucida Bright Math OT}
  \setmonofont{Lucida Grande Mono DK}
  \setsansfont{FiraSans}
  [
    Extension = .otf,
    UprightFont = *-Book,
    ItalicFont = *-BookItalic,
    BoldFont = *-SemiBold,
    BoldItalicFont = *-SemiBoldItalic,
    Scale = 0.98,
    Numbers = OldStyle
  ]
  \newfontfamily\fullcaps{FiraSans}
  [
    Extension = .otf,
    UprightFont = *-Book,
    Scale = 0.98,
    Numbers = Uppercase
  ]
  \usepackage[babel=true]{microtype}
  \usepackage{icomma}

  %% Couleurs
  \usepackage{xcolor}
  \definecolor{comments}{rgb}{0.5,0.55,0.6} % commentaires
  \definecolor{link}{rgb}{0,0.4,0.6}        % liens internes
  \definecolor{url}{rgb}{0.6,0,0}           % liens externes
  \definecolor{citation}{rgb}{0,0.5,0}      % citations
  \colorlet{regex}{Snow4}                   % coins des regex
  \colorlet{regexmatch}{Snow3}              % correspondances regex
  \colorlet{codebg}{LightYellow1}           % fond code R
  \colorlet{indexedface}{Orchid2}     % cellule indexée (face)
  \colorlet{indexedtop}{Orchid3}      % cellule indexée (dessus)
  \colorlet{indexedside}{Orchid4}     % cellule indexée (côté)
  \colorlet{mindexedface}{MediumOrchid2} % cellule indexée médium (face)
  \colorlet{mindexedtop}{MediumOrchid3}  % cellule indexée médium (dessus)
  \colorlet{mindexedside}{MediumOrchid4} % cellule indexée médium (côté)
  \colorlet{dindexedface}{DarkOrchid2} % cellule indexée foncé (face)
  \colorlet{dindexedtop}{DarkOrchid3}  % cellule indexée foncé (dessus)
  \colorlet{dindexedside}{DarkOrchid4} % cellule indexée foncé (côté)
  \colorlet{prob}{orange}                   % encadrés «problème»
  \colorlet{lineno}{gray}                   % numéros de lignes
  \colorlet{TFFrameColor}{black}            % encadrés (cadre)
  \colorlet{TFTitleColor}{white}            % encadrés (titre)
  \definecolor{rouge}{rgb}{0.90,0,0.1}      % rouge bandeau UL
  \definecolor{or}{rgb}{1,0.8,0}            % or bandeau UL

  %% Hyperliens
  \usepackage[backref=page]{hyperref}
  \hypersetup{%
    pdfauthor = \theauthor,
    pdftitle = \thetitle,
    colorlinks = true,
    linktocpage = true,
    urlcolor = {url},
    linkcolor = {link},
    citecolor = {citation},
    pdfpagemode = {UseOutlines},
    pdfstartview = {Fit}}
  \renewcommand*{\backrefalt}[4]{%
    \ifcase #1 %
      Aucune citation.%
    \or
      Cité à la page #2.%
    \else
      Cité aux pages #2.%
    \fi
  }
  \setlength{\XeTeXLinkMargin}{1pt}

  %% Affichage de la table des matières du PDF
  \usepackage{bookmark}
  \bookmarksetup{%
    open = true,
    depth = 3,
    numbered = true}

  %% Étiquettes de \autoref (redéfinitions compatibles avec babel).
  %% Attention! Les % à la fin des lignes sont importants sinon des
  %% blancs apparaissent dès que la commande \selectlanguage est
  %% utilisée... comme dans la bibliographie, par exemple.
  \addto\extrasfrench{%
    \def\algorithmeautorefname{algorithme}%
    \def\algorithmebisautorefname{algorithme}%
    \def\appendixautorefname{annexe}%
    \def\figureautorefname{figure}%
    \def\exerciceautorefname{exercice}%
    \def\subfigureautorefname{figure}%
    \def\subsectionautorefname{section}%
    \def\subtableautorefname{tableau}%
    \def\tableautorefname{tableau}%
  }

  %% Table des matières et al. (inspiré de classicthesis.sty)
  \renewcommand{\cftchapterleader}{\hspace{1.5em}}
  \renewcommand{\cftchapterafterpnum}{\cftparfillskip}
  \renewcommand{\cftsectionleader}{\hspace{1.5em}}
  \renewcommand{\cftsectionafterpnum}{\cftparfillskip}
  \renewcommand{\cfttableleader}{\hspace{1.5em}}
  \renewcommand{\cfttableafterpnum}{\cftparfillskip}
  \renewcommand{\cftfigureleader}{\hspace{1.5em}}
  \renewcommand{\cftfigureafterpnum}{\cftparfillskip}
  \setlength{\cftbeforechapterskip}{1.0em plus 0.5em minus 0.5em}
  \setlength{\cftfigurenumwidth}{3.2em}

  %% Création d'une «liste des vidéos» similaires à la
  %% liste des tableaux et à la liste des figures.
  \newcommand{\listvideosname}{Liste des vidéos}
  \newlistof{listofvideos}{lov}{\listvideosname}
  \renewcommand{\afterlovtitle}{%
    \afterchaptertitle{\normalfont\normalsize Le numéro indiqué à
      gauche est celui de la section dans laquelle se trouve le bloc
      signalétique.\label{listofvideos}}\vspace{10pt}}

  %% Titres des chapitres
  \chapterstyle{hangnum}
  \renewcommand{\chaptitlefont}{\normalfont\Huge\sffamily\bfseries\raggedright}

  %% Marges, entêtes et pieds de page
  \setlength{\marginparsep}{7mm}
  \setlength{\marginparwidth}{13mm}
  \setlength{\headwidth}{\textwidth}
  \addtolength{\headwidth}{\marginparsep}
  \addtolength{\headwidth}{\marginparwidth}

  %% Titres des sections et sous-sections
  \setsecheadstyle{\normalfont\Large\sffamily\bfseries\raggedright}
  \setsubsecheadstyle{\normalfont\large\sffamily\bfseries\raggedright}
  \maxsecnumdepth{subsection}
  \setsecnumdepth{subsection}

  %% Listes. Paramétrage avec enumitem.
  \setlist[enumerate]{leftmargin=*,align=left}
  \setlist[enumerate,2]{label=\alph*),labelsep=*,leftmargin=1.5em}
  \setlist[enumerate,3]{label=\roman*),labelsep=*,leftmargin=1.5em,align=right}
  \setlist[itemize]{leftmargin=*,align=left}
  \setlist[trivlist]{nosep}

  %% Options de babel
  \frenchbsetup{CompactItemize=false,%
    ThinSpaceInFrenchNumbers=true,
    ItemLabeli=\mantriangleright,
    ItemLabelii=\textendash,
    og=«, fg=»}
  \def\frenchtablename{{\scshape Tab.}}
  \def\frenchfigurename{{\scshape Fig.}}
  \def\frenchlistfigurename{Liste des figures}

  %%% =========================
  %%%  Sections de code source
  %%% =========================

  %% Syntaxe de R avec ajouts de quelques mots clés.
  \lstloadlanguages{R,sh}
  \lstdefinelanguage{Renhanced}[]{R}{%
    otherkeywords = {!,!=,~,$,*,\&,\%/\%,\%*\%,\%\%,<-,<<-},%$
    morekeywords={colMeans,colSums,head,is.na,is.null,mapply,ms,na.rm,%
      nlmin,rep_len,replicate,row.names,rowMeans,rowSums,seq_len,%
      seq_along,sys.time,system.time,tail,which.max,which.min,letters,%
      LETTERS,Trig},
    deletekeywords={c,start,q,_},
    alsoletter={._\%},
    alsoother={:\$},
    moredelim=[il]{\#-*-}}
  \lstdefinelanguage{shenhanced}[]{sh}{%
    morekeywords={grep},
    moredelim=[il]{\#-*-}}

  %% Mise en forme du code source.
  %%
  %% Les numéros de lignes sont des hyperliens vers le point dans le
  %% document où l'on y fait référence dans une boite \gotorbox (voir
  %% plus loin).
  %%
  %% Pour y parvenir, j'utilise deux étiquettes pour une ligne: une
  %% basée sur un "nom" utilisé dans la rédaction, et une autre
  %% générée automatiquement à partir du numéro de chapitre et du
  %% numéro de ligne.
  %%
  %% Solution basée sur https://tex.stackexchange.com/q/191771
  \lstset{%
    language=Renhanced,
    extendedchars=true,
    basicstyle=\small\ttfamily\NoAutoSpacing,
    commentstyle=\color{comments}\slshape,
    keywordstyle=\mdseries,
    showstringspaces=false,
    numbers=left,
    numberstyle={%
      \color{lineno}\tiny\ttfamily%
      \ifnum\value{lstnumber}=\getrefnumber{code:\thechapter:\thelstnumber}%
        \renewcommand*\thelstnumber{\hyperlink{goto:\thechapter:\the\value{lstnumber}}{\bfseries\arabic{lstnumber}}}%
      \fi},
    firstnumber=\scriptfirstline,
    escapechar=`,
    index=[1][keywords],
    indexstyle=\indexcode}

  %% Commandes pour créer les références et liens vers des numéros de
  %% lignes.
  \makeatletter
  \newcommand\labelline[1]{%
    \def\@currentlabel{\thelstnumber}%
    \label{lst:#1}\label{code:\thechapter:\the\value{lstnumber}}}
  \makeatother
  \newcommand{\reflines}[1]{%
    \hypertarget{goto:\thechapter:\getrefnumber{lst:#1}}{\ref{lst:#1}}--%
    \hypertarget{goto:\thechapter:\getrefnumber{lst:#1:fin}}{\ref{lst:#1:fin}}}

  %% L'entête des fichiers de script n'est pas affiché dans le
  %% document.
  \def\scriptfirstline{11}      % nombre magique!

  %%% =========================
  %%%  Nouveaux environnements
  %%% =========================

  %% Environnements de type théorème
  \theoremstyle{definition}
  \newtheorem{algorithme}{Algorithme}[chapter]
  \newtheorem*{algorithme*}{Algorithme} % pour les solutions
  \newtheorem{algorithmebis}{Algorithme}[algorithme]
  \renewcommand{\thealgorithmebis}{\thealgorithme\alph{algorithmebis}}

  %% Redéfinition de l'environnement titled-frame de framed.sty avec
  %% deux modifications: épaisseur des filets réduite de 2pt à 1pt;
  %% "(suite)" plutôt que "(cont)" dans la barre de titre
  %% lorsque l'encadré se poursuit après un saut de page.
  \renewenvironment{titled-frame}[1]{%
    \def\FrameCommand{\fboxsep8pt\fboxrule1pt
      \TitleBarFrame{\textbf{#1}}}%
    \def\FirstFrameCommand{\fboxsep8pt\fboxrule1pt
      \TitleBarFrame[$\blacktriangleright$]{\textbf{#1}}}%
    \def\MidFrameCommand{\fboxsep8pt\fboxrule1pt
      \TitleBarFrame[$\blacktriangleright$]{\textbf{#1\ (suite)}}}%
    \def\LastFrameCommand{\fboxsep8pt\fboxrule1pt
      \TitleBarFrame{\textbf{#1\ (suite)}}}%
    \MakeFramed{\advance\hsize-16pt \FrameRestore}}%
  {\endMakeFramed}

  %% Liste d'objectifs au début des chapitres dans un encadré
  %% basé sur titled-frame, ci-dessus.
  \newenvironment{objectifs}{%
    \begin{titled-frame}{\rule[-7pt]{0pt}{20pt}\sffamily Objectifs du chapitre}
      \setlength{\parindent}{0pt}
      \small\sffamily
      \begin{itemize}[nosep]}%
      {\end{itemize}\end{titled-frame}}

  %% Problèmes (mises en situation) des chapitres: énoncé au début du
  %% chapitre; astuces en cours de chapitre; solution à la fin
  %% du chapitre.
  \newenvironment{prob-enonce}{%
    \begin{emphbox}[prob]{{\normalfont\faCogs}\; Énoncé du problème}}%
    {\end{emphbox}}
  \newenvironment{prob-astuce}{%
    \begin{emphbox}[prob]{{\normalfont\faBolt}\; Astuce}}%
    {\end{emphbox}}
  \newenvironment{prob-solution}{%
    \begin{emphbox}[prob]{{\normalfont\faLightbulbO}\; Solution du problème}}%
    {\end{emphbox}}

  %% Environnements de Sweave. Les environnements Sinput et Soutput
  %% utilisent Verbatim (de fancyvrb). On les réinitialise pour
  %% enlever la configuration par défaut de Sweave, puis on réduit
  %% l'écart entre les blocs Sinput et Soutput.
  \DefineVerbatimEnvironment{Sinput}{Verbatim}{}
  \DefineVerbatimEnvironment{Soutput}{Verbatim}{}
  \fvset{listparameters={\setlength{\topsep}{0pt}}}

  %% L'environnement Schunk est complètement redéfini en un hybride
  %% des environnements snugshade* et leftbar de framed.sty.
  \makeatletter
  \renewenvironment{Schunk}{%
    \setlength{\topsep}{1pt}
    \def\FrameCommand##1{\hskip\@totalleftmargin
       \vrule width 2pt\colorbox{codebg}{\hspace{3pt}##1}%
      % There is no \@totalrightmargin, so:
      \hskip-\linewidth \hskip-\@totalleftmargin \hskip\columnwidth}%
    \MakeFramed {\advance\hsize-\width
      \@totalleftmargin\z@ \linewidth\hsize
      \advance\labelsep\fboxsep
      \@setminipage}%
  }{\par\unskip\@minipagefalse\endMakeFramed}
  \makeatother

  %% Exercices et réponses
  \Newassociation{sol}{solution}{solutions}
  \newcounter{exercice}[chapter]
  \renewcommand{\theexercice}{\thechapter.\arabic{exercice}}
  \newenvironment{exercice}[1][]{%
    \begin{list}{}{%
        \refstepcounter{exercice}
        \ifthenelse{\equal{#1}{nosol}}{%
          \renewcommand{\makelabel}{\bfseries\theexercice}}{%
          \hypertarget{ex:\theexercice}{}
          \Writetofile{solutions}{\protect\hypertarget{sol:\theexercice}{}}
          \renewcommand{\makelabel}{%
            \bfseries\protect\hyperlink{sol:\theexercice}{\theexercice}}}
        \settowidth{\labelwidth}{\bfseries\theexercice}
        \setlength{\leftmargin}{\labelwidth}
        \addtolength{\leftmargin}{\labelsep}
        \setlist[enumerate,1]{label=\alph*),labelsep=*,leftmargin=1.5em}
        \setlist[enumerate,2]{label=\roman*),labelsep=0.5em,align=right}}
    \item}%
    {\end{list}}
  \renewenvironment{solution}[1]{%
    \begin{list}{}{%
        \renewcommand{\makelabel}{%
          \bfseries\protect\hyperlink{ex:#1}{#1}}
        \settowidth{\labelwidth}{\bfseries #1}
        \setlength{\leftmargin}{\labelwidth}
        \addtolength{\leftmargin}{\labelsep}
        \setlist[enumerate,1]{label=\alph*),labelsep=*,leftmargin=1.5em}
        \setlist[enumerate,2]{label=\roman*),labelsep=0.5em,align=right}}
    \item}
    {\end{list}}

  %% Corriger un bogue dans answers apparu dans TeX Live 2020.
  %% https://tex.stackexchange.com/a/537873
  \makeatletter
  \renewcommand{\Writetofile}[2]{%
    \@bsphack
    \Iffileundefined{#1}{}{%
      \Ifopen{#1}{%
        {%
          \begingroup
          % \let\protect\string %%% <----- removed
          \Ifanswerfiles{%
            \protected@iwrite{\@nameuse{#1@file}}{\afterassignment\let@protect@string}{#2}%
          }{}%
          \endgroup
        }%
      }{}%
    }%
    \@esphack
  }
  \def\let@protect@string{\let\protect\string}
  \makeatother

  %% Algorithmes en pseudocode
  \newenvironment{pseudocode}{%
    \begin{list}{}{%
      \renewcommand{\makelabel}{}
      \setlength{\topsep}{2pt}
      \setlength{\labelwidth}{0pt}
      \setlength{\leftmargin}{0pt}}
    \item}%
  {\end{list}}

  %% Listes de commandes
  \newenvironment{ttscript}[1]{%
    \begin{list}{}{%
        \setlength{\labelsep}{1.5ex}
        \settowidth{\labelwidth}{\texttt{#1}}
        \setlength{\leftmargin}{\labelwidth}
        \addtolength{\leftmargin}{\labelsep}
        \setlength{\parsep}{0.5ex plus0.2ex minus0.2ex}
        \setlength{\itemsep}{0.3ex}
        \renewcommand{\makelabel}[1]{\vphantom{|}##1\hfill}}}
    {\end{list}}

  %%% =====================
  %%%  Nouvelles commandes
  %%% =====================

  %% Noms de fonctions, code, etc.
  \newcommand{\code}[1]{\texttt{#1}}
  \newcommand{\pkg}[1]{\textbf{#1}}

  %% Expressions régulières; utilise le paquetage local
  %% ullrcorners pour tracer les coins autour de l'expression et le
  %% paquetage soul pour mettre les correspondances en surbrillance.
  \newcommand{\regex}[1]{\ullrcorners{regex}{\code{#1}}}
  \newcommand{\chaine}[1]{\code{#1}}
  \newcommand{\matches}{\sethlcolor{regexmatch}\hl}

  %% J'ai besoin d'afficher le symbole d'espace de \verb* à
  %% l'intérieur d'expressions régulières, donc en argument à \regex.
  %% Pour ce faire, le symbole est sauvegardé à l'intérieur d'une
  %% boite. [https://tex.stackexchange.com/a/50805]
  \newsavebox{\spacebox}
  \begin{lrbox}{\spacebox}
    \verb*! !
  \end{lrbox}
  \newcommand{\aspace}{\usebox{\spacebox}}

  %% Hyperlien avec symbole de lien externe juste après; second
  %% argument peut être vide pour afficher l'url comme lien
  %% [https://tex.stackexchange.com/q/53068/24355 pour procédure de
  %% test du second paramètre vide]
  \newcommand{\link}[2]{%
    \def\param{#2}%
    \ifx\param\empty
      \href{#1}{\nolinkurl{#1}~{\smaller\faExternalLink*}}%
    \else
      \href{#1}{#2~{\smaller\faExternalLink*}}%
    \fi
  }

  %% Boite additionnelle (basée sur awesomebox.sty) pour liens vers
  %% des vidéos et ajout à la liste des vidéos.
  \newcounter{videos}[chapter]
  \newcommand{\videobox}[2]{%
    \refstepcounter{videos}
    \addcontentsline{lov}{figure}{\protect\numberline{\thesection}\protect#1}
    \awesomebox[url]{\aweboxrulewidth}{\faYoutube}{url}{#2}}

  %% Boites additionnelles (basées sur awesomebox.sty) pour remarques
  %% spécifiques à macOS et Windows.
  \newcommand{\macosbox}[1]{%
    \awesomebox{\aweboxrulewidth}{\faApple}{black}{#1}}
  \newenvironment{macosblock}{%
    \begin{awesomeblock}{\aweboxrulewidth}{\faApple}{black}}
    {\end{awesomeblock}}
  \newcommand{\windowsbox}[1]{%
    \awesomebox{\aweboxrulewidth}{\faWindows}{black}{#1}}

  %% Boite additionnelle (basée sur awesomebox.sty) pour changements
  %% au fil de la lecture.
  \newcommand{\gotorbox}[1]{%
    \awesomebox{\aweboxrulewidth}{\faMapSigns}{black}{#1}}

  %% Boite pour le nom du fichier de script correspondant au début des
  %% sections d'exemples.
  \newcommand{\scriptfile}[1]{%
    \begingroup
    \noindent
    \mbox{%
      \makebox[3mm][l]{\raisebox{-0.5pt}{\small\faChevronCircleDown}}\;%
      \smaller[1] {\sffamily Fichier d'accompagnement} {\ttfamily #1}}
    \endgroup}

  %% Lien vers GitLab dans la page de notices
  \newcommand{\viewsource}[1]{%
    \href{#1}{\faGitlab\ Voir sur GitLab}}

  %% «Logo» C++
  \newcommand{\Rplus}{\protect\hspace{-.05em}\protect\raisebox{.2ex}{\smaller\textbf{+}}}
  \newcommand{\Cpp}{\mbox{C\Rplus\Rplus}\xspace}

  %% Raccourcis usuels vg
  \newcommand{\pt}{{\scriptscriptstyle \Sigma}}
  \newcommand{\mat}[1]{\mathbf{#1}}
  \newcommand{\var}[1]{\operatorname{Var} [ #1 ]}

  %%% =======
  %%%  Index
  %%% =======
  \newcommand{\bfhyperpage}[1]{\textbf{\hyperpage{#1}}}
  \renewcommand{\preindexhook}{%
    \label{chap:index}
    Les numéros de page en caractères gras indiquent les pages où les
    concepts sont introduits, définis ou expliqués.\vskip\onelineskip}
  \newcommand{\Index}[1]{\index{#1|bfhyperpage}}

  %% Les entrées d'index générées par listings avec \indexcode,
  %% ci-dessus, causent problème lorsque le mot clé est suivi d'une
  %% virgule. La redéfinition de \lst@nolig règle le problème.
  %% https://tex.stackexchange.com/q/510104
  \makeatletter
  \newcommand{\indexcode}[2][]{%
    \begingroup
    \def\lst@nolig{}%
    \def\param{#1}%
    \ifx\param\empty
      \index{#2@\code{#2}}%
    \else
      \index{#2@\code{#2} (#1)}%
    \fi
    \endgroup
  }
  \makeatother
  \newcommand{\Indexcode}[2][]{%
    \def\param{#1}%
    \ifx\param\empty
      \Index{#2@\code{#2}}%
    \else
      \Index{#2@\code{#2} (#1)}%
    \fi
  }
  \newcommand{\icode}[2][]{\indexcode[#1]{#2}\code{#2}}
  \newcommand{\Icode}[2][]{\Indexcode[#1]{#2}\code{#2}}

  %% Configuration des titres des regroupements par lettre et
  %% inscription dans la tdm du PDF; tiré de memoir. Voir aussi le
  %% fichier .ist.
  \makeatletter
  \newcommand{\doidxbookmark}[1]{{\def\@tempa{Symbols}\def\@tempb{#1}%
      \centering\bfseries \ifx\@tempa\@tempb %
      Symboles
      \phantomsection%
      \pdfbookmark[1]{Symboles}{symbols-idx}%
      \else
      #1%
      \phantomsection%
      \pdfbookmark[1]{#1}{#1-idx}%
      \fi%
      \nopagebreak\par}}
  \makeatother

  %% Création de l'index; options de imakeidx
  \makeindex[intoc,columns=2,options=-s \jobname.ist]

  %%% =======
  %%%  Varia
  %%% =======

  %% Sous-figures
  \newsubfloat{figure}

  %% Boites utilisées pour la construction d'images
  \newsavebox{\planche}
  \newsavebox{\planchep}
  \newsavebox{\supports}
  \newsavebox{\peintre}
  \newsavebox{\cube}
  \newsavebox{\mcube}
  \newsavebox{\dcube}
  \newsavebox{\ligne}
  \newsavebox{\colonne}
  \newsavebox{\carotte}
  \newsavebox{\matrice}
  \newsavebox{\trancheh}
  \newsavebox{\tranchev}
  \newsavebox{\tableau}
  \newsavebox{\cubedouble}
  \newsavebox{\polygone}
  \newsavebox{\gpolygone}
  \newsavebox{\nametag}
  \newsavebox{\ccube}
  \newsavebox{\mccube}
  \newsavebox{\dccube}
  \newsavebox{\cligne}
  \newsavebox{\ccolonne}
  \newsavebox{\ccarotte}
  \newsavebox{\cmatrice}
  \newsavebox{\ctrancheh}
  \newsavebox{\ctranchev}
  \newsavebox{\ctableau}

  %% Style de la bibliographie
  \bibliographystyle{francais}

  %% Permettre la coupure des url aux traits d'union
  %% https://tex.stackexchange.com/a/256190
  \def\UrlBreaks{\do\/\do-}

  %% Compteur bidon pour hyperlien vers des flottants non numérotés
  %% https://stackoverflow.com/q/2745731
  \newcounter{dummy}

  %% Longueurs pour la composition des pages couvertures avant et
  %% arrière.
  \newlength{\banderougewidth} \newlength{\banderougeheight}
  \newlength{\bandeorwidth}    \newlength{\bandeorheight}
  \newlength{\imageheight}
  \newlength{\logoheight}

  %% Aide pour la césure
  \hyphenation{%
    con-nai-tre
    con-sole
    cons-tante
    cons-tantes
    con-tenu
    con-trôle
    nom-bre
    puis-que
  }

  \includeonly{donnees}

\begin{document}

\frontmatter

\pagestyle{empty}

%%%
%%% Page de titre
%%%

\begingroup
\TPGrid{3}{36}
\textblockorigin{0mm}{0mm}
\setlength{\parindent}{0mm}
\setlength{\imageheight}{29\TPVertModule}
\setlength{\banderougewidth}{2\TPHorizModule}
\setlength{\banderougeheight}{\TPVertModule}
\setlength{\bandeorwidth}{\TPHorizModule}
\setlength{\bandeorheight}{\banderougeheight}
\setlength{\logoheight}{2.5\TPVertModule}
\setlength{\gapwidth}{1.5pt}
\addtolength{\bandeorwidth}{-\gapwidth}
\addtolength{\imageheight}{-\gapwidth}

\def\titlefmt{%
  \sffamily\bfseries\fontsize{52}{52}\selectfont\thetitle}
\def\authorfmt{%
  \sffamily\mdseries\fontsize{30}{38}\selectfont\theauthor}
\def\affiliation{%
  \sffamily\mdseries\fontsize{22}{26}\selectfont
  Professeur titulaire \\
  École d'actuariat, Université Laval}
\def\edition{%
  \sffamily\mdseries\fontsize{22}{26}\selectfont
  Édition {\fullcaps\year}.\month}

%% bandeau identitaire
\begin{textblock*}{\paperwidth}[0,1](0mm,30\TPVertModule)
  \textcolor{rouge}{\rule{\banderougewidth}{\banderougeheight}}% % bande rouge
  \rule{\gapwidth}{0pt}%                                         % filet
  \textcolor{or}{\rule{\bandeorwidth}{\bandeorheight}}           % bande or
\end{textblock*}

%% logo UL
\begin{textblock*}{\TPHorizModule}(2\TPHorizModule,31\TPVertModule)
  \rule{\gapwidth}{0pt}%                                         % filet
  \includegraphics[height=\logoheight,%
                   keepaspectratio=true]{ul_p}
\end{textblock*}

%% image de fond
\begin{textblock*}{\paperwidth}(0mm,0mm)
  \includegraphics[height=\imageheight,%
                   width=\paperwidth,%
                   trim=555 185 251 0,clip]{Xerus_inauris_1}
\end{textblock*}

%% titre
\begin{textblock*}{\paperwidth}(0.3\TPHorizModule,4.5\TPVertModule)
  \titlefmt
\end{textblock*}

%% auteur
\begin{textblock*}{\paperwidth}(0.3\TPHorizModule,7.5\TPVertModule)
  \authorfmt
\end{textblock*}

\null\cleardoublepage

%%%
%%% Page frontispice
%%%

%% titre
\begin{textblock*}{\paperwidth}(0.3\TPHorizModule,4.5\TPVertModule)
  \titlefmt
\end{textblock*}

%% auteur
\begin{textblock*}{2\TPHorizModule}(0.3\TPHorizModule,7.5\TPVertModule)
  \authorfmt
\end{textblock*}

%% affiliation
\begin{textblock*}{2\TPHorizModule}(0.3\TPHorizModule,9\TPVertModule)
  \affiliation
\end{textblock*}

%% édition
\begin{textblock*}{1.7\TPHorizModule}(0.3\TPHorizModule,30\TPVertModule)
  \edition
\end{textblock*}
\endgroup

%%% Local Variables:
%%% mode: latex
%%% TeX-master: "programmer-avec-r"
%%% TeX-engine: xetex
%%% End:

\null\cleartoverso              % cf. section 2.2 textpos.pdf

%%% Copyright (C) 2018 Vincent Goulet
%%%
%%% Ce fichier fait partie du projet
%%% «Programmer avec R»
%%% https://gitlab.com/vigou3/programmer-avec-r
%%%
%%% Cette création est mise à disposition selon le contrat
%%% Attribution-Partage dans les mêmes conditions 4.0
%%% International de Creative Commons.
%%% https://creativecommons.org/licenses/by-sa/4.0/

\begingroup
\calccentering{\unitlength}
\begin{adjustwidth*}{\unitlength}{-\unitlength}
  \setlength{\parindent}{0pt}
  \setlength{\parskip}{\baselineskip}
  \small

  \raisebox{-2.5mm}{%
    \includegraphics[height=7mm,keepaspectratio=true]{images/by-sa}} %
  \raisebox{-2.5mm}{%
    \includegraphics[height=7mm,keepaspectratio=true]{images/by-nc-sa}} %
  {\theauthor}, {\year}

  {\textcopyright} {\year} par {\theauthor}. «\thetitle», à
  l'exception les sections
  \ref*{sec:git:getting_started}--\ref*{sec:git:branching}, est mis à
  disposition selon le contrat
  \href{https://creativecommons.org/licenses/by-sa/4.0/deed.fr}{%
    Attribution-Partage dans les mêmes conditions 4.0 International}
  de Creative Commons. En vertu de ce contrat, vous êtes libre de:
  \begin{itemize}
  \item \textbf{partager} --- reproduire, distribuer et communiquer
    l'œuvre;
  \item \textbf{remixer} --- adapter l'œuvre;
  \item utiliser cette œuvre à des fins commerciales.
  \end{itemize}
  Selon les conditions suivantes:

  \begin{tabularx}{\linewidth}{@{}lX@{}}
    \raisebox{-9mm}[0mm][13mm]{%
      \includegraphics[height=11mm,keepaspectratio=true]{images/by}}
    & \textbf{Attribution} --- Vous devez créditer l'œuvre, intégrer
    un lien vers le contrat et indiquer si des modifications ont été
    effectuées à l'œuvre. Vous devez indiquer ces informations par
    tous les moyens possibles, mais vous ne pouvez suggérer que
    l'Offrant vous soutient ou soutient la façon dont vous avez utilisé
    son œuvre. \\
    \raisebox{-9mm}{\includegraphics[height=11mm,keepaspectratio=true]{images/sa}}
    & \textbf{Partage dans les mêmes conditions} --- Dans le cas où vous
    modifiez, transformez ou créez à partir du matériel composant
    l'œuvre originale, vous devez diffuser l'œuvre modifiée dans
    les même conditions, c'est à dire avec le même contrat avec lequel
    l'œuvre originale a été diffusée.
  \end{tabularx}

  Les sections
  \ref*{sec:git:getting_started}--\ref*{sec:git:branching} sont
  dérivées de «Pro Git» de Scott Chacon et Ben Straub sous contrat
  \href{https://creativecommons.org/licenses/by-nc-sa/3.0/deed.fr}{%
    Attribution-Pas d’utilisation commerciale-Partage dans les mêmes
    conditions 3.0 non transposé} de Creative Commons. Ce contrat
  ajoute aux conditions ci-dessus:

  \begin{tabularx}{\linewidth}{@{}lX@{}}
    \raisebox{-9mm}[0mm][13mm]{%
      \includegraphics[height=11mm,keepaspectratio=true]{images/nc}}
    & \textbf{Pas d’utilisation commerciale} --- Vous n'êtes pas
    autorisé à faire un usage commercial de cette œuvre, tout ou
    partie du matériel la composant.
  \end{tabularx}

  \textbf{Code source} \\
  \viewsource{\reposurl}

  \textbf{Couverture} \\
  Groupe d'écureuils de terre du Cap (\emph{Xerus inauris}) sortant de
  leur terrier près de Solitaire, dans le désert du Namib, en Namibie.

  Crédit photo: {\textcopyright} Hans Hillewaert,
  \href{https://creativecommons.org/licenses/by-sa/3.0/deed.fr}{CC BY-SA 3.0}, via
  \href{https://commons.wikimedia.org/w/index.php?curid=2429056}{Wikimedia Commons}.
\end{adjustwidth*}
\endgroup

%%% Local Variables:
%%% mode: latex
%%% TeX-engine: xetex
%%% TeX-master: "programmer-avec-r"
%%% coding: utf-8
%%% End:

\clearpage

\pagestyle{companion}

\chapter*{Introduction}
\addcontentsline{toc}{chapter}{Introduction}
\markboth{Introduction}{Introduction}


Il existe déjà de multiples ouvrages traitant de S-Plus ou \textsf{R}.
Dans la majorité des cas, toutefois, ces deux logiciels sont présentés
dans le cadre d'applications statistiques spécifiques. Le présent
ouvrage se concentre plutôt sur l'apprentissage du langage de
programmation sous-jacent aux diverses fonctions statistiques, le S.

Cette seconde édition est principalement une réorganisation du contenu
de la première édition. Les chapitres sur la régression linéaire et
les séries chronologiques ont été éliminés et les annexes C et D
ont été déplacées dans le corps du document. Nous avons également
corrigé plusieurs coquilles suite à la révision effectuée par
Mme~Mireille Côté.

Le texte est la somme de notes et d'exercices de cours donnés par
l'auteur à l'École d'actuariat de l'Université Laval. Les six premiers
chapitres, qui constituent le c{\oe}ur du document, proviennent d'une
partie d'un cours où l'accent est mis sur l'apprentissage d'un
(deuxième) langage de programmation par des étudiants de premier cycle
en sciences actuarielles. Les applications numériques et statistiques
de S-Plus et \textsf{R} présentées aux chapitres \ref{optimisation},
\ref{rng} et \ref{simulation} sont étudiées plus tard dans le cursus
universitaire.

Les cours d'introduction au langage S sont donnés à raison de une
heure par semaine de cours magistral suivie de deux heures en
laboratoire d'informatique. C'est ce qui explique la structure des six
premiers chapitres: les éléments de théorie, contenant peu voire aucun
exemple, sont présentés en rafale en classe. Puis, lors des séances de
laboratoire, les étudiantes et étudiants sont appelés à lire et
exécuter les exemples se trouvant à la fin des chapitres. Chaque
section d'exemples couvre l'essentiel des concepts présentés dans le
chapitre et les complémente souvent. L'étude de ces sections fait donc
partie intégrante de l'apprentissage du langage S.

Le texte des sections d'exemples est disponible en format électronique
dans le site Internet
\begin{quote}
  \url{http://vgoulet.act.ulaval.ca/intro_S}
\end{quote}

Certains exemples et exercices trahissent le premier public de ce
document: on y fait à l'occasion référence à des concepts de base de
la théorie des probabilités et des mathématiques financières. Les
contextes actuariels demeurent néanmoins peu nombreux et ne devraient
généralement pas dérouter le lecteur pour qui ces notions sont moins
familières.

Les chapitres \ref{optimisation} (fonctions d'optimisation), \ref{rng}
(générateurs de nombres aléatoires) et \ref{simulation} (planification
d'une simulation en S) sont structurés de manière plus classique,
notamment parce que le texte y est en prose.

Le texte prend parti en faveur de l'utilisation de GNU Emacs et du
mode ESS pour l'édition de code S.  Les annexes contiennent de
l'information sur l'utilisation de S-Plus et \textsf{R} avec cet
éditeur.

Dans la mesure du possible, cet ouvrage tâche de présenter les
environnements S-Plus et \textsf{R} en parallèle, en soulignant leurs
différences s'il y a lieu. Les informations propres à S-Plus ou à
\textsf{R} sont d'ailleurs signalées en marge par les marques «S$+$»
et «\textsf{R}», respectivement. Étant donné la nette préférence de
l'auteur pour \textsf{R}, les divers extraits de code ont généralement
été exécutés avec ce moteur S.

À moins d'erreurs et d'omissions (que les lecteurs sont invités à nous
faire connaître), les informations données à propos de S-Plus sont
exactes pour les versions 6.1 (Linux et Windows), 6.2 Student Edition
(Windows) et 7.0 (Linux et Windows). Pour \textsf{R}, la version 2.4.1
(Linux et Windows), soit la plus récente lors de la rédaction, a été
utilisée comme référence.

On notera enfin que cet ouvrage n'a aucune prétention d'exhaustivité.
C'est ce qui explique les nombreux renvois au livre de \cite{MASS},
plus complet.

L'auteur tient à remercier M.~Mathieu Boudreault pour sa collaboration
dans la rédaction des exercices.

\begin{flushright}
  \itshape
  Vincent Goulet \url{<vincent.goulet@act.ulaval.ca>} \\
  Québec, janvier 2007
\end{flushright}


%%% Local Variables:
%%% mode: latex
%%% TeX-master: "introduction_programmation_S"
%%% coding: utf-8
%%% End:


\tableofcontents
\cleartorecto
\listoftables
\cleartorecto
\listoffigures
\cleartorecto
\listofvideos

\chapter*{Comment (vraiment) bien \\ utiliser ce document}

\begin{enumerate}
\item Lire l'introduction.
\item Consulter le document dans une visionneuse PDF (Aperçu, Acrobat
  Reader) plutôt que dans un navigateur web.
\item Suivre le guide dans la lecture des chapitres. Aux avis signalés
  par le symbole {\faMapSigns}, passer à l'étude du code R des
  sections d'exemples pour revenir ensuite au texte.
\item Lire les commentaires dans les fichiers de script R et prendre
  le temps d'analyser et de bien comprendre les résultats.
\item Travailler sur les exercices seulement après avoir complété les
  étapes précédentes.
\end{enumerate}

%%% Local Variables:
%%% mode: latex
%%% TeX-engine: xetex
%%% TeX-master: "programmer-avec-r"
%%% coding: utf-8
%%% End:


%% Citation de D.E. Knuth (recto)
\cleartorecto
\thispagestyle{empty}
\begin{vplace}[0.45]
  \begin{quote}
    L'ultime test pour vérifier si je comprends quelque chose
    consiste à l'expliquer à un ordinateur. Je peux vous dire quelque
    chose et vous hocherez la tête, mais sans que je sois certain d'avoir
    bien expliqué. L'ordinateur, lui, ne hoche pas la tête. Il répète
    exactement ce que je lui ai dit. Vous pouvez faire illusion dans
    plusieurs aspects de la vie, mais pas avec les ordinateurs.
    \sourceatright{Donald E.~Knuth}
  \end{quote}
\end{vplace}

%% Vignette tirée de xkcd.com (verso)
\cleartoverso
\thispagestyle{empty}
\begin{vplace}[0.45]
  \centering
  \begin{minipage}{0.9\linewidth}
    \setkeys{Gin}{width=\textwidth}
    \includegraphics{images/compiling.png} \\
    \footnotesize\sffamily%
    Tiré de \href{https://xkcd.com/303/}{XKCD.com}
  \end{minipage}
\end{vplace}

\mainmatter

%%% Copyright (C) 2017 Vincent Goulet
%%%
%%% Ce fichier fait partie du projet «Programmer avec R»
%%% http://github.com/vigou3/programmer-avec-r
%%%
%%% Cette création est mise à disposition selon le contrat
%%% Attribution-Partage dans les mêmes conditions 4.0
%%% International de Creative Commons.
%%% http://creativecommons.org/licenses/by-sa/4.0/

\chapter{Éléments d'informatique pour programmeurs}
\label{chap:informatique}

% Intro du chapitre

\section{Bref historique des langages de programmation}
\label{sec:informatique:historique}

Ada Lovelace (1815--1852) est généralement reconnue comme la première
auteure d'un algorithme et de ce que l'on appelle aujourd'hui un
programme informatique. C'est en son honneur qu'a été nommé le langage
Ada conçu en réponse à un cahier de charges du département de la
Défense des États-Unis au début des années 1980.

Franchissons d'un bond les premiers jours de l'informatique et de la
programmation pour arriver aux ordinateurs électriques modernes, dans
les années 1940. On programme alors généralement ceux-ci en
\emph{assembleur}, un langage de très bas niveau facilement
interprétable par la machine, mais difficile à lire par des humains,
comme l'extrait de programme de la
\autoref{fig:informatique:assembleur} le démontre bien.

\begin{figure}
  \centering
  \includegraphics[trim=0 475 0 0, clip=true]{Motorola_6800_Assembly_Language.png}
  \caption[Programme en assembleur pour un microprocesseur 8~bits
  Motorola 6800.]{Extrait d'un programme en assembleur pour un
    microprocesseur 8~bits Motorola 6800. {\small Source:
      \link{https://commons.wikimedia.org/wiki/File\%3AMotorola_6800_Assembly_Language.png}{Wikimedia
        Commons}}}
  \label{fig:informatique:assembleur}
\end{figure}

Les premiers langages créés pour transmettre des instructions à un
ordinateur apparaissent dans les années 1950. Ils sont à l'origine
intimement liés à l'architecture d'un ordinateur: à chaque type
d'ordinateur son langage de programmation. Certaines contraintes ---
ou exigences --- des langages de l'époque proviennent aussi du support
physique alors utilisé pour stocker les programmes: les cartes
perforées.

\subsection{Fortran}
\label{sec:informatique:historique:fortran}

En 1954, l'ingénieur de IBM John Backus publie les spécifications du
langage FORTRAN (\emph{FORmula TRANslating System}). Le premier
compilateur voit le jour deux ans plus tard. Fortran (avec des
minuscules à partir de 1977) deviendra rapidement le langage standard
dans le calcul scientifique.

Plus d'un demi-siècle plus tard, l'empreinte de Fortran demeure
importante, notamment grâce aux bibliothèques d'algèbre linéaire
BLAS\footnote{%
  \emph{Basic Linear Algebra Subprograms};
  \link{http://www.netlib.org/blas}{}.} %
et LAPACK\footnote{%
  \emph{Linear Algebra PACKage};
  \link{http://www.netlib.org/lapack}{}.} %
auxquelles ont recours la plupart des progiciels scientifiques, dont
R. Le langage est toujours utilisé en calcul haute performance et pour
mesurer le rendement des superordinateurs.

\begin{figure}[t]
  \notebox{Dans le très recommandable film Les figures de l'ombre
    (\emph{Hidden Figures}, 2016), une des héroïnes entreprend de
    s'attaquer à la programmation des nouveaux ordinateurs de la NASA.
    On peut alors voir qu'elle apprend le Fortran.}
\end{figure}

\subsection{Lisp}
\label{sec:informatique:historique:lisp}

Le deuxième langage le plus ancien toujours largement diffusé est Lisp
(\emph{LISt Processing}). Créé par John McCarthy en 1958 en tant que
modèle pratique pour représenter des programmes, le Lisp est devenu le
langage de choix pour la recherche et les applications en intelligence
artificielle. Le terme Lisp désigne aujourd'hui une famille de
langages comprenant de nombreux dialectes, dont Common Lisp, Scheme et
Emacs Lisp.

Le Lisp se distingue en outre par une syntaxe simple en notation
préfixée (voir encadré), le support pour la programmation
fonctionnelle (\autoref{sec:informatique:paradigmes}) et la faculté de
manipuler le code source en tant que structure de données. Autre trait
distinctif: la syntaxe du Lisp fait un usage immodéré des parenthèses.

Le Lisp est entouré d'une aura de beauté et d'élégance dont peu
d'autres langages peuvent se targuer. Citons Eric Raymond dans
\link{http://www.catb.org/esr/faqs/hacker-howto.html}{\emph{How to
    Become a Hacker}}:
\begin{quote}
  Il faut apprendre le Lisp pour l'extraordinaire expérience d'éveil
  [\emph{enlightenment experience}] que procure le fait de finalement
  le comprendre; cette expérience fera de vous un meilleur programmeur
  pour toujours, même si vous n'avez plus vraiment à utiliser le Lisp.
\end{quote}

Pour illustrer encore davantage la place toute particulière qu'occupe
le Lisp en programmation, mentionnons également l'aphorisme selon
lequel «ceux qui ne connaissent pas le Lisp sont condamnés à le
réinventer», d'une certaine façon la version courte de la célèbre
\link{https://en.wikipedia.org/wiki/Greenspun\%27s_tenth_rule}{\emph{Greenspun's
    tenth rule of programming}}: %
\begin{quote}
  Tout programme en C ou en Fortran suffisamment complexe contient une
  mise en œuvre mal spécifiée, pleine de bogues et lente de la moitié
  de Common Lisp.
\end{quote}
(Fait amusant à noter: il n'y a pas d'autres lois que la dixième, de
l'\link{http://philip.greenspun.com/bboard/q-and-a-fetch-msg?msg_id=000tgU}{aveu
  même de l'auteur}.)

\begin{figure}[t]
  \label{fig:informatique:notations}
  \setlength{\FrameRule}{1pt}
  \lstset{backgroundcolor=\color{codebg},
    frame=lr, rulecolor=\color{codebg},
    xleftmargin=3.4pt, xrightmargin=3.4pt}
  \begin{emphbox}{\mdseries Notations infixée, préfixée et suffixée}
    Les notations infixée (\emph{infix}), préfixée (\emph{prefix}) et
    suffixée (\emph{postfix}) sont trois manières différentes, mais
    équivalentes, d'écrire des expressions en mathématiques ou en
    programmation.

    Par exemple, l'opération d'addition de deux opérandes $x$ et $y$
    s'écrit en notation infixée
\begin{lstlisting}
x + y
\end{lstlisting}
    En notation préfixée, aussi appelée notation polonaise,
    l'opérateur est placé \emph{avant} les opérandes:
\begin{lstlisting}
+ x y
\end{lstlisting}
    On l'aura compris, en notation suffixée, ou notation polonaise
    inversée, l'opérateur apparait \emph{après} les opérandes:
\begin{lstlisting}
x y +
\end{lstlisting}
    Nous sommes davantage habitués à lire la notation infixée, quoique
    la notation préfixée nous soit familière pour les opérateurs à un
    seul opérande (comme la négation) ou pour les appels de fonctions.
    La notation suffixée n'a jamais recours aux parenthèses.

    La compagnie HP commercialise de très prisées calculatrices
    scientifiques utilisant la notation suffixée (libellées RPN pour
    \emph{Reverse Polish Notation}) depuis 1972.
  \end{emphbox}
\end{figure}

\subsection{COBOL}
\label{sec:informatique:historique:cobol}

Le troisième langage développé dans les années 1950 et toujours en
usage de nos jours est COBOL (\emph{COmmon Business Oriented
  Language}). Ce langage spécialisé dans les applications de gestion a
été créé en 1959 par un comité formé pour proposer un langage commun
pour l'administration américaine.

Le COBOL reste très utilisé dans de grandes entreprises, notamment
dans les institutions financières. La légende urbaine veut d'ailleurs
que les programmeurs COBOL soient comparativement très bien rémunérés
aujourd'hui sous l'effet combiné de leur rareté et de l'importance
opérationnelle des applications qu'ils doivent maintenir.

\subsection{Algol}
\label{sec:informatique:historique:algol}

Dès la fin de la décennie 1950, un comité de scientifiques se réunit à
Zurich pour concevoir ce que l'on voudrait voir devenir le langage de
programmation standard. De ces rencontres naîtra Algol
(\emph{ALGorithmic Oriented Language}) en 1958. Comme la plupart des
tentatives de définition d'un standard, c'est un échec: le langage est
populaire dans les milieux académiques, mais restera peu utilisé dans
les applications commerciales.

Cela dit, on doit à Algol plusieurs innovations importantes, de telle
sorte qu'un grand nombre des langages qui verront le jour par la suite
seront considérés comme ses descendants; le poster
\link{http://www.oreilly.com/go/languageposter}{\emph{History of
    Programming Languages}} de O'Reilly Media illustre ce fait à
merveille. \citet{Hoare:1973} a d'ailleurs cette jolie formule:
\begin{quote}
  Voici un langage très en avance sur son temps, il n'a pas seulement
  été une amélioration de ses prédécesseurs mais aussi une
  amélioration de presque tous ses successeurs.
\end{quote}


\begin{figure}[t]
  \centering
  \begin{minipage}{0.9\linewidth}
    \setkeys{Gin}{width=\textwidth}
    \includegraphics{standards} \\
    \footnotesize\sffamily%
    Tiré de \href{http://xkcd.com/927/}{XKCD.com}
  \end{minipage}
\end{figure}

\subsection{APL}
\label{sec:informatique:historique:apl}

Notre historique ne serait pas complet sans un mot sur APL (\emph{A
  Programming Language}, qui l'aurait cru). Même s'il n'a jamais connu
une diffusion importante, ce langage conçu par Kenneth Iverson autour
de 1962 n'en a pas moins eu une influence considérable sur la manière
de penser et de représenter les opérations mathématiques sur les
tableaux à plusieurs dimensions.

Doté d'une large gamme de symboles pour représenter des opérations et
d'une syntaxe tout à fait particulière --- les expressions sont
exécutées de droite à gauche! --- l'APL est remarquablement concis et
puissant; voir la \autoref{fig:informatique:apl} pour un aperçu. Le
revers de la cette médaille et ce qui a assurément nui à son adoption
à large échelle, c'est la difficulté que l'on éprouve à relire le
code. Assez pour que d'aucun qualifient l'APL de «langage à écriture
seulement».

APL a pendant longtemps été un langage très prisé par les actuaires,
aussi subsiste-t-il du code dans certaines compagnies d'assurance. Le
modèle de traitement des vecteurs, matrices et tableaux de l'APL a
servi d'inspiration pour la conception du langage S --- nous y
reviendrons au \autoref{chap:presentation}.

Le langage continue sa vie aujourd'hui principalement sous forme de son
successeur, J.

\begin{figure}
  \centering
  \scalebox{0.4}{\includegraphics{APL_intro}}
  \caption[Opérations simples en APL.]{Opérations simples en APL. De
    haut en bas: génération des nombres de $1$ à $5$ et stockage dans
    la variable \code{x}; affichage du contenu de \code{x}; somme des
    éléments de \code{x}; moyenne des éléments de \code{x}. {\small Source:
    François-Dominique,
    \link{https://commons.wikimedia.org/w/index.php?curid=43207460}{Wikimedia
      Commons}, CC BY-SA 4.0}}
  \label{fig:informatique:apl}
\end{figure}

\subsection{C}
\label{sec:informatique:historique:c}

Le langage C a été inventé en 1972 chez Bell Labs par Ken Thompson et
Dennis Ritchie afin de réécrire le système d'exploitation UNIX
(\autoref{sec:informatique:os:unix}). Il demeure beaucoup utilisé pour
la programmation système: le noyau de systèmes d'exploitation comme
Windows et Linux sont développés en grande partie en C.

Le C est un langage de programmation généraliste considéré, selon les
standards actuels, comme de bas niveau. Pour illustrer, l'utilisateur
doit programmer des traitements comme la libération de la mémoire, la
vérification de la validité des indices sur les tableaux, l'ouverture
et la fermeture des fichiers, etc.

Le langage demeure l'un des plus utilisés dans le monde et son
influence est considérable. De nombreux langages plus modernes comme
\Cpp, C\# et Java reprennent des aspects de C. Le C est également
beaucoup utilisé pour le calcul numérique intensif, où il s'est en
quelque sorte substitué à Fortran. La plupart des progiciels
scientifiques --- dont R, encore une fois --- offrent la possibilité
d'appeler du code C lorsque la rapidité de calcul devient un enjeu. Le
gain en temps d'exécution doit toutefois être suffisamment grand pour
compenser le temps de développement plus long qu'exige le C par
rapport à des langages de plus haut niveau.

\begin{figure}[t]
  \notebox{Dans l'ouvrage désormais classique de \cite{KandR:1978}, le
    premier exemple d'un programme C affiche le message «\emph{hello,
      world}» à l'écran. Ça deviendra ensuite une tradition de
    démontrer le fonctionnement ou la syntaxe d'un langage avec cet
    exemple.}
\end{figure}

\subsection{Autres jalons}
\label{sec:informatique:historique:autres}

Nous nous sommes attardés jusqu'ici à des langages de programmation
vieux de plus de 40 ans à cause de leur importance historique et parce
qu'ils sont toujours utilisés couramment. De nombreux autres langages
ont vu le jour depuis, si bien qu'ils se comptent aujourd'hui par
milliers. En voici quelques autres ayant occupé une place
prépondérante dans l'histoire.

\begin{itemize}
\item Comme son nom l'indique, \textbf{\Cpp} (Bjarne Stroustrup, 1980)
  est un dérivé du C qui lui ajoute, en autres choses, la
  programmation orientée objet. Certains considèrent que {\Cpp}
  devrait être le point d'entrée de toute personne voulant débuter à
  programmer avec des langages de la famille du C. Le code C est
  compatible avec le \Cpp.
\item Principal représentant de l'ère Internet des années 1990,
  \textbf{Java} (James Gosling, 1995) a été conçu pour que le code
  écrit dans ce langage puisse s'exécuter sur n'importe quelle
  plateforme informatique sans nécessiter une nouvelle compilation. Il
  est donc très populaire dans les applications web ou embarquées. Sa
  syntaxe est fortement inspirée du \Cpp.
\item \textbf{Visual Basic} (Microsoft, 1991) permet de développer des
  applications de manière interactive en disposant des composantes sur
  un canevas. Le langage est aujourd'hui discontinué, mais son dérivé
  \textbf{Visual Basic for Applications} (VBA) demeure beaucoup
  utilisé dans les applications de la suite bureautique Office.
\item \textbf{Python} (Guido van Rossum, 1991) est un langage de haut
  niveau, orienté objet, multiplateforme et sous licence libre. C'est
  un des langages les plus utilisés aujourd'hui pour le calcul
  scientifique et l'analyse de données massives.
\end{itemize}

Certains langages de programmation ont une vocation généraliste,
certains visent des niches particulières, alors que d'autres cherchent
surtout à faire progresser l'état des connaissances dans la théorie
des langages. Quoi qu'il en soit, un langage de programmation demeure
un outil et, en informatique comme dans d'autres domaines, il convient
de choisir le meilleur outil pour accomplir une tâche donnée.

Nous invitons les lecteurs intéressés à en savoir davantage sur
l'histoire des langages de programmation en général, ou sur l'un ou
l'autre des langages mentionnés ci-dessus en particulier, à débuter
par les très complètes entrées de Wikipedia en
\link{https://fr.wikipedia.org/wiki/Histoire_des_langages_de_programmation}{français}
et en
\link{https://en.wikipedia.org/wiki/History_of_programming_languages}{anglais}.


\section{Langages compilés et interprétés}
\label{sec:informatique:compile_vs_interprete}

Tous les langages de programmation\footnote{%
  À l'exception du langage machine, mais rares sont les personnes qui
  souhaitent ne programmer qu'avec des $0$ et des $1$.} %
requièrent un traitement afin que les programmes écrits avec ceux-ci
puissent être exploités par un ordinateur. Il existe deux grandes
façons d'effectuer ce traitement: par \emph{compilation} et par
\emph{interprétation}.

Un compilateur est un programme informatique qui transforme en langage
machine le code source rédigé dans un langage donné. On dit alors de
ce langage qu'il est \emph{compilé}. Le compilateur produit un fichier
informatique que l'ordinateur peut ensuite exécuter directement. Sur
la plateforme Windows, on reconnait notamment ces fichiers par leurs
extensions \code{.exe} ou \code{.dll}. Les noyaux d'à peu près toutes
les applications sur nos ordinateurs sont des fichiers compilés.

Un interprète (ou interpréteur), quant à lui, est un programme qui
analyse, traduit et exécute le code source d'un programme
informatique, et ce, à chaque fois que le programme doit être exécuté.
Un langage qui nécessite l'intermédiaire d'un interprète est dit
\emph{interprété}. Un programme écrit dans un langage interprété ne
peut être distribué sans son interprète.

Le traitement additionnel entre le code source et le langage machine
fait généralement en sorte que les langages interprétés sont plus
lents que les langages compilés à l'exécution. En revanche, une foule
de détails de mise en œuvre sont pris en charge par l'interprète, ce
qui réduit le temps de développement.


\section{Paradigmes de programmation}
\label{sec:informatique:paradigmes}

Un programme informatique est toujours écrit pour résoudre un problème
donné. Or, comme pour tout problème dans la vie, la solution à
celui-ci peut être formulée de différentes manières. Le style
fondamental avec lequel on exprime la solution est appelé le
\emph{paradigme} de programmation.

Il existe plusieurs paradigmes de programmation. Nous nous
contenterons ici de n'en présenter que quatre.

\begin{description}
\item[Impératif] On indique à l'ordinateur les opérations à exécuter
  et l'ordre dans lequel les exécuter. C'est le paradigme le plus
  intuitif.
\item[Déclaratif] On indique cette fois à l'ordinateur ce que l'on
  souhaite obtenir comme résultat, mais sans préciser comme y
  parvenir. Il est laissé à la mise en œuvre du langage de déterminer
  la meilleure méthode de résolution du problème.

  Par exemple, pour extraire d'une base de données \code{etudiants} le
  prénom de toutes les personnes de plus de 25~ans, on écrira dans un
  langage déclaratif comme le SQL:
\begin{Schunk}
\begin{Sinput}
SELECT prenom FROM etudiants WHERE age > 25
\end{Sinput}
\end{Schunk}
  Nulle part n'est-il précisé comment le programme devrait procéder à
  l'extraction.
\item[Fonctionnel] Un programme est une suite d'appels de fonctions,
  comme en mathématiques. L'exécution d'une fonction n'a pas d'impact
  sur les autres fonctions\footnote{%
    Autrement on dit qu'une fonction a un effet de bord (\emph{side
      effect}.} %
  Les opérations complexes sont réalisées en combinant les fonctions,
  de manière analogue à la composition de fonctions $g \circ f$.
\item[Orienté objet] Un programme est conçu comme un ensemble de blocs
  logiciels (les objets) qui interagissent entre eux. Une
  \emph{méthode} applique un traitement différent à un objet
  selon sa \emph{classe}. Ce paradigme est particulièrement utilisé
  dans les grands et complexes projets informatiques.
\end{description}

La plupart des langages d'usage courant combinent d'office, ou
laissent au programmeur la possibilité de combiner, plusieurs
paradigmes. Par exemple, le langage {\Cpp} est à la fois un langage
impératif et orienté objet.


\section{Sémantique et syntaxe}
\label{sec:informatique:semantique}

Il y a plusieurs parallèles à dresser entre les langages de
programmation et les langues parlées ou écrites: leur apprentissage
requiert de la pratique; on en connait jamais un trop grand nombre; le
premier est plus difficile à apprendre que les suivants.
L'informatique a également emprunté à la linguistique les notions de
\emph{sémantique} et de \emph{syntaxe}.

La sémantique est l'étude de ce que signifie un message ou un
programme informatique, c'est-à-dire de ce qu'il transmet ou exécute.
La syntaxe, quant à elle, étudie la structure du message ou du
programme. En simplifiant, disons que, pour une langue donnée, un
dictionnaire permet de connaître la sémantique, alors qu'une grammaire
décrit la syntaxe \citep{Hebenstreit:semantique}.

Par exemple, exprimons le message «J'ai soif» en anglais. Dire
«\emph{I am hungry}» relèverait d'une erreur de sémantique, puisque la
signification du message s'en trouve changée. En revanche, avec
«\emph{I have thirsty}» le bon message se rendra, mais peut-être pas
sans que l'erreur de syntaxe n'ait d'abord suscité une hésitation chez
l'interlocuteur.

Comme pour une langue, la maîtrise d'un langage de programmation exige
de respecter à la fois les règles de sémantique et les règles de
syntaxe du langage. Le second volet est relativement facile: si le
programme compile ou s'exécute, c'est généralement que la syntaxe est
respectée. Le respect de la sémantique, ou de la manière propre d'un
langage d'exprimer des idées, demande plus d'efforts et de pratique.


\section{Systèmes d'exploitation}
\label{sec:informatique:os}

Le système d'exploitation (\emph{operating system}, OS) est un
ensemble de programmes qui gère les ressources matérielles et
logicielles d'un ordinateur. Premier programme exécuté par
l'ordinateur lors de la mise en marche, le système d'exploitation
reçoit les demandes de ressources des applications --- espace mémoire,
unité de calcul, disque dur, communication avec les périphériques,
utilisation du réseau --- et les alloue en fonction de l'état du
système et des demandes concurrentes des autres applications. Tous les
logiciels applicatifs requièrent un système d'exploitation pour
fonctionner.

Il existe une grande variété de systèmes d'exploitation. Les plus
connus aujourd'hui sont, pour les ordinateurs personnels et les
serveurs, Windows de Microsoft, macOS de Apple et Linux, un logiciel
libre créé et toujours administré par Linus Torvalds. Pour les
appareils mobiles, deux systèmes d'exploitation se partagent
l'essentiel du marché: iOS et Android. Le premier est dérivé de macOS
et le second, de Linux.

Les systèmes macOS et Linux appartiennent à la famille des systèmes
Unix. Cela ne laisse donc en fait que deux grandes classes de systèmes
d'exploitation couramment utilisés: Windows et Unix.

\subsection{Windows}
\label{sec:informatique:os:windows}

Le système d'exploitation Windows de Microsoft équipe près de 90~\%
des ordinateurs personnels dans le monde. Dans les circonstances, rares
sont les personnes qui n'ont jamais utilisé le système, ne serait-ce
qu'à petites doses.

À l'origine, en 1985, Windows n'était qu'une interface graphique
superposée au véritable système d'exploitation des ordinateurs de type
«compatibles IBM PC», le DOS. Les versions grand public depuis
Windows~2000 sont plutôt basées sur le noyau de Windows~NT, un système
conçu à partir de zéro et lancé au milieu des années 1990 en tant que
système pour les entreprises.

Le plus grand avantage de Windows est son ubiquité. À peu près toutes
les applications, sauf peut-être celles s'adressant à un segment de
marché très précis, sont disponibles sous Windows, en code natif, via
un port depuis Unix, ou encore par l'intermédiaire d'une couche
logicielle d'émulation comme Cygwin ou MinGW/MSYS.

Un système d'exploitation multiutilisateur prend pour acquis que
plusieurs utilisateurs se partageront les ressources de l'ordinateur,
simultanément lorsqu'il s'agit d'un serveur, ou les uns après les
autres dans le cas d'un ordinateur personnel. Dans un tel système, on
établit des règles strictes d'accès aux ressources du système (seul
l'administrateur peut installer ou supprimer des applications) et
d'accès aux ressources des autres utilisateurs (Marianne n'a pas
accès aux fichiers d'Alexandre et vice-versa).

Les versions de Windows multiutilisateurs sont apparues relativement
tard sur le marché grand public, et ce, sans que le paradigme ne soit
imposé aux utilisateurs. Résultat: la plupart des utilisateurs de
Windows s'octroient à la ronde les droits d'administrateur et, pire,
plusieurs applications nécessitent toujours un accès total au système.
C'est par ces failles que se propagent virus et autres logiciels
malveillants sur la plateforme.

Les systèmes Windows sont livrés avec une trousse de développement
très restreinte. Les programmeurs doivent donc installer les outils
nécessaires, souvent sous forme de suites logicielles monolithiques,
ce qui entraîne dédoublements et, parfois, conflits.

\subsection{Unix}
\label{sec:informatique:os:unix}

Le terme «Unix» --- ou UNIX qui, en majuscules, est une marque déposée
de Open Group --- recoupe une famille de systèmes d'exploitation
dérivant tous du premier système Unix développé chez Bell Labs dans
les années 1970 par Kenneth Thompson, Dennis Ritchie et plusieurs
autres.

Les systèmes Unix partagent un certain nombre de caractéristiques
communes. D'abord, ils sont intrinsèquement multitâches et
multiutilisateurs. Ensuite, la «philosophie Unix» commande la
modularité: le système d'exploitation fournit aux utilisateurs
(via un interpréteur de commande, ou \emph{shell}) et aux applications
une multitude de petits outils très spécialisés qui ne réalisent à peu
près qu'une seule tâche. Le système rend ensuite simple, via un
mécanisme dit de «redirection» (\emph{pipe}) de les combiner pour
effectuer des opérations plus complexes.

Par exemple, pour extraire le mot \emph{redirection} du code source
contenant le précédent paragraphe, nous pourrions utiliser l'outil
\code{grep} pour extraire la ligne où se trouve le mot, puis l'outil
\code{cut} pour isoler le mot qui nous intéresse et, enfin, l'outil
\code{tr} pour supprimer les guillemets autour du mot.\footnote{%
  Dans la mesure où le mot recherché est le quatrième de la ligne dans
  le fichier source \code{informatique.tex}, la commande résultante
  est:
  \code{grep "«redirection»" informatique.tex | cut -d " " -f 4 | tr -d «»}\,.
  Le caractère \code{|} est la commande de redirection.}

À l'exception notable des Mac, Unix demeure assez peu répandu sur les
ordinateurs personnels. En revanche le système d'exploitation, en
particulier sa variante libre Linux, équipe l'immense majorité des
stations de travail, serveurs et supercalculateurs. De plus, bon
nombre d'applications scientifiques et de langages de programmation
sont d'abord développés sur des systèmes Unix. Pour ces
raisons, nous estimons utile pour des programmeurs de connaître les
quelques «Unix-ismes» ci-dessous.

\begin{itemize}
\item La barre oblique (\code{/}) est utilisée pour séparer les
  dossiers dans les chemins d'accès aux fichiers (plus de détails à la
  \autoref{sec:informatique:fs}).
\item Tout utilisateur dispose d'un espace disque réservé identifié
  comme le «répertoire personnel» ou le «dossier de départ»;
  \emph{home directory}, en anglais. Ce répertoire est situé dans
  \code{/Users} sous macOS et habituellement dans \code{/home} sous
  Linux.
\item À l'invite de commande et dans la configuration de plusieurs
  applications, le caractère \,\verb=~=\, représente le répertoire
  personnel.
\item Toujours à l'invite de commande, il est possible d'appuyer sur
  \code{TAB} pour compléter les noms de fichiers, de répertoires ou de
  commandes.
\end{itemize}



%%% Local Variables:
%%% mode: latex
%%% TeX-engine: xetex
%%% TeX-master: "programmer-avec-r"
%%% End:

%%% Copyright (C) 2020 Vincent Goulet
%%%
%%% Ce fichier fait partie du projet
%%% «Programmer avec R»
%%% https://gitlab.com/vigou3/programmer-avec-r
%%%
%%% Cette création est mise à disposition sous licence
%%% Attribution-Partage dans les mêmes conditions 4.0
%%% International de Creative Commons.
%%% https://creativecommons.org/licenses/by-sa/4.0/

\chapter{Algorithmes et algorithmique}
\label{chap:algorithmes}

\def\scriptfilename{\currfilebase.R}

\begin{objectifs}
\item Résoudre un problème sous la forme d'un algorithme en langage
  naturel ou en pseudocode.
\item Utiliser des instructions conditionnelles dans un algorithme.
\item Utiliser la récursion et l'itération pour effectuer des tâches
  répétitives dans un algorithme.
\item Évaluer la performance d'un algorithme à l'aide de la notation
  $O()$.
\end{objectifs}

Marianne et Alexandre doivent fournir une solution logicielle à un
problème donné. Sitôt la lecture de l'énoncé du problème complétée,
ils se précipitent sur le clavier pour coder leur solution. Rapidement
bloqués, pensant qu'ils butent sur un problème de syntaxe ou
d'utilisation d'une fonction, ils demandent de l'aide. Lorsqu'on leur
demande qu'est-ce qu'ils essaient de faire, quelle est la nature de leur
solution, ils ne peuvent répondre que par un long silence.

Dans mon rôle de professeur, j'ai été témoin de la scène ci-dessus à
de nombreuses reprises. Marianne et Alexandre font une erreur commune
chez les programmeurs --- et pas que chez les débutants: commencer à
coder sans avoir au préalable suffisamment réfléchi à la solution. Il
leur manque un plan clair des étapes à suivre pour résoudre le
problème. Or, ce plan, ou la recette à suivre pour transformer les
entrants d'un problème en une solution, c'est
l'\index{algorithme}\emph{algorithme}. Travailler sans algorithme,
c'est comme naviguer sans boussole.

Puisque programmer requiert des algorithmes, l'étude de la
programmation doit par conséquent s'accompagner de notions
d'\index{algorithmique}\emph{algorithmique}, la science qui étudie les
algorithmes et les structures de données. L'algorithmique est une
discipline riche en techniques ingénieuses et en analyses
mathématiques poussées. Connaitre ses principes de base et étudier
certains algorithmes classiques vous permettra de mieux planifier vos
méthodes de résolution de problème. En effet, un bon algorithme permet
de résoudre en quelques secondes un problème qui pourrait autrement
prendre des années.

L'étude des algorithmes s'accompagne habituellement de celle des
structures de données ou, en d'autres termes, de la manière
d'organiser les données dans un ordinateur. Comme nous le verrons au
\autoref{chap:donnees}, il existe bien différentes structures de
données en R, mais leur réelle mise en œuvre demeure tout à fait
transparente pour les programmeuses et les programmeurs. C'est
pourquoi nous ferons l'impasse sur des notions que l'on retrouve dans
tous les ouvrages classiques d'algorithmique, comme le tableau
(\emph{array}), la liste chainée (\emph{list}), l'arbre (\emph{tree})
ou la table de hachage (\emph{hashtable}).

Le présent chapitre est inspiré de \citet{Sussman:scheme:1996},
\citet{Knuth:ACP:vol1:1997} et \citet{Stephens:algorithms:2013}.


\section{Définition et analogie}
\label{sec:algorithmes:definition}

Un \index{algorithme}algorithme est une procédure de calcul
permettant de résoudre un problème bien spécifié. L'algorithme
explique, de manière non ambigüe et dans un nombre d'opérations fini,
comment, à partir d'entrants, obtenir l'extrant solution du problème.

Pour se mériter un algorithme, un problème doit renfermer une dose
minimale de complexité: personne n'écrit un algorithme pour extraire
le quatrième élément d'un vecteur de données. On suppose que cela fait
partie de la définition de vecteur et que vous savez comment effectuer
l'opération dès lors que vous connaissez le langage de programmation
sous-jacent.

Dans la même veine, un algorithme ne dépend pas du langage de
programmation employé, car il fournit les étapes pour résoudre un
problème, non pas leur mise en œuvre. En d'autres termes, un
algorithme explique \emph{ce qu'il faut faire} pour résoudre un
problème et non \emph{comment le faire}.

Illustrons le concept d'algorithme par l'un de ses plus illustres
représentants: l'\Index{Euclide@Euclide (algorithme)}algorithme
d'Euclide (c.~300 av. J.-C.), une procédure de calcul du plus grand
commun diviseur (PGCD) de deux nombres entiers $m$ et $n$. Selon
\citet{Knuth:ACP:vol2:1997}, il s'agirait de l'un des plus anciens
algorithmes. Il repose sur l'idée que le PGCD de $m$ et $n$ est aussi
le PGCD de $n$ et du reste de la division de $m$ par $n$.

\begin{algorithme}[Algorithme d'Euclide]
  \label{algo:algorithmes:euclide}
  Calculer le plus grand commun diviseur de deux entiers positifs $m$
  et $n$ avec, sans perte de généralité, $n < m$.
  \begin{enumerate}
  \item Diviser $m$ par $n$ et poser $r$ égal au reste de la division.
    (Nous avons alors $0 \leq r < n$.)
  \item Si $r = 0$, retourner la valeur $n$.
  \item Poser $m \leftarrow n$, $n \leftarrow r$, puis retourner à
    l'étape 1.
  \end{enumerate}
\end{algorithme}

Suivons les étapes de l'\index{Euclide@Euclide (algorithme)}algorithme
d'Euclide avec $m = 544$ et $n = 119$. À l'étape~1, nous déterminons
que $544/119 = 4 + 68/119$, donc nous posons $r \leftarrow 68$.
Puisque $r \neq 0$, l'étape~2 de l'algorithme ne s'applique pas. Nous
passons donc à l'étape~3 et nous posons $m \leftarrow 119$ et
$n \leftarrow 68$. De retour à l'étape~1, nous déterminons cette fois
que $r \leftarrow 51$, ce qui mène encore à l'étape~3 et
$m \leftarrow 68$ et $n \leftarrow 51$. Après un autre cycle, nous en
sommes à $m \leftarrow 51$ et $n \leftarrow 17$. Cette fois, nous
posons $r \leftarrow 0$ à l'étape~1 et la procédure s'arrête ensuite à
l'étape~2. Le plus grand commun diviseur de $544$ et $119$ est $17$.
Le \autoref{tab:algorithmes:euclide} dresse le sommaire des opérations
ci-dessus.

\begin{table}
  \centering
  \caption{Calcul du plus grand commun diviseur de $544$ et $119$ avec
    l'algorithme d'Euclide}
  \label{tab:algorithmes:euclide}
  \begin{tabular}{crrr}
    \toprule
    \text{Étape} & $m$ & $n$ & $r$ \\
    \midrule
    $1$ & $544$ & $119$ & $68$ \\
    $3$ & $119$ & $ 68$ & $68$ \\
    $1$ & $119$ & $ 68$ & $51$ \\
    $3$ & $ 68$ & $ 51$ & $51$ \\
    $1$ & $ 51$ & $ 17$ & $ 0$ \\
    $2$ &   --- & $ 17$ &  --- \\
    \bottomrule
  \end{tabular}
\end{table}

\warningbox{L'ordre des opérations à l'étape~3 de
  l'\autoref{algo:algorithmes:euclide} est important. En effet,
  l'action «Poser $m \leftarrow n$, $n \leftarrow r$» est très
  différente de «Poser $n \leftarrow r$, $m \leftarrow n$» puisque
  dans cette dernière, la valeur de $n$ est perdue après la première
  réaffectation. La seconde action est équivalente à «Poser
  $n \leftarrow r$, $m \leftarrow r$».}

Selon \citet[section~1.1]{Knuth:ACP:vol1:1997}, un algorithme devrait
réunir cinq caractéristiques.\footnote{%
  Liste adaptée de la version française fournie par
  \citet{Wikipedia:Algorithme}.}
\begin{enumerate}
\item \emph{Finitude}. Un algorithme doit toujours se terminer après
  un nombre fini d’étapes. Ce nombre peut toutefois devenir très
  grand.
\item \emph{Définition précise}. Chaque étape d'un algorithme doit
  être définie précisément; les actions à réaliser doivent être
  spécifiées rigoureusement et sans ambigüité pour chaque cas.
\item \emph{Entrées}. Un algorithme comporte aucune, une ou plusieurs
  \emph{entrées}, des quantités fournies à l'algorithme avant qu'il ne
  commence ou qui sont allouées dynamiquement durant son exécution.
  Ces entrées proviennent d'un ensemble d'objets bien spécifié. Par
  exemple, l'\autoref{algo:algorithmes:euclide} comporte deux entrées,
  les nombres $m$ et $n$, provenant de l'ensemble des entiers
  positifs.
\item \emph{Sorties}. Un algorithme comporte une ou plusieurs
  \emph{sorties}: des quantités ayant une relation spécifiée avec les
  entrées. L'\autoref{algo:algorithmes:euclide} a une seule sortie,
  soit la valeur $n$ de l'étape~2, le plus grand commun diviseur des
  entrées.
\item \emph{Efficacité}. On s'attend généralement d'un algorithme
  qu'il soit \emph{efficace} dans le sens où toutes les opérations
  qu'il doit accomplir sont suffisamment élémentaires pour pouvoir
  être en principe réalisées dans une durée finie par une personne
  munie de papier et d'un crayon.
\end{enumerate}

Pour bien comprendre la nature et la portée d'un algorithme, il est
assez utile de le comparer à une recette de cuisine. Dans une recette,
les entrées sont les ingrédients, la procédure les diverses étapes et
la sortie, le plat préparé.

Comme un algorithme, une recette repose sur un certain nombre
d'instructions élémentaires qui se passent d'explication, comme
saisir, mijoter, mélanger, etc. La recette est aussi indépendante des
ustensiles utilisés comme l'algorithme l'est du langage de
programmation.

Surtout, surtout, une recette explique les étapes à suivre pour
obtenir le plat voulu, mais, sauf peut-être pour certaines opérations
plus délicates, elle n'explique pas \emph{comment} réaliser chaque
étape. S'il faut ajouter du lait dans un mélange à gâteau, la recette
n'indique pas de sortir le lait du réfrigérateur, de vérifier la date
de péremption et de verser le lait dans la tasse à mesurer jusqu'à la
ligne correspondant à la quantité voulue! Un algorithme ne contient pas
plus d'instructions pour vérifier la validité des entrées ou pour
effectuer des opérations de base comme les opérations arithmétiques,
la lecture de données, l'indiçage d'un vecteur, etc.

En revanche, on s'attend généralement d'un algorithme, qui est destiné
à un ordinateur, qu'il soit plus précis qu'une recette. En effet, un
ordinateur éprouverait beaucoup de difficulté avec une instruction
comme «saler et poivrer au goût»!

\tipbox{C'est une erreur commune dans la rédaction d'un algorithme de
  fournir des détails de mise en œuvre superflus, de surcroit dans un
  langage de programmation bien précis. Lorsque vous écrirez un
  algorithme, demandez-vous si vous êtes en train d'expliquer comment
  mesurer une quantité de lait avec une tasse à mesurer de marque
  ACME.}


\section{Pseudocode}
\label{sec:algorithmes:pseudocode}

Puisqu'il s'agissait du premier exemple d'algorithme complet, j'ai
choisi de présenter l'\autoref{algo:algorithmes:euclide} dans un
format très près du langage naturel.

Une autre manière très populaire d'exprimer les algorithmes consiste à
utiliser du \index{pseudocode}\emph{pseudocode}, un «langage»
similaire à un langage de programmation, mais toujours sans référence
à un langage en particulier. L'écriture en pseudocode permet parfois
de mieux prendre la mesure de la structure et des détails d'un
algorithme avant d'en attaquer la mise en œuvre.

Voici de nouveau l'\index{Euclide@Euclide (algorithme)}algorithme
d'Euclide, présenté cette fois en pseudocode. Portez-y une attention
toute particulière puisqu'il ne s'agit pas d'une transcription directe
de l'\autoref{algo:algorithmes:euclide}. Vérifiez que les deux
versions de l'algorithme retournent bien le même résultat.

\begin{algorithme}[Algorithme d'Euclide, version en pseudocode]
  \label{algo:algorithmes:euclide:pseudo}
  Calculer le plus grand commun diviseur de deux entiers positifs $m$
  et $n$ avec, sans perte de généralité, $n < m$.
\begin{Schunk}
\begin{Verbatim}[commandchars=\\\{\}]
PGCD(entier m, entier n)
  Tant que (n \neq 0)
    r \leftarrow m mod n
    m \leftarrow n
    n \leftarrow r
  Fin Tant que
  Retourner m
Fin PGCD
\end{Verbatim}
\end{Schunk}
\end{algorithme}

\tipbox{L'opération \index{modulo}modulo, notée en mathématiques
  $m \mod n$, représente le reste de la division de $m$ par $n$. Par
  exemple, tel que calculé précédemment, $544 \mod 119 = 68$.}

Il n'existe pas réellement de normes de composition pour le
pseudocode. L'\autoref{algo:algorithmes:euclide:pseudo} reprend
toutefois quelques-unes des conventions les plus usuelles.

La procédure de calcul faisant l'objet de l'algorithme est rédigée
sous forme d'une fonction\footnote{%
  C'est la terminologie de R. Une unité de code qui effectue un
  traitement donné porte divers noms selon le langage de
  programmation: fonction, routine, sous-routine, méthode, procédure,
  sous-procédure\dots}. %
L'algorithme débute par la \index{signature}\emph{signature} de la
fonction (son nom avec le nom et le type de tous les arguments). Le
code est \index{indentation}indenté (décalé vers la droite) après la
signature de la fonction pour montrer qu'il fait partie du corps de la
fonction.

Le bloc délimité par les instructions \code{Tant que
  (\meta{condition})} et \code{Fin Tant que} forme une
\index{boucle}\emph{boucle} dont le contenu --- lui aussi indenté ---
est exécuté tant et aussi longtemps que la \meta{condition} est vraie.

L'instruction \code{Retourner} met immédiatement fin à la fonction,
ici en retournant la valeur \code{m}.

Enfin, on présente le pseudocode comme du code informatique dans une
police non proportionnelle qui préserve l'alignement vertical des
caractères.


\section{Évaluation conditionnelle}
\label{sec:algorithmes:if-else}

Plusieurs algorithmes, même les plus simples, nécessitent de pouvoir
effectuer des opérations différentes selon le résultat d'un test. Par
exemple, pour effectuer le banal calcul de la valeur absolue d'un
nombre, nous devons vérifier si le nombre est positif, négatif ou nul
pour retourner un résultat différent en fonction de la règle suivante:
\begin{equation*}
  \lvert x \rvert =
  \begin{cases}
    x, & \text{si } x > 0 \\
    0, & \text{si } x = 0 \\
    -x, & \text{si } x < 0.
  \end{cases}
\end{equation*}

De telles constructions requièrent des \emph{instructions
  conditionnelles} qui permettent d'effectuer différents calculs selon
le résultat d'une condition booléenne (vraie ou fausse). Il en existe
quatre grands types.

\begin{itemize}
\item Les instructions conditionnelles à un volet de la forme
  \begin{quote}
    Si \meta{condition}, alors \meta{conséquence}
  \end{quote}
  Dans cette construction, la \meta{conséquence} est évaluée seulement
  lorsque la \meta{condition} est vraie. Dans le cas contraire, il ne
  se passe rien et le déroulement de l'algorithme se poursuit
  normalement. Nous avons utilisé une telle instruction conditionnelle
  à l'étape~2 de l'\autoref{algo:algorithmes:euclide}. La locution
  «alors» est souvent omise dans les algorithmes comme, d'ailleurs,
  dans la syntaxe de plusieurs langages de programmation.
\item Les instructions conditionnelles à deux volets de la forme
  \begin{quote}
    Si \meta{condition}, alors \meta{conséquence}, \\
    sinon \meta{alternative}
  \end{quote}
  Cette construction permet d'effectuer un choix entre deux actions:
  ou bien la \meta{condition} est vraie et la \meta{conséquence} est
  évaluée, ou bien la \meta{condition} est fausse et c'est alors
  l'\meta{alternative} qui est évaluée. Le déroulement de l'algorithme
  se poursuit après l'exécution de l'une ou l'autre des deux actions.
\item Les instructions conditionnelles imbriquées de la forme
  \begin{quote}
    Si \meta{condition$_1$}, alors \meta{conséquence$_1$}, \\
    sinon si \meta{condition$_2$}, alors \meta{conséquence$_2$}, \\
    sinon \meta{alternative}
  \end{quote}
  Une instruction conditionnelle imbriquée permet de choisir entre
  trois actions: la \meta{conséquence$_1$} est évaluée lorsque la
  \meta{condition$_1$} est vraie; la \meta{conséquence$_2$} est évaluée
  lorsque la \meta{condition$_1$} est fausse et que la
  \meta{condition$_2$} est vraie; l'\meta{alternative} est évaluée
  lorsque la \meta{condition$_1$} et la \meta{condition$_2$} sont toutes
  les deux fausses. Pour choisir entre plus de trois actions
  possibles, il suffit d'ajouter des clauses «sinon si\dots, alors\dots».
\item Une variante de l'instruction précédente qui facilite la lecture
  lorsqu'il y a un grand nombre d'actions possibles:
  \begin{quote}
    Selon que \\
    \meta{condition$_1$}, alors \meta{action$_1$}; \\
    \meta{condition$_2$}, alors \meta{action$_2$}; \\
    \meta{condition$_3$}, alors \meta{action$_3$}; \\
    \dots \\
    autrement, \meta{action par défaut}
  \end{quote}
  Seule l'\meta{action} correspondant à la \meta{condition} vraie est
  exécutée. Si aucune \meta{condition} n'est vraie, c'est
  l'\meta{action par défaut} qui est exécutée. La terminologie
  française que j'utilise ci-dessus est plutôt rare dans la
  littérature. Dans les langages de programmation, le mot-clé utilisé
  pour ce type de construction est généralement l'un ou l'autre
  de \icode{switch}, \code{case} ou \code{cond}.
\end{itemize}

Afin d'illustrer ce qui précède, exprimons le calcul de la valeur
absolue sous forme d'algorithme. Je fournis d'un coup les versions en
langage naturel et en pseudocode.

\begin{algorithme}
  \label{algo:algorithmes:abs}
  Calculer la valeur absolue d'un nombre réel $x$.

  \noindent
  \begin{minipage}[t]{0.48\linewidth}
    \begin{enumerate}
    \item Si $x > 0$, retourner $x$; sinon si $x = 0$, retourner $0$;
      sinon retourner $-x$.
    \end{enumerate}
  \end{minipage}
  \hfill
  \begin{minipage}[t]{0.48\linewidth}
    \begin{Schunk}
\begin{Verbatim}
abs(réel x)
  Si (x > 0)
    Retourner x
  Sinon si (x = 0)
    Retourner 0
  Sinon
    Retourner -x
Fin abs
\end{Verbatim}
    \end{Schunk}
  \end{minipage}
\end{algorithme}

Pouvons-nous simplifier un peu tout cela? Vous aurez sans doute
remarqué que le cas $x = 0$ ne nécessite pas vraiment de traitement
particulier. Il peut être combiné avec le cas $x > 0$ ou avec le cas
$x < 0$ sans changer le résultat. C'est une leçon que vous devez
retenir: un algorithme simple et efficace ne correspond pas toujours
directement à une formule mathématique.

Voici une version simplifiée de l'algorithme du calcul de la valeur
absolue. J'en ai inversé la logique. Présenté ainsi, l'algorithme met
l'accent sur le fait qu'un traitement ne doit être effectué que pour
les nombres négatifs.

\begin{algorithmebis}
  \label{algo:algorithmes:abs:simplifie}
  Calculer la valeur absolue d'un nombre réel $x$.

  \noindent
  \begin{minipage}[t]{0.48\linewidth}
    \begin{enumerate}
    \item Si $x < 0$, retourner $-x$; sinon retourner $x$.
    \end{enumerate}
  \end{minipage}
  \hfill
  \begin{minipage}[t]{0.48\linewidth}
    \begin{Schunk}
\begin{Verbatim}
abs(réel x)
  Si (x < 0)
    Retourner -x
  Sinon
    Retourner x
Fin abs
\end{Verbatim}
    \end{Schunk}
  \end{minipage}
\end{algorithmebis}

Nous pouvons aussi voir la fonction valeur absolue comme une fonction
identité (qui retourne son argument inchangé), sauf lorsque l'argument
est négatif. Ceci résulte en une autre version de l'algorithme qui se
passe d'une clause «sinon».\footnote{%
  Pensez-y: la clause «sinon» n'est jamais nécessaire lorsque la
  \meta{conséquence} d'une instruction conditionnelle à deux volets
  contient une instruction «Retourner».}

\begin{algorithmebis}
  \label{algo:algorithmes:abs:simplifie:bis}
  Calculer la valeur absolue d'un nombre réel $x$.

  \noindent
  \begin{minipage}[t]{0.48\linewidth}
    \begin{enumerate}
    \item Si $x < 0$, retourner $-x$.
    \item Retourner $x$.
    \end{enumerate}
  \end{minipage}
  \hfill
  \begin{minipage}[t]{0.48\linewidth}
    \begin{Schunk}
\begin{Verbatim}
abs(réel x)
  Si (x < 0)
    Retourner -x
  Retourner x
Fin abs
\end{Verbatim}
    \end{Schunk}
  \end{minipage}
\end{algorithmebis}

Jusqu'à maintenant, nous avons réussi à énoncer la \meta{condition}
des instructions conditionnelles uniquement à partir d'opérateurs
mathématiques simples comme $<$, $=$ et $>$. Pour construire des
conditions composées, il faut avoir recours aux opérateurs logiques
ET, OU et NON.
\begin{itemize}
\item \meta{e$_1$} ET \meta{e$_2$} est vraie lorsque les expressions
  \meta{e$_1$} et \meta{e$_2$} sont toutes deux vraies; sinon la
  condition est fausse.
\item \meta{e$_1$} OU \meta{e$_2$} est vraie lorsque l'une ou l'autre
  des expressions \meta{e$_1$} et \meta{e$_2$} est vraie; la condition
  est fausse uniquement lorsque \meta{e$_1$} et \meta{e$_2$} sont
  toutes deux fausses.
\item NON \meta{e} est vraie lorsque l'expression \meta{e} est fausse,
  et fausse autrement.
\end{itemize}

Par exemple, la condition que l'on écrit mathématiquement $5 < x < 10$
s'exprime ainsi en algorithmique: $x > 5$ ET $x < 10$.


\section{Itération et récursion}
\label{sec:algorithmes:iteration}

La plupart des tâches que nous confions à des ordinateurs sont
répétitives. Après tout, c'est ce dans quoi ils excellent puisque,
contrairement aux humains, ils ne se lassent jamais de répéter
toujours la même opération.

Il existe deux grands types de processus répétitifs: les processus
\emph{itératifs} (ou qui utilisent
l'\Index{iteration@itération}itération) et les processus
\emph{récursifs} (ou qui utilisent la
\Index{recursion@récursion}\index{recursivite@récursivité|see{récursion}}récursion).

Sortons un instant du contexte de la programmation informatique pour
illustrer chaque type de processus. Nous voulons indiquer à une
personne la procédure à suivre pour peinturer une clôture longue de
$n$ planches. (Pour simplifier, nous allons omettre de peinturer les
supports horizontaux.)

Dans l'approche itérative, nous supposons que les planches sont
numérotées de $1$ à $n$, puis nous indiquons à la personne d'appliquer
de la peinture sur la planche $i$, pour $i$ allant de $1$ à $n$.
Transcrites en pseudocode, les instructions vont comme suit.
\begin{Schunk}
\begin{Verbatim}
Peinturer(entier n)
  Pour i de 1 à n
    AppliquerPeintureSurPlanche(i)
  Fin Pour
Fin Peinturer
\end{Verbatim}
\end{Schunk}

La \autoref{fig:algorithmes:peinture-iter} illustre le processus itératif.

\begin{figure}
  \centering
  \setlength{\unitlength}{0.99mm}
  %% planche non peinturée
\savebox{\planche}(2,12)[bl]{%
  \color{gray}
  \polygon*(0,2)(0,3.5)(4,3.5)(4,2)
  \polygon*(0,7.5)(0,9)(4,9)(4,7.5)
  \color{lightgray}
  \moveto(1,0)
  \lineto(1,11)
  \lineto(2,12)
  \lineto(3,11)
  \lineto(3,0)
  \closepath
  \fillpath}

%% planche peinturée
\savebox{\planchep}(2,12)[bl]{%
  \color{gray}
  \polygon*(0,2)(0,3.5)(4,3.5)(4,2)
  \polygon*(0,7.5)(0,9)(4,9)(4,7.5)
  \color[RGB]{204,76,57}
  \moveto(1,0)
  \lineto(1,11)
  \lineto(2,12)
  \lineto(3,11)
  \lineto(3,0)
  \closepath
  \fillpath}

%% peintre
\savebox{\peintre}{%
  \reflectbox{\includegraphics[height=16\unitlength,keepaspectratio]{images/peintre}}}

  \begin{picture}(20,96)(-6,-1)
  \multiput( 0,80)(4,0){5}{\usebox{\planche}}
  \put     (-6,79){\usebox{\peintre}}

  \multiput( 0,64)(4,0){1}{\usebox{\planchep}}
  \multiput( 4,64)(4,0){4}{\usebox{\planche}}
  \put     (-6,63){\usebox{\peintre}}

  \multiput( 0,48)(4,0){2}{\usebox{\planchep}}
  \multiput( 8,48)(4,0){3}{\usebox{\planche}}
  \put     (-2,47){\usebox{\peintre}}

  \multiput( 0,32)(4,0){3}{\usebox{\planchep}}
  \multiput(12,32)(4,0){2}{\usebox{\planche}}
  \put     ( 2,31){\usebox{\peintre}}

  \multiput( 0,16)(4,0){4}{\usebox{\planchep}}
  \multiput(16,16)(4,0){1}{\usebox{\planche}}
  \put     ( 6,15){\usebox{\peintre}}

  \multiput( 0,0)(4,0){5}{\usebox{\planchep}}
  \put     (10,-1){\usebox{\peintre}}
\end{picture}

  \caption{Illustration d'un processus itératif de peinture d'une
    clôture}
  \label{fig:algorithmes:peinture-iter}
\end{figure}

L'approche récursive est plus difficile à bien saisir de prime abord.
Une vieille blague dit d'ailleurs que pour comprendre la récursion, il
faut d'abord comprendre la récursion! Dans cette approche, le
processus de peinturer une clôture de $n$ planches est découpé en deux
étapes: appliquer de la peinture sur la première planche et peinturer
une clôture de $n - 1$ planches, et ce, tant qu'il y a des planches à
peinturer. Examinez bien le pseudocode correspondant ci-dessous.
\begin{Schunk}
\begin{Verbatim}
Peinturer(entier n)
  Si (n = 0)
    Travail terminé
  AppliquerPeintureSurPlanche(1) ET Peinturer(n - 1)
Fin Peinturer
\end{Verbatim}
\end{Schunk}

La personne qui suit les instructions récursives ne peut commencer à
appliquer de la peinture tant qu'elle n'a pas rencontré la condition
d'arrêt (l'extrémité de la clôture) puisque la description du travail
à effectuer dépend d'elle-même. Une fois arrivée au critère d'arrêt,
la personne dispose finalement de toutes les informations pour faire
son travail, qu'elle effectue finalement à rebours. Les processus
récursifs se caractérisent par ces phases d'\index{expansion}expansion
(ou d'accumulation d'information) et de \index{contraction}contraction
(ou d'exécution des instructions). La
\autoref{fig:algorithmes:peinture-recursive} illustre le processus
récursif.

\begin{figure}
  \centering
  \setlength{\unitlength}{0.99mm}
  %% planche non peinturée
\savebox{\planche}(2,12)[bl]{%
  \color{gray}
  \polygon*(0,2)(0,3.5)(4,3.5)(4,2)
  \polygon*(0,7.5)(0,9)(4,9)(4,7.5)
  \color{lightgray}
  \moveto(1,0)
  \lineto(1,11)
  \lineto(2,12)
  \lineto(3,11)
  \lineto(3,0)
  \closepath
  \fillpath}

%% planche peinturée
\savebox{\planchep}(2,12)[bl]{%
  \color{gray}
  \polygon*(0,2)(0,3.5)(4,3.5)(4,2)
  \polygon*(0,7.5)(0,9)(4,9)(4,7.5)
  \color[RGB]{204,76,57}
  \moveto(1,0)
  \lineto(1,11)
  \lineto(2,12)
  \lineto(3,11)
  \lineto(3,0)
  \closepath
  \fillpath}

%% peintre
\savebox{\peintre}{%
  \reflectbox{\includegraphics[height=16\unitlength,keepaspectratio]{images/peintre}}}

  \begin{picture}(64,176)(-6,-1)
  \multiput( 1,160)(4,0){5}{\usebox{\planche}}
  \put     (-6,159){\usebox{\peintre}}

  \multiput( 1,144)(11,0){1}{\usebox{\planche}}
  \multiput(12,144)(4,0){4}{\usebox{\planche}}
  \put     ( 5,143){\usebox{\peintre}}

  \multiput( 1,128)(11,0){2}{\usebox{\planche}}
  \multiput(23,128)(4,0){3}{\usebox{\planche}}
  \put     (16,127){\usebox{\peintre}}

  \multiput( 1,112)(11,0){3}{\usebox{\planche}}
  \multiput(34,112)(4,0){2}{\usebox{\planche}}
  \put     (27,111){\usebox{\peintre}}

  \multiput( 1,96)(11,0){4}{\usebox{\planche}}
  \multiput(45,96)(6,0){1}{\usebox{\planche}}
  \put     (38,95){\usebox{\peintre}}

  \multiput( 1,80)(11,0){5}{\usebox{\planche}}
  \put     (49,79){\usebox{\peintre}}

  \multiput( 1,64)(11,0){4}{\usebox{\planche}}
  \multiput(45,64)(6,0){1}{\usebox{\planchep}}
  \put     (38,63){\usebox{\peintre}}

  \multiput( 1,48)(11,0){3}{\usebox{\planche}}
  \multiput(34,48)(4,0){2}{\usebox{\planchep}}
  \put     (27,47){\usebox{\peintre}}

  \multiput( 1,32)(11,0){2}{\usebox{\planche}}
  \multiput(23,32)(4,0){3}{\usebox{\planchep}}
  \put     (16,31){\usebox{\peintre}}

  \multiput( 1,16)(11,0){1}{\usebox{\planche}}
  \multiput(12,16)(4,0){4}{\usebox{\planchep}}
  \put     ( 5,15){\usebox{\peintre}}

  \multiput( 1,0)(4,0){5}{\usebox{\planchep}}
  \put     (-6,-1){\usebox{\peintre}}
\end{picture}

  \caption[Illustration d'un processus récursif de peinture d'une
  clôture]{Illustration d'un processus récursif de peinture d'une
    clôture. La séparation d'une planche symbolise la mise en attente
    de cette étape.}
  \label{fig:algorithmes:peinture-recursive}
\end{figure}

Nous pouvons aussi comparer les deux types de processus en observant
ce qui se passerait si la personne devait interrompre son travail pour
le confier à quelqu'un d'autre. Dans l'approche itérative, elle
pourrait simplement indiquer qu'elle en était à appliquer de la
peinture sur la planche $m$ dans le processus de peinturer une clôture
de $n$ planches. Le travail pourrait ensuite se poursuivre à partir de
la planche $m + 1$. Dans l'approche récursive, cependant, la personne
ne pourrait uniquement dire qu'elle en était à peinturer une clôture
de $m$ planches, car elle aurait aussi bien pu se trouver dans la
phase d'expansion que dans celle de contraction. Elle devrait donc
aussi fournir des informations --- qu'elle a mémorisées --- sur
l'«état» dans lequel elle se trouvait au moment d'interrompre son
travail.

À ce stade, vous pourriez vous demander pourquoi donc parler de
\index{recursion@récursion}récursion puisque ça semble si compliqué,
voire inefficace. Le fait est que la solution de plusieurs problèmes
s'exprime beaucoup plus simplement de manière récursive que de manière
itérative. Pensez au tri d'une main de cartes: il s'agit de placer la
plus petite carte au début de la main, puis de trier le reste de la
main, et ainsi de suite jusqu'à ce que la main soit triée au complet.
L'\index{Euclide@Euclide (algorithme)}algorithme d'Euclide est
également récursif. Examinez de nouveau
l'\autoref{algo:algorithmes:euclide}: l'étape~2 n'est rien d'autre que
le critère d'arrêt et l'étape~3, le rappel de la procédure avec des
arguments différents.

Dans l'absolu, les procédures récursives sont moins efficaces en
termes de temps de calcul et de consommation de ressources, notamment
parce qu'elles nécessitent une \index{pile}pile (\emph{stack}) pour
accumuler les informations durant la phase d'expansion. En pratique,
toutefois, ces inconvénients tendent à diminuer ou à disparaitre avec
l'optimisation des compilateurs de plusieurs langages de programmation.
De plus, la propriété de
\index{recursion@récursion!terminale}récursion terminale
(\emph{tail-recursion}) de certains langages, dont
\index{Scheme}Scheme, fait en sorte que certaines procédures
récursives sont en tous points équivalentes à des procédures
itératives. Nous n'irons pas plus loin, mais si le sujet vous
intéresse, je vous recommande de lire la section~1.2.1 de
\citet{Sussman:scheme:1996}. L'article de
\citet{Wikipedia:Recursion_terminale} constitue également un bon point
de départ pour explorer le sujet plus en profondeur.

Je ne pourrais prétendre avoir écrit un ouvrage d'introduction à la
programmation sans traiter de l'exemple classique quand il s'agit de
comparer \index{iteration@itération}itération et
\index{recursion@récursion}récursion: le calcul de la
\Index{factorielle}factorielle. La fonction factorielle est définie
par
\begin{equation*}
  n! = n \cdot (n - 1) \cdot (n - 2) \cdots 3 \cdot 2 \cdot 1.
\end{equation*}
Or, puisque
\begin{equation*}
  n! = n \cdot [(n - 1) \cdot (n - 2) \dots 3 \cdot 2 \cdot 1] =
  n \cdot (n - 1)!,
\end{equation*}
la fonction est récursive et une manière toute naturelle de calculer
la factorielle d'un nombre utilise la récursion. Exprimons cela  sous
forme d'algorithme. Le pseudocode sied un peu mieux aux algorithmes
récursifs.

\begin{algorithme}[Factorielle, approche récursive]
  \label{algo:algorithmes:factorielle-recursive}
  Calculer la \index{factorielle}factorielle d'un entier positif $n$.
  \begin{Schunk}
\begin{Verbatim}[commandchars=\\\{\}]
factorial(entier n)
  Si (n = 1)
    Retourner 1
  Retourner n \times factorial(n - 1)
Fin factorial
\end{Verbatim}
  \end{Schunk}
\end{algorithme}

La \autoref{fig:algorithmes:factorial} illustre le processus de calcul
récursif de la valeur de $6!$. Remarquez comme la structure est
similaire à celle du processus de peinture d'une clôture de la
\autoref{fig:algorithmes:peinture-recursive}.

\begin{figure}
  \centering
  \begin{minipage}{0.67\linewidth}
\begin{Verbatim}[commandchars=\\\{\}]
factorial(6)
6 \times factorial(5)
6 \times (5 \times factorial(4))
6 \times (5 \times (4 \times factorial(3)))
6 \times (5 \times (4 \times (3 \times factorial(2))))
6 \times (5 \times (4 \times (3 \times (2 \times factorial(1)))))
6 \times (5 \times (4 \times (3 \times (2 \times 1))))
6 \times (5 \times (4 \times (3 \times 2)))
6 \times (5 \times (4 \times 6))
6 \times (5 \times 24)
6 \times 120
720
\end{Verbatim}
  \end{minipage}
  \caption{Processus récursif du calcul de $6!$}
  \label{fig:algorithmes:factorial}
\end{figure}

Nous pouvons également voir $n!$ comme le résultat du produit de $1$
par $2$, puis le produit du résultat par $3$, puis par $4$, et ainsi
de suite jusqu'à $n$. Avec cette approche, nous devons enregistrer le
résultat courant du produit ainsi qu'un
\index{compteur}\emph{compteur} qui nous permet de savoir où nous en
sommes dans le calcul. Cela devrait vous rappeler l'approche itérative
de la peinture d'une clôture, car c'est bien ce dont il s'agit: le
calcul de la \index{factorielle}factorielle de manière itérative.
Exprimons cette fois l'algorithme en langage naturel.

\begin{algorithme}[Factorielle, approche itérative]
  \label{algo:algorithmes:factorielle-iterative}
  Calculer la \index{factorielle}factorielle d'un entier positif $n$.
  \begin{enumerate}
  \item Poser $p \leftarrow i \leftarrow 1$.
  \item Si $i = n$, retourner la valeur $p$.
  \item Incrémenter le compteur $i$: $i \leftarrow i + 1$.
  \item Poser $p \leftarrow p \times i$, puis retourner à l'étape 2.
  \end{enumerate}
\end{algorithme}

La \autoref{fig:algorithmes:fact-iter} illustre le processus de calcul
itératif de la valeur de $6!$. Sa structure est cette fois similaire à
celle de \autoref{fig:algorithmes:peinture-iter}. Remarquez comment
les valeurs de $p$ et de $i$ enregistrent l'état du processus, ce qui
nous permettrait de l'interrompre et de le relancer à tout moment.

\begin{figure}
  \centering
  \begin{minipage}{0.25\linewidth}
\begin{Verbatim}[commandchars=\\\{\}]
factorial(6)
i \leftarrow 1; p \leftarrow 1
i \leftarrow 2; p \leftarrow 2
i \leftarrow 3; p \leftarrow 6
i \leftarrow 4; p \leftarrow 24
i \leftarrow 5; p \leftarrow 120
i \leftarrow 6; p \leftarrow 720
720
\end{Verbatim}
  \end{minipage}
  \caption{Processus itératif du calcul de $6!$}
  \label{fig:algorithmes:fact-iter}
\end{figure}


\section{Nombre d'opérations}
\label{sec:algorithmes:bigO}

Supposons que vous devez calculer, pour un jeu de données quelconque,
l'écart moyen des données supérieures à $d$ par rapport à cette
valeur. Le résultat est l'\emph{espérance résiduelle empirique} des
données évaluée à $d$. L'algorithme de haut niveau suivant permet
d'effectuer ce calcul.

\begin{algorithme}
  \label{algo:algorithmes:emrl}
  Calculer l'espérance résiduelle empirique à $d$ de données $x_1, x_2,
  \dots, x_n$.
  \begin{enumerate}
  \item Sélectionner les données supérieures à $d$.
  \item Soustraire $d$ de chaque valeur retenue à l'étape 1.
  \item Retourner la moyenne des valeurs obtenues à l'étape 2.
  \end{enumerate}
\end{algorithme}

Nous pouvons aussi concevoir une variante de l'algorithme ci-dessus en
y apportant une toute petite modification.

\begin{algorithmebis}
  \label{algo:algorithmes:emrl:bis}
  Calculer l'espérance résiduelle empirique à $d$ de données $x_1, x_2,
  \dots, x_n$.
  \begin{enumerate}
  \item Sélectionner les données supérieures à $d$.
  \item Effectuer la moyenne des valeurs obtenues à l'étape 1.
  \item Retourner la différence entre la moyenne obtenue à l'étape 2
    et $d$.
  \end{enumerate}
\end{algorithmebis}

Mathématiquement, les deux approches sont tout à fait équivalentes.
Cependant, avec l'\autoref{algo:algorithmes:emrl}, nous effectuons une
soustraction pour chaque valeur supérieure à $d$ que compte le jeu de
données. Avec l'\autoref{algo:algorithmes:emrl:bis}, nous n'effectuons
qu'une seule soustraction. Si le jeu de données compte un million
d'entrées supérieures à $d$, c'est un million de soustractions contre
une seule. En temps de calcul, cela ne représente qu'un écart de
quelques centièmes de secondes sur un ordinateur moderne, mais ces
fractions de secondes peuvent finir par faire une différence
importante lorsque la taille des jeux de données augmente ou lorsque
l'opération se répète un très grand nombre de fois.

Le dénombrement du nombre d'opérations requis par un algorithme est un
aspect important de leur analyse. Il se fait généralement en notation
$O(f(n))$ qui signifie «de l'ordre de $f(n)$», où $f(n)$ est une
fonction d'un paramètre $n$ qui mesure la taille du problème. Dans
l'exemple ci-dessus, $n$ est le nombre d'entrées dans le jeu de
données. Pour un algorithme de calcul du produit de deux matrices, $n$
pourrait être le nombre de lignes des matrices.

Attardons-nous à un second exemple intéressant: l'élévation d'une
valeur $b$ à la puissance $n$. L'approche la plus intuitive
consisterait à utiliser la relation
\begin{equation*}
  b^n = b (b (b \cdots (b))) = b \cdot b^{n - 1},
\end{equation*}
ce qui mène à un premier algorithme.

\begin{algorithme}
  \label{algo:algorithmes:puissance}
  Élever un nombre $b$ à la puissance entière positive $n$.
  \begin{Schunk}
\begin{Verbatim}[commandchars=\\\{\}]
Puissance(nombre réel b, entier positif n)
  Si (n = 0)
    Retourner 1
  Retourner b \times Puissance(b, n - 1)
Fin Puissance
\end{Verbatim}
  \end{Schunk}
\end{algorithme}

Vous aurez tout de suite identifié que
l'\autoref{algo:algorithmes:puissance} est récursif. Le nombre
d'opérations effectué avec cet algorithme est directement
proportionnel à $n$. Par exemple, il requiert $5$ opérations pour
élever un nombre à la puissance $6$. On dit donc que l'ordre de
grandeur de l'algorithme est $O(n)$.

L'\autoref{algo:algorithmes:puissance} peut convenir lorsque $n$ est
petit, mais peut-on faire mieux pour une «grande» valeur de $n$?
Imaginez que vous ne disposez que d'une calculatrice munie des
opérations arithmétiques de base et du carré, et que vous devez élever
un nombre à la puissance, disons, $21$. Comment feriez-vous pour
réduire le nombre d'opérations à entrer au clavier?

\notebox{Cet exemple est moins fantaisiste qu'il n'y parait! Jusqu'au
  milieu des années 1990, les étudiants en actuariat ne disposaient
  que de ce type de calculatrice pour les examens professionnels. Il
  s'agissait d'un modèle de Texas Instruments affectueusement surnommé
  «TI-zéro».}

Vous avez pensé à un algorithme? Votre idée consiste fort probablement
à diviser la puissance par deux autant de fois que possible et à
élever au carré par la suite. Pour $n = 21$, cela donne
\begin{align*}
  b^{21}
  &= b (b^{20}) \\
  & = b (b^{10})^2 \\
  &= b ((b^5)^2)^2 \\
  & = b ((b (b^4))^2)^2 \\
  &= b ((b (b^2)^2)^2)^2.
\end{align*}
Ce calcul se traduit en un nouvel algorithme récursif comme suit.

\begin{algorithmebis}
  \label{algo:algorithmes:puissance:bis}
  Élever un nombre $b$ à la puissance entière positive $n$.
  \begin{Schunk}
\begin{Verbatim}[commandchars=\\\{\}]
Puissance(nombre réel b, entier positif n)
  Si (n = 0)
    Retourner 1
  Si (n est pair)
    Retourner (Puissance(b, n/2))^2
  Si (n est impair)
   Retourner b \times Puissance(b, n - 1)
Fin Puissance
\end{Verbatim}
  \end{Schunk}
\end{algorithmebis}

Là où l'\autoref{algo:algorithmes:puissance} requiert $20$ opérations
pour élever un nombre à la puissance $21$,
l'\autoref{algo:algorithmes:puissance:bis} n'en demande que $6$. Comme
la puissance est divisée par deux à répétition, le nombre d'opérations
est proportionnel au logarithme (en base deux) de $n$. Pour vous en
convaincre, observez que le calcul de $b^{2n}$ effectue une seule
opération de plus que le calcul de $b^n$. L'ordre de grandeur de
l'algorithme est donc $O(\log n)$.

La mesure d'ordre de grandeur est très approximative. Par exemple, un
processus nécessitant $n^2$ étapes, un processus en nécessitant
$1000n^2$ et un autre en nécessitant $3n^2 + 5n + 17$ ont tous un
ordre de grandeur $O(n^2)$ puisque les constantes ne sont pas prises
en compte.

Quelques ordres de grandeur surviennent plus souvent dans l'étude des
algorithmes. En voici une liste partielle, question de vous offrir un
peu de perspective si vous lisez que la performance d'un algorithme
d'ordre $O(n^3)$ est tout à fait acceptable.

\begin{itemize}
\item $O(1)$. Un algorithme d'ordre $O(1)$ prend toujours le même
  temps pour effectuer ses calculs, peu importe la taille du problème.
  Il va sans dire que de tels algorithmes sont excessivement rares. Ils
  effectuent en général des tâches très simples.
\item $O(\log n)$. Tel que mentionné précédemment, l'ordre de grandeur
  $O(\log n)$ survient dès lors que la taille du problème est divisée
  par deux à chaque étape.
\item $O(\sqrt{n})$. Les algorithmes d'ordre $O(\sqrt{n})$ sont rares.
  La fonction racine carrée croit lentement, mais néanmoins plus
  rapidement que la fonction logarithme.
\item $O(n)$. Nous avons rencontré un exemple d'algorithme $O(n)$
  ci-dessus. Les algorithmes avec cette performance demeurent
  acceptables étant donné le rythme de croissance somme toute
  raisonnable de la fonction identité.
\item $O(n \log n)$. Les algorithmes qui répètent $n$ fois un calcul
  d'ordre $O(\log n)$ ont une performance d'ordre $O(n \log n)$.
  C'est le cas de plusieurs algorithmes de tri.
\item $O(n^2)$. Les algorithmes qui effectuent un calcul pour toutes
  les paires d'entrées ont une performance $O(n^2)$. Ici, ça commence
  à devenir lent.
\item $O(n!)$. Attention, danger. La fonction factorielle croit si
  rapidement qu'un algorithme ayant ce type de performance ne peut
  résoudre que de petits problèmes.
\end{itemize}

Pour en savoir plus sur les principes de base du dénombrement des
opérations d'un algorithme, consultez
\citet[chapitre~1]{Stephens:algorithms:2013}.



\section{Exercices}
\label{sec:algorithmes:exercices}

\Opensolutionfile{solutions}[solutions-algorithmes]

\begin{Filesave}{solutions}
\section*{Chapitre \ref*{chap:algorithmes}}
\addcontentsline{toc}{section}{Chapitre \protect\ref*{chap:algorithmes}}

\begingroup
%% modifications locales au style des listings
% \lstset{%
%   frame=single,
%   numbers=none,
%   basicstyle=\normalfont\ttfamily\NoAutoSpacing}

\end{Filesave}

\begin{exercice}
  Écrire l'algorithme d'une fonction \code{quadroots} qui prend en
  arguments les coefficients (réels) $a$, $b$ et $c$ d'un polynôme du
  second degré $a x^2 + b x + c$ et qui calcule les racines du
  polynôme avec la formule bien connue
  \begin{equation*}
    \frac{-b \pm \sqrt{b^2 - 4ac}}{2a}.
  \end{equation*}
  La fonction doit retourner deux, une ou aucune racine selon que le
  discriminant $b^2 - 4ac$ est positif, nul ou négatif, dans l'ordre.
  \begin{sol}
    Voici des versions en langage naturel et en pseudocode.

    \begin{algorithme}
      Calculer les racines du polynôme $a x^2 + b x + c$ à partir des
      coefficients réels $a$, $b$ et $c$.
      \begin{enumerate}[1.]
      \item Poser $d \leftarrow b^2 - 4ac$.
      \item Si $d < 0$, retourner un résultat vide.
      \item Si $d = 0$, retourner $-b/(2a)$.
      \item Retourner $(-b - \sqrt{d})/(2a)$ et
        $(-b + \sqrt{d})/(2a)$.
      \end{enumerate}
    \end{algorithme}

    \begin{algorithme}
      Calculer les racines du polynôme $a x^2 + b x + c$ à partir des
      coefficients réels $a$, $b$ et $c$.
      \begin{Schunk}
\begin{Verbatim}[commandchars=\\\{\}]
quadroots(réel a, réel b, réel c)
  d \leftarrow b^2 - 4 \times a \times c
  Si (d < 0)
    Retourner NULL
  Si (d = 0)
    Retourner -b/(2 \times a)
  x \leftarrow (-b - sqrt(d))/(2 \times a)
  y \leftarrow (-b + sqrt(d))/(2 \times a)
  Retourner x et y
Fin quadroots
\end{Verbatim}
      \end{Schunk}
    \end{algorithme}
  \end{sol}
\end{exercice}

\begin{exercice}
  Expliquer dans vos mots pourquoi la clause «sinon» n'est jamais
  nécessaire lorsque la \meta{conséquence} d'une instruction
  conditionnelle à deux volets contient une instruction «Retourner».
  \begin{sol}
    Si la \meta{condition} d'une instruction conditionnelle est
    vraie, la fonction se termine avec l'instruction «Retourner».
    Toutes les instructions qui suivent la clause «Si
    \meta{condition}, alors \dots» ne peuvent donc être évaluées que
    lorsque la \meta{condition} est fausse. Inutile de les protéger
    par une clause «sinon».

    Par exemple, supposons qu'une fonction contient les instructions
    suivantes:
    \begin{Schunk}
\begin{Verbatim}[commandchars=\\\{\}]
Si \meta{condition}
  \meta{conséquence}
sinon
  \meta{alternative}
\meta{instruction}
\end{Verbatim}
    \end{Schunk}
    Avec une telle structure, l'\meta{instruction} sera évaluée en
    tout temps, que la \meta{condition} soit vraie ou fausse. Or, si la
    \meta{conséquence} contient «Retourner», tant l'\meta{alternative}
    que l'\meta{instruction} ne seront évaluée que lorsque la
    \meta{condition} est fausse. La formulation ci-dessous est donc
    tout à fait équivalente.
    \begin{Schunk}
\begin{Verbatim}[commandchars=\\\{\}]
Si \meta{condition}
  \meta{conséquence}
\meta{alternative}
\meta{instruction}
\end{Verbatim}
    \end{Schunk}
  \end{sol}
\end{exercice}

\begin{exercice}
  Plusieurs algorithmes comportent une étape d'échange de deux
  valeurs, que l'on pourrait symboliser par $m \leftrightarrow n$.
  \begin{enumerate}
  \item Expliquer comment effectuer l'opération $m \leftrightarrow n$
    sans perdre l'une ou l'autre des valeurs $m$ et $n$.
    \emph{Indice}: il faut avoir recours à une troisième variable
    temporaire, disons $t$.
  \item Expliquer comment réordonner les valeurs $(a, b, c, d)$ en
    $(b, c, d, a)$ en effectuant un minimum de remplacements.
  \end{enumerate}
  \begin{sol}
    \begin{enumerate}
    \item Les étapes à suivre sont, dans l'ordre: $t \leftarrow m$, $m
      \leftarrow n$, $n \leftarrow t$.
    \item Une solution qui ne requiert que cinq remplacements est:
      $t \leftarrow a$, $a \leftarrow b$, $b \leftarrow c$,
      $c \leftarrow d$, $d \leftarrow t$.
    \end{enumerate}
  \end{sol}
\end{exercice}

\begin{exercice}
  Soit $p$ et $q$ des expressions logiques. Compléter les tables de
  vérité ci-dessous.
  \begin{multicols}{2}
    \begin{enumerate}
    \item
      \begin{tabular}{ccc}
        \toprule
        $p$ & $q$ & $p$ ET $q$ \\
        \midrule
        V & V \\
        V & F \\
        F & V \\
        F & F \\
        \bottomrule
      \end{tabular}
      \columnbreak
    \item
      \begin{tabular}{ccc}
        \toprule
        $p$ & $q$ & $p$ OU $q$ \\
        \midrule
        V & V \\
        V & F \\
        F & V \\
        F & F \\
        \bottomrule
      \end{tabular}
    \end{enumerate}
  \end{multicols}
  \begin{sol}
    Nous pouvons compléter les tables de vérité directement à partir
    des définitions des opérateurs ET et OU fournies à la
    \autoref{sec:algorithmes:if-else}:
    \begin{multicols}{2}
      \begin{enumerate}
      \item
        \begin{tabular}{ccc}
          \toprule
          $p$ & $q$ & $p$ ET $q$ \\
          \midrule
          V & V & V \\
          V & F & F \\
          F & V & F \\
          F & F & F \\
          \bottomrule
        \end{tabular}
        \columnbreak
      \item
        \begin{tabular}{ccc}
          \toprule
          $p$ & $q$ & $p$ OU $q$ \\
          \midrule
          V & V & V \\
          V & F & V \\
          F & V & V \\
          F & F & F \\
          \bottomrule
        \end{tabular}
      \end{enumerate}
    \end{multicols}
  \end{sol}
\end{exercice}

\begin{exercice}
  La notation en algèbre booléenne est plus compacte et lisible que
  celle utilisée à la \autoref{sec:algorithmes:if-else}. Si $p$ et $q$
  sont des expressions logiques, on représente généralement $p$ ET $q$
  par $p \wedge q$, $p$ OU $q$ par $p \vee q$, et NON $p$ par
  $\overline{p}$.

  Soit $p$, $q$ et $r$ des expressions logiques. Démontrer à l'aide de
  tables de vérité les identités logiques suivantes.
  \begin{enumerate}
  \item $p \wedge (q \vee r) = (p \wedge q) \vee (p \wedge r)$
  \item $p \vee (q \wedge r) = (p \vee q) \wedge (p \vee r)$
  \item $\overline{p \wedge q} = \overline{p} \vee \overline{q}$
  \item $\overline{p \vee q} = \overline{p} \wedge \overline{q}$
  \item $(p \wedge q) \vee (\overline{p} \wedge r) = (p \wedge q) \vee (\overline{p} \wedge
    r) \vee (q \wedge r)$
  \end{enumerate}
  \begin{sol}
    Tel que proposé dans l'énoncé, tous les résultats se démontrent à
    l'aide de tables de vérité.
    \begin{enumerate}
    \item
      \begin{tabularx}{\linewidth}{*{7}{c}>{\centering\arraybackslash}X}
        \toprule
        $p$ & $q$ & $r$ & $p \wedge q$ & $p \wedge r$ & $q \vee r$ & $p \wedge (q \vee r)$ & $(p \wedge q) \newline\vee (p \wedge r)$
        \\
        \midrule
        V & V & V & V & V & V & V & V \\
        V & V & F & V & F & V & V & V \\
        V & F & V & F & V & V & V & V \\
        V & F & F & F & F & F & F & F \\
        F & V & V & F & F & V & F & F \\
        F & V & F & F & F & V & F & F \\
        F & F & V & F & F & V & F & F \\
        F & F & F & F & F & F & F & F \\
        \bottomrule
      \end{tabularx}
    \item
      \begin{tabularx}{\linewidth}{*{7}{c}>{\centering\arraybackslash}X}
        \toprule
        $p$ & $q$ & $r$ & $p \vee q$ & $p \vee r$ & $q \wedge r$ & $p \vee (q \wedge r)$ & $(p \vee q) \newline\wedge (p \vee r)$
        \\
        \midrule
        V & V & V & V & V & V & V & V \\
        V & V & F & V & V & F & V & V \\
        V & F & V & V & V & F & V & V \\
        V & F & F & V & V & F & V & V \\
        F & V & V & V & V & V & V & V \\
        F & V & F & V & F & F & F & F \\
        F & F & V & F & V & F & F & F \\
        F & F & F & F & F & F & F & F \\
        \bottomrule
      \end{tabularx}
    \item
      \begin{tabular}{*{6}{c}}
        \toprule
        $p$ & $q$ & $\overline{p}$ & $\overline{q}$ & $\overline{p \wedge q}$ & $\overline{p} \vee \overline{q}$ \\
        \midrule
        V & V & F & F & F & F \\
        V & F & F & V & V & V \\
        F & V & V & F & V & V \\
        F & F & V & V & V & V \\
        \bottomrule
      \end{tabular}
    \item
      \begin{tabular}{*{6}{c}}
        \toprule
        $p$ & $q$ & $\overline{p}$ & $\overline{q}$ & $\overline{p \vee q}$ & $\overline{p} \wedge \overline{q}$ \\
        \midrule
        V & V & F & F & F & F \\
        V & F & F & V & F & F \\
        F & V & V & F & F & F \\
        F & F & V & V & V & V \\
        \bottomrule
      \end{tabular}
    \item
      \begin{tabularx}{\linewidth}{*{6}{c}>{\centering\arraybackslash}X>{\centering\arraybackslash}X}
        \toprule
        $p$ & $q$ & $r$ & $p \wedge q$ & $\overline{p} \wedge r$ & $q \wedge r$ & $(p \wedge q) \newline\vee (\overline{p} \wedge r)$ & $(p \wedge q) \newline\vee (\overline{p} \wedge r) \newline\vee (q \wedge r)$
        \\
        \midrule
        V & V & V & V & F & V & V & V \\
        V & V & F & V & F & F & V & V \\
        V & F & V & F & F & F & F & F \\
        V & F & F & F & F & F & F & F \\
        F & V & V & F & V & V & V & V \\
        F & V & F & F & F & F & F & F \\
        F & F & V & F & V & F & V & V \\
        F & F & F & F & F & F & F & F \\
        \bottomrule
      \end{tabularx}
    \end{enumerate}
  \end{sol}
\end{exercice}

\begin{exercice}
  Écrire une version récursive présentée en pseudocode de l'algorithme
  d'\index{Euclide@Euclide (algorithme)}Euclide.
  \begin{sol}
    Tel que mentionné à la \autoref{sec:algorithmes:iteration}, les
    algorithmes récusifs s'expriment généralement mieux en pseudocode.
    L'algorithme d'\index{Euclide@Euclide (algorithme)}Euclide est
    particulièrement simple exprimé ainsi.
    \begin{Schunk}
\begin{Verbatim}
PGCD(entier m, entier n)
  Si (n = 0) retourner m
  Retourner PGCD(n, m mod n)
Fin PGCD
\end{Verbatim}
    \end{Schunk}
  \end{sol}
\end{exercice}

\begin{exercice}
  La suite de \index{Fibonacci}Fibonacci est une suite de
  nombres entiers très connue. Les deux premiers termes de la suite
  sont $0$ et $1$ et tous les autres sont la somme des deux termes
  précédents. Mathématiquement, les valeurs de la suite de Fibonacci
  sont données par la fonction
  \begin{align*}
    f(0) &= 0 \\
    f(1) &= 1 \\
    f(n) & = f(n - 1) + f(n - 2), \quad n \geq 2.
  \end{align*}
  Le quotient de deux termes successifs converge vers
  $\phi = (1 + \sqrt{5})/2$, le nombre d'or.\footnote{%
    On prête au nombre d'or toutes sortes de propriétés, certaines
    même d'ordre mystique. Il est en tout cas au cœur d'un des romans
    les plus vendus de tous les temps, \emph{Da~Vinci Code} (Dan
    Brown, 2003).} %
  Le calcul de $f(n)$ se prête tout naturellement à une mise en œuvre
  récursive.
  \begin{enumerate}
  \item Proposer un algorithme récursif pour calculer la valeur $f(n)$
    de la suite de Fibonacci.
  \item Calculer les 10 premiers nombres de Fibonacci.
  \item L'algorithme récursif direct de la suite de Fibonacci est
    terriblement inefficace puisqu'il répète plusieurs calculs un
    grand nombre de fois. Pour vous en convaincre, écrivez au long
    tous les calculs requis par l'algorithme pour calculer $f(5)$.
  \item Proposer un algorithme itératif en langage naturel pour
    calculer les nombres de Fibonacci.
  \end{enumerate}
  \begin{sol}
    \begin{enumerate}
    \item L'algorithme ci-dessous permet de calculer les nombres de
      Fibonacci de manière récursive.
      \begin{algorithme}[Suite de Fibonacci]
        \label{algo:algorithmes:Fibonacci}
        Calculer la valeur $f(n)$ de la suite de
        \index{Fibonacci}Fibonacci pour un entier non négatif $n$.
        \begin{Schunk}
\begin{Verbatim}
Fib(entier n)
  Si (n = 0), retourner 0
  Si (n = 1), retourner 1
  Retourner Fib(n - 2) + Fib(n - 1)
Fin Fib
\end{Verbatim}
        \end{Schunk}
      \end{algorithme}
    \item Les 10 premiers nombres de Fibonacci sont:
      $0, 1, 1, 2, 3, 5, 8, 13, 21, 34$.
    \item En décomposant la méthode de calcul des nombres de Fibonacci
      employée par l'\autoref{algo:algorithmes:Fibonacci}, on obtient:
      \begin{align*}
        f(5)
        &= f(4) + f(3) \\
        &= (f(3) + f(2)) + (f(2) + f(1)) \displaybreak[0]\\
        &= ((f(2) + f(1)) + (f(1) + f(0))) \\
        &\phantom{=} + ((f(1) + f(0)) + f(1)) \displaybreak[0]\\
        &= (((f(1) + f(0)) + f(1)) + (f(1) + f(0))) \\
        &\phantom{=} + ((f(1) + f(0)) + f(1)).
      \end{align*}
      Plusieurs calculs sont redondants. Il faut calculer $f(1)$ pas
      moins de cinq fois et le calcul de $f(3)$ dès la première
      égalité ci-dessus n'est qu'une reprise d'un calcul effectué par
      ailleurs. En fait, cet algorithme est souvent utilisé comme cas
      pathologique dans les ouvrages d'algorithmique.
    \item L'algorithme ci-dessous permet de calculer les nombres de
      Fibonacci de manière itérative.
      \begin{algorithme}[Suite de Fibonacci]
        \label{algo:algorithmes:Fibonacci-iter}
        Calculer la valeur $f(n)$ de la suite de
        \index{Fibonacci}Fibonacci pour un entier non négatif $n$.
        \begin{enumerate}[1.]
        \item Poser $a \leftarrow 0$, $b \leftarrow 1$ et
          $i \leftarrow 2$.
        \item Si $n = 0$, retourner $a$.
        \item Si $n = 1$, retourner $b$.
        \item Poser $r \leftarrow a + b$.
        \item Si $i = n$, retourner la valeur $r$.
        \item Poser $a \leftarrow b$ et $b \leftarrow r$,
        \item Incrémenter le compteur $i$: $i \leftarrow i + 1$, puis
          retouner à l'étape 4.
        \end{enumerate}
      \end{algorithme}
    \end{enumerate}
  \end{sol}
\end{exercice}

\begin{exercice}
  L'algorithme suivant permet de calculer la valeur de la fonction de
  \index{Ackermann (fonction)}Ackermann $A(x, y)$ pour des entiers non
  négatifs $x$ et $y$ telle que définie dans
  \citet{Sussman:scheme:1996}.

  \begin{Schunk}
\begin{Verbatim}[commandchars=\\\{\}]
A(entier x, entier y)
  Si (y = 0)
    Retourner 0
  Sinon Si (x = 0)
    Retourner 2 \times y
  Sinon Si (y = 1)
    Retourner 2
  Sinon
    Retourner A(x - 1, A(x, y - 1))
Fin A
\end{Verbatim}
  \end{Schunk}

  Calculer la valeur de $A(1, 10)$.
  \begin{sol}
    On a
    \begin{align*}
      A(1, 10)
      &= A(0, A(1, 9)) \\
      &= A(0, A(0, A(1, 8))) \\
      &= A(0, A(0, A(0, A(1, 7)))) \\
      &= A(0, A(0, A(0, A(0, A(1, 6))))) \displaybreak[0] \\
      &= A(0, A(0, A(0, A(0, A(0, A(1, 5)))))) \\
      &= A(0, A(0, A(0, A(0, A(0, A(0, A(1, 4))))))) \\
      &= A(0, A(0, A(0, A(0, A(0, A(0, A(0, A(1, 3)))))))) \\
      &= A(0, A(0, A(0, A(0, A(0, A(0, A(0, A(0, A(1, 2))))))))) \displaybreak[0] \\
      &= A(0, A(0, A(0, A(0, A(0, A(0, A(0, A(0, A(0, A(1, 1)))))))))) \\
      &= A(0, A(0, A(0, A(0, A(0, A(0, A(0, A(0, A(0, 2))))))))) \\
      &= A(0, A(0, A(0, A(0, A(0, A(0, A(0, A(0, 4)))))))) \\
      &= A(0, A(0, A(0, A(0, A(0, A(0, A(0, 8))))))) \displaybreak[0] \\
      &= A(0, A(0, A(0, A(0, A(0, A(0, 16)))))) \\
      &= A(0, A(0, A(0, A(0, A(0, 32))))) \\
      &= A(0, A(0, A(0, A(0, 64)))) \\
      &= A(0, A(0, A(0, 128))) \displaybreak[0] \\
      &= A(0, A(0, 256)) \\
      &= A(0, 512) \\
      &= \nombre{1024}.
    \end{align*}
  \end{sol}
\end{exercice}

\begin{exercice}
  Soit un algorithme avec $n$ entrées qui génère des valeurs pour
  chacun des carrés unitaires à la surface d'un cube
  $n \times n \times n$. Calculer la performance de l'algorithme. %
  \citep[Tiré de][]{Stephens:algorithms:2013}
  \begin{sol}
    Si le cube a $n$ unités de côté, sa surface est $6 n^2$. La
    performance de l'algorithme est donc $O(n^2)$. Nous pouvons
    obtenir cette réponse simplement en remarquant que l'algorithme
    doit calculer une aire.
  \end{sol}
\end{exercice}

\begin{exercice}
  Soit un algorithme avec $n$ entrées qui génère des valeurs pour
  chacun des cubes unitaires sur le pourtour d'un cube
  $n \times n \times n$, tel qu'illustré à la
  \autoref{fig:algorithmes:cubes-pourtour}. Calculer la performance de
  l'algorithme. %
  \citep[Tiré de][]{Stephens:algorithms:2013}
  \begin{sol}
    Le cube $n \times n \times n$ compte $8$ coins et $12$ côtés. La
    longueur de chaque côté en excluant les extrémités est $n - 2$. Le
    nombre total de cubes sur le pourtour est donc
    $12(n - 2) + 8 = 12n - 16$. La performance de l'algorithme est
    donc $O(n)$. Nous pouvons aussi obtenir cette réponse en
    remarquant que l'algorithme doit en définitive calculer une
    longueur.
  \end{sol}

  \begin{figure}
    \centering
    \setlength{\unitlength}{5mm}
    \thicklines
    \savebox{\cube}(1,1.5)[bl]{%
  \color{white}
  \polygon*(0,0)(0,1)(1,1)(1,0)
  \color{lightgray}
  \polygon*(0,1)(0.5,1.3)(1.5,1.3)(1,1)
  \color{gray}
  \polygon*(1,1)(1.5,1.3)(1.5,0.3)(1,0)
  \color{black}
  \polyline(0,0)(0,1)(1,1)(1,0)(0,0)
  \polyline(0,1)(0.5,1.3)(1.5,1.3)(1,1)(0,1)
  \polyline(1,1)(1.5,1.3)(1.5,0.3)(1,0)(1,1)
}

    \begin{picture}(11,3.9)(-4.5,0)
      \multiput(1.0,0.6)(1,0){3}{\usebox{\cube}}
      \multiput(0.5,0.3)(2,0){2}{\usebox{\cube}}
      \multiput(0.0,0.0)(1,0){3}{\usebox{\cube}}

      \multiput(1.0,1.6)(2,0){2}{\usebox{\cube}}
      \multiput(0.0,1.0)(2,0){2}{\usebox{\cube}}

      \multiput(1.0,2.6)(1,0){3}{\usebox{\cube}}
      \multiput(0.5,2.3)(2,0){2}{\usebox{\cube}}
      \multiput(0.0,2.0)(1,0){3}{\usebox{\cube}}

      \put(-4.5,1.5){\makebox(0,0)[l]{$n = 3$}}
    \end{picture} \\
    \bigskip
    \begin{picture}(11,5.2)(-4,0)
      \multiput(1.5,0.9)(1,0){4}{\usebox{\cube}}
      \multiput(1.0,0.6)(3,0){2}{\usebox{\cube}}
      \multiput(0.5,0.3)(3,0){2}{\usebox{\cube}}
      \multiput(0.0,0.0)(1,0){4}{\usebox{\cube}}

      \multiput(1.5,1.9)(3,0){2}{\usebox{\cube}}
      \multiput(0.0,1.0)(3,0){2}{\usebox{\cube}}

      \multiput(1.5,2.9)(3,0){2}{\usebox{\cube}}
      \multiput(0.0,2.0)(3,0){2}{\usebox{\cube}}

      \multiput(1.5,3.9)(1,0){4}{\usebox{\cube}}
      \multiput(1.0,3.6)(3,0){2}{\usebox{\cube}}
      \multiput(0.5,3.3)(3,0){2}{\usebox{\cube}}
      \multiput(0.0,3.0)(1,0){4}{\usebox{\cube}}

      \put(-4,2){\makebox(0,0)[l]{$n = 4$}}
    \end{picture} \\
    \bigskip
    \begin{picture}(11,6.5)(-3.5,0)
      \multiput(2.0,1.2)(1,0){5}{\usebox{\cube}}
      \multiput(1.5,0.9)(4,0){2}{\usebox{\cube}}
      \multiput(1.0,0.6)(4,0){2}{\usebox{\cube}}
      \multiput(0.5,0.3)(4,0){2}{\usebox{\cube}}
      \multiput(0.0,0.0)(1,0){5}{\usebox{\cube}}

      \multiput(2.0,2.2)(4,0){2}{\usebox{\cube}}
      \multiput(0.0,1.0)(4,0){2}{\usebox{\cube}}

      \multiput(2.0,3.2)(4,0){2}{\usebox{\cube}}
      \multiput(0.0,2.0)(4,0){2}{\usebox{\cube}}

      \multiput(2.0,4.2)(4,0){2}{\usebox{\cube}}
      \multiput(0.0,3.0)(4,0){2}{\usebox{\cube}}

      \multiput(2.0,5.2)(1,0){5}{\usebox{\cube}}
      \multiput(1.5,4.9)(4,0){2}{\usebox{\cube}}
      \multiput(1.0,4.6)(4,0){2}{\usebox{\cube}}
      \multiput(0.5,4.3)(4,0){2}{\usebox{\cube}}
      \multiput(0.0,4.0)(1,0){5}{\usebox{\cube}}

      \put(-3.5,2.5){\makebox(0,0)[l]{$n = 5$}}
    \end{picture}
    \caption{Illustrations pour un algorithme qui génère des valeurs
      sur le pourtour d'un cube}
    \label{fig:algorithmes:cubes-pourtour}
  \end{figure}
\end{exercice}

\begin{exercice}
  Soit un algorithme avec $n$ entrées qui génère des valeurs pour
  chacun des carrés unitaires d'une pyramide comme dans le jeu
  classique \link{https://fr.wikipedia.org/wiki/Q*bert}{Q*bert}, tel
  qu'illustré à la \autoref{fig:algorithmes:q*bert}. Supposer que les
  cubes invisibles sont présents de telle sorte que les structures
  sont pleines. Calculer la performance de l'algorithme. %
  \citep[Tiré de][]{Stephens:algorithms:2013}
  \begin{sol}
    L'étage $k$ d'une pyramide de $n$ étage compte $k(k + 1)/2$ cubes.
    Le nombre total de cubes est donc égal à
    \begin{equation*}
      \sum_{k = 1}^n \frac{k (k + 1)}{2}.
    \end{equation*}
    À l'aide des identités
    \begin{align*}
      \sum_{k = 1}^n k &= \frac{n (n + 1)}{2} \\
      \sum_{k = 1}^n k^2 &= \frac{n (n + 1) (2n + 1)}{6},
    \end{align*}
    on obtient avec un peu d'algèbre que
    \begin{equation*}
      \sum_{k = 1}^n \frac{k (k + 1)}{2} = \frac{n^3 + 3 n^2 + 2n}{6}.
    \end{equation*}
    La performance de l'algorithme est donc $O(n^3)$, une réponse
    intuitivement claire dans la mesure où l'algorithme calcule un
    volume.
  \end{sol}

  \begin{figure}
    \centering
    \setlength{\unitlength}{8mm}
    \thicklines
    \savebox{\cube}(1,1.5)[bl]{%
  \color{white}
  \polygon*(0,0)(0,1)(1,1)(1,0)
  \color{lightgray}
  \polygon*(0,1)(0.5,1.3)(1.5,1.3)(1,1)
  \color{gray}
  \polygon*(1,1)(1.5,1.3)(1.5,0.3)(1,0)
  \color{black}
  \polyline(0,0)(0,1)(1,1)(1,0)(0,0)
  \polyline(0,1)(0.5,1.3)(1.5,1.3)(1,1)(0,1)
  \polyline(1,1)(1.5,1.3)(1.5,0.3)(1,0)(1,1)
}

    \begin{minipage}{0.48\linewidth}
      \begin{picture}(7,1.3)(-3,0)
        \put(0.0,0.0){\usebox{\cube}}

        \put(-3.0,0.5){\makebox(0,0)[l]{$n = 1$}}
      \end{picture}
    \end{minipage}
    \hfill
    \begin{minipage}{0.48\linewidth}
      \begin{picture}(7,2.6)(-2,0)
        \multiput(0.0,0.0)(1.5,0.3){2}{\usebox{\cube}}
        \multiput(0.5,1.3)(1.5,0.3){1}{\usebox{\cube}}

        \put(-2.0,1.3){\makebox(0,0)[l]{$n = 2$}}
      \end{picture}
    \end{minipage} \\
    \bigskip
    \begin{minipage}{0.48\linewidth}
      \begin{picture}(7,3.9)(-2,0)
        \multiput(0.0,0.0)(1.5,0.3){3}{\usebox{\cube}}
        \multiput(0.5,1.3)(1.5,0.3){2}{\usebox{\cube}}
        \multiput(1.0,2.6)(1.5,0.3){1}{\usebox{\cube}}

        \put(-2.0,1.8){\makebox(0,0)[l]{$n = 3$}}
      \end{picture}
    \end{minipage}
    \hfill
    \begin{minipage}{0.48\linewidth}
      \begin{picture}(7,5.2)(-1,0)
        \multiput(0.0,0.0)(1.5,0.3){4}{\usebox{\cube}}
        \multiput(0.5,1.3)(1.5,0.3){3}{\usebox{\cube}}
        \multiput(1.0,2.6)(1.5,0.3){2}{\usebox{\cube}}
        \multiput(1.5,3.9)(1.5,0.3){1}{\usebox{\cube}}

        \put(-1.0,2.6){\makebox(0,0)[l]{$n = 4$}}
      \end{picture}
    \end{minipage}
    \caption{Illustrations pour un algorithme qui génère des valeurs
      pour une pyramide comme dans le jeu Q*bert}
    \label{fig:algorithmes:q*bert}
  \end{figure}
\end{exercice}

\begin{exercice}
  Comparer à l'aide d'un graphique pour $n \in (0, 20)$ le rythme de
  croissance des fonctions $\log n$, $\sqrt{n}$, $n$, $n \log n$,
  $n^2$ et $n!$.
  \begin{sol}
    Vous pouvez tracer un graphique regroupant toutes les fonctions
    dans R avec les expressions suivantes.
    \begin{Schunk}
\begin{Sinput}
> plot(0:20, 0:20, type = "n")
> curve(log(x), add = TRUE, col = 1)
> curve(sqrt(x), add = TRUE, col = 2)
> curve(I(x), add = TRUE, col = 3)
> curve(x * log(x), add = TRUE, col = 4)
> curve(x^2, add = TRUE, col = 5)
> curve(factorial(x), add = TRUE, col = 6)
\end{Sinput}
    \end{Schunk}
  \end{sol}
\end{exercice}


\begin{Filesave}{solutions}
\endgroup                               % fin de \lstset local
\end{Filesave}

\Closesolutionfile{solutions}

%%% Local Variables:
%%% mode: latex
%%% TeX-engine: xetex
%%% TeX-master: "programmer-avec-r"
%%% coding: utf-8
%%% End:

\section{Pr�sentation du langage S}


\subsection{Le langage S}

\begin{frame}
  \frametitle{Un langage pour �programmer avec des donn�es�}

  \begin{itemize}
  \item D�velopp� chez Bell Laboratories
  \item Langage de programmation complet et autonome
  \item Inspir� de plusieurs langages, dont l'APL et le Lisp:
    \begin{itemize}
    \item interpr�t� (et non compil�)
    \item sans d�claration obligatoire des variables
    \item bas� sur la notion de vecteur
    \item particuli�rement puissant pour les applications
      math�matiques et statistiques (et donc actuarielles)
    \end{itemize}
  \end{itemize}
\end{frame}


\subsection{Les moteurs S}

\begin{frame}
  \frametitle{Deux �moteurs� ou dialectes du langage S}

  \begin{itemize}
  \item Jusqu'� r�cemment, le plus connu �tait S-PLUS de Insightful
    Corporation, maintenant S$+$ de Timbco Software
  \item Le leader est maintenant \textsf{R}, ou GNU S, est une version
    libre (\emph{Open Source}) �\emph{not unlike S}�
  \item Environnements int�gr�s de manipulation de donn�es, de calcul
    et de pr�paration de graphiques
  \end{itemize}
\end{frame}


\subsection[Interfaces]{Interfaces pour S-Plus et R}

\begin{frame}
  \frametitle{D'abord des applications en ligne de commande}

  \begin{itemize}
  \item S$+$ poss�de une interface graphique �labor�e qui permet
    d'utiliser le logiciel sans trop conna�tre le langage de
    programmation
  \item \textsf{R} dispose �galement d'une interface graphique sous
    Windows et Mac OS
  \item L'�dition s�rieuse de code S b�n�ficie cependant grandement
    d'un bon �diteur de texte
  \end{itemize}
\end{frame}

\begin{frame}
  \begin{itemize}
  \item Question 6.2 de la foire aux questions (FAQ) de
    \textsf{R}: �Devrais-je utiliser \textsf{R} � l'int�rieur de
    Emacs?�
  \item R�ponse: �Oui, tout � fait�
  \item Nous apprendrons � utiliser \textsf{R} � l'int�rieur de GNU
    Emacs avec le mode ESS
  \item Plusieurs autres options disponibles: choisissez celle que
    vous pr�f�rez
  \end{itemize}
\end{frame}


\subsection[Emacs et ESS]{Installation de Emacs avec ESS}

\begin{frame}
  \frametitle{Installation de Emacs avec ESS}

  \begin{itemize}
  \item Pour une installation simplifi�e de GNU Emacs et ESS,
    consulter le site Internet
    \begin{center}
      \url{http://vgoulet.act.ulaval.ca/emacs/}
    \end{center}
  \item Voir l'annexe A du document d'accompagnement pour les plus
    importantes commandes
  \end{itemize}
\end{frame}


\subsection[D�marrer et quitter]{D�marrer et quitter R}

\begin{frame}
  \frametitle{D�marrer et quitter R}

  \begin{itemize}
  \item Pour d�marrer \textsf{R} \R � l'int�rieur de Emacs:

    \texttt{M-x R RET}

    puis sp�cifier un dossier de travail
  \item Une console \textsf{R} est ouverte dans un \emph{buffer} nomm�
    \texttt{*R*}
  \item Pour quitter, deux options:
    \begin{enumerate}
    \item taper \fonction{q()} � la ligne de commande
    \item dans Emacs, faire \ess{C-c C-q}
    \end{enumerate}
  \end{itemize}
\end{frame}


\subsection{Strat�gies de travail}

\begin{frame}
  \frametitle{Deux grandes fa�ons de travailler avec \textsf{R}}

  \begin{enumerate}
  \item Le code est virtuel et les objets sont r�els
  \item Le code est r�el et les objets sont virtuels
  \end{enumerate}
\end{frame}

\begin{frame}
  \frametitle{Code virtuel, objets r�els}

  \begin{itemize}
  \item C'est l'approche qu'encouragent les interfaces graphiques,
    mais aussi la moins pratique � long terme
  \item On entre des expressions directement � la ligne de commande
    pour les �valuer imm�diatement
  \item Les objets cr��s au cours d'une session de travail sont
    sauvegard�s
  \item Par contre, le code utilis� pour cr�er ces objets est perdu
    lorsque l'on quitte \textsf{R}, � moins de sauvegarder celui-ci
    dans des fichiers
  \end{itemize}
\end{frame}

\begin{frame}
  \frametitle{Code r�el, objets virtuels}

  \begin{itemize}
  \item C'est l'approche que nous favoriserons
  \item Le travail se fait essentiellement dans des fichiers de script
    (de simples fichiers de texte) dans lesquels sont sauvegard�es les
    expressions (parfois complexes!) et le code des fonctions
    personnelles
  \item Les objets sont cr��s au besoin en ex�cutant le code
  \end{itemize}
\end{frame}

\begin{frame}
  \frametitle{Emacs permet un travail efficace}
  \begin{itemize}
  \item D�marrer un processus \textsf{R} et sp�cifier le dossier de
    travail
  \item Ouvrir un fichier de script avec \ess{C-x C-f} (pour cr�er un
    nouveau fichier, ouvrir un fichier qui n'existe pas)
  \item Positionner le curseur sur une expression et faire \ess{C-c
      C-n} pour l'�valuer
  \item Le r�sultat appara�t dans le \emph{buffer} \texttt{*R*}
  \end{itemize}
\end{frame}


\subsection[Gestion des projets]{Gestion des projets ou environnements de
  travail}

\begin{frame}
  \frametitle{Gestion des projets ou environnements de travail}

  \begin{itemize}[<+->]
  \item \textsf{R} travaille dans un dossier et non avec des fichiers
    individuels
  \item Dans \textsf{R}, \R les objets cr��s sont conserv�s en m�moire
    \begin{itemize}
    \item sauvegard�s en quittant l'application, ou
    \item avec la commande \fonction{save.image()}
    \end{itemize}
  \item L'environnement de travail (\emph{workspace}) est sauvegard�
    dans le fichier \texttt{.RData} dans le dossier de travail
  \end{itemize}
\end{frame}

\begin{frame}
  \frametitle{Comment d�terminer le dossier de travail}
  \begin{itemize}
  \item Avec Emacs et ESS on doit sp�cifier le dossier de travail �
    chaque fois que l'on d�marre un processus R
  \item Les interfaces graphiques permettent �galement de sp�cifier le
    dossier de travail au lancement de l'application
  \end{itemize}
\end{frame}


\subsection{Consulter l'aide en ligne}

\begin{frame}[fragile]
  \frametitle{La premi�re source d'aide}

  \begin{itemize}
  \item Rubriques d'aide contiennent une foule d'informations ainsi
    que des exemples d'utilisation
  \item Leur consultation est tout � fait essentielle
  \item Pour consulter la rubrique d'aide de la fonction \code{foo},
    on peut entrer � la ligne de commande
\begin{Sinput}
> ?foo
\end{Sinput}
  \item Dans Emacs, \code{C-c C-v foo RET} ouvrira
    la rubrique d'aide de la fonction \code{foo} dans un nouveau
    \emph{buffer}
  \end{itemize}
\end{frame}

\subsection[Documentation]{O� trouver de la documentation}

\begin{frame}
  \frametitle{Plusieurs ressources disponibles}

  \begin{itemize}
  \item S$+$ est livr� avec quatre livres, mais aucun ne s'av�re
    vraiment utile pour apprendre le langage S
  \item Plusieurs livres --- en versions papier ou �lectronique,
    gratuits ou non --- ont �t� publi�s sur S$+$ et/ou \textsf{R}
  \item Liste exhaustive dans le site du projet \textsf{R}
  \end{itemize}
\end{frame}


%%% Local Variables:
%%% mode: latex
%%% TeX-master: "introduction_programmation_S_diapos"
%%% End:

\chapter{Bases du langage S}
\label{bases}


Ce chapitre pr�sente les bases du langage S, soit les notions
d'expression et d'affectation, la description d'un objet S et les
mani�res de cr�er les objets les plus usuels lorsque le S est utilis�
comme langage de programmation.

\section{Commandes S}
\label{bases:commandes}

Toute commande S est soit une \emph{affectation}\index{affectation},
soit une \emph{expression}\index{expression}.
\begin{itemize}
\item Normalement, une expression est imm�diatement �valu�e et le
  r�sultat est affich� � l'�cran:
\begin{Schunk}
\begin{Sinput}
> 2 + 3
\end{Sinput}
\begin{Soutput}
[1] 5
\end{Soutput}
\begin{Sinput}
> pi
\end{Sinput}
\begin{Soutput}
[1] 3.141593
\end{Soutput}
\begin{Sinput}
> cos(pi/4)
\end{Sinput}
\begin{Soutput}
[1] 0.7071068
\end{Soutput}
\end{Schunk}
\item Lors d'une affectation, une expression est �valu�e, mais le
  r�sultat est stock� dans un objet (variable) et rien n'est affich� �
  l'�cran. Le symbole d'affectation est \Fonction{<-} (ou
  \Fonction{->}).
\begin{Schunk}
\begin{Sinput}
> a <- 5
> a
\end{Sinput}
\begin{Soutput}
[1] 5
\end{Soutput}
\begin{Sinput}
> b <- a
> b
\end{Sinput}
\begin{Soutput}
[1] 5
\end{Soutput}
\end{Schunk}
\item �viter d'utiliser l'op�rateur \fonction{=} pour affecter une
  valeur � une variable, puisqu'il ne fonctionne que dans certaines
  situations seulement.
\item Dans S-Plus (mais plus dans \textsf{R} depuis la version 1.8.0),
  on peut �galement affecter avec le caract�re �\fonction{\_}�, mais cet
  emploi est fortement d�courag� puisqu'il rend le code difficile �
  lire. Dans le mode ESS de Emacs, taper ce caract�re g�n�re carr�ment
  \verb*| <- |.
\end{itemize}

\begin{astuce}
  Il arrive fr�quemment que l'on souhaite affecter le r�sultat d'un
  calcul dans un objet et en m�me temps voir ce r�sultat. Pour ce
  faire, placer l'affectation entre parenth�ses (l'op�ration
  d'affectation devient alors une nouvelle expression):
\begin{Schunk}
\begin{Sinput}
> (a <- 2 + 3)
\end{Sinput}
\begin{Soutput}
[1] 5
\end{Soutput}
\end{Schunk}
\end{astuce}


\section{Conventions pour les noms d'objets}
\index{noms d'objets!conventions}
\label{bases:noms}

Les caract�res permis pour les noms d'objets sont les lettres a--z,
A--Z, les chiffres 0--9 et le point �.�. Le caract�re \R �\code{\_}�
est maintenant permis dans \textsf{R}, mais son utilisation est
d�courag�e.

\begin{itemize}
\item Les noms d'objets ne peuvent commencer par un chiffre.
\item Le S est sensible � la casse, ce qui signifie que \code{foo},
  \code{Foo} et \code{FOO} sont trois objets distincts. Un moyen
  simple d'�viter des erreurs li�es � la casse consiste � n'employer
  que des lettres minuscules.
\item Certains noms sont utilis�s par le syst�me, aussi vaut-il mieux
  �viter de les utiliser. En particulier, �viter d'utiliser
  \begin{center}
    \code{c}, \code{q}, \code{t}, \code{C}, \code{D},
    \code{I}, \code{diff}, \code{length}, \code{mean},
    \code{pi}, \code{range}, \code{var}.
  \end{center}
\item \index{noms d'objets!r�serv�s} Certains mots sont r�serv�s pour
  le syst�me et il est interdit de les utiliser comme nom d'objet. Les
  mots r�serv�s sont:
  \begin{center}
    \code{Inf}, \code{NA}, \code{NaN},
    \code{NULL} \\
    \code{break}, \code{else},
    \code{for}, \code{function}, \code{if}, \code{in},
    \code{next}, \code{repeat}, \code{return}, \code{while}.
  \end{center}
\item Dans S-Plus 6.1 et plus, \Splus
  \code{T}\index{T@\code{T}|see{\code{TRUE}}} et \objet{TRUE} (vrai),
  ainsi que \code{F}\index{F@\code{F}|see{\code{FALSE}}} et
  \objet{FALSE} (faux) sont �galement des noms r�serv�s.
\item Dans \textsf{R}, \R les noms \code{TRUE} et \code{FALSE}
  sont �galement r�serv�s. Les variables \code{T} et \code{F}
  prennent par d�faut les valeurs \code{TRUE} et \code{FALSE},
  respectivement, mais peuvent �tre r�affect�es.
\begin{Schunk}
\begin{Sinput}
> T
\end{Sinput}
\begin{Soutput}
[1] TRUE
\end{Soutput}
\end{Schunk}
\begin{Schunk}
\begin{Sinput}
> TRUE <- 3
\end{Sinput}
\end{Schunk}
\begin{Schunk}
\begin{Soutput}
Erreur dans TRUE <- 3 : membre gauche de
l'assignation (do_set) incorrect
\end{Soutput}
\end{Schunk}
\begin{Schunk}
\begin{Sinput}
> (T <- 3)
\end{Sinput}
\begin{Soutput}
[1] 3
\end{Soutput}
\end{Schunk}
\end{itemize}


\section{Les objets S}
\label{bases:objets}

Tout dans le langage S est un objet, m�me les fonctions et les
op�rateurs. Les objets poss�dent au minimum un \emph{mode} et une
\emph{longueur}.

\begin{itemize}
\item Le mode d'un objet est obtenu avec la fonction \Fonction{mode}.
\begin{Schunk}
\begin{Sinput}
> v <- c(1, 2, 5, 9)
> mode(v)
\end{Sinput}
\begin{Soutput}
[1] "numeric"
\end{Soutput}
\end{Schunk}
\item La longueur d'un objet est obtenue avec la fonction
  \Fonction{length}.
\begin{Schunk}
\begin{Sinput}
> length(v)
\end{Sinput}
\begin{Soutput}
[1] 4
\end{Soutput}
\end{Schunk}
\item Certains objets sont �galement dot�s d'un ou plusieurs
  \emph{attributs}.
\end{itemize}


\subsection{Modes et types de donn�es}

Le mode\Index{mode} prescrit ce qu'un objet peut contenir. � ce titre,
un objet ne peut avoir qu'un seul mode. Le tableau \ref{tab:modes}
contient la liste des modes disponibles en S. � chacun de ces modes
correspond une fonction du m�me nom servant � cr�er un objet de ce
mode.

\begin{table}
  \centering
  \begin{tabular}{ll}
    \toprule
    \Mode{numeric}   & nombres r�els \\
    \Mode{complex}   & nombres complexes \\
    \Mode{logical}   & valeurs bool�ennes (vrai/faux) \\
    \Mode{character} & cha�nes de caract�res \\
    \Mode{function}  & fonction \\
    \Mode{list}      & donn�es quelconques \\
    \bottomrule
  \end{tabular}
  \caption{Modes disponibles et contenus correspondants}
  \label{tab:modes}
\end{table}

\subsection{Longueur}

La longueur\Index{longueur} d'un objet est �gale au nombre d'�l�ments
qu'il contient.

\begin{itemize}
\item La longueur d'une cha�ne de caract�res est toujours 1. Un
  objet de mode \code{character} doit contenir plusieurs cha�nes
  de caract�res pour que sa longueur soit sup�rieure � 1.
\begin{Schunk}
\begin{Sinput}
> v <- "actuariat"
> length(v)
\end{Sinput}
\begin{Soutput}
[1] 1
\end{Soutput}
\begin{Sinput}
> v <- c("a", "c", "t", "u", "a", "r", "i", 
+     "a", "t")
> length(v)
\end{Sinput}
\begin{Soutput}
[1] 9
\end{Soutput}
\end{Schunk}
\item Un objet peut �tre de longueur 0 et doit alors �tre interpr�t�
  comme un contenant vide\index{vide|see{\code{NULL}}}.
\begin{Schunk}
\begin{Sinput}
> v <- numeric(0)
> length(v)
\end{Sinput}
\begin{Soutput}
[1] 0
\end{Soutput}
\end{Schunk}
\end{itemize}


\subsection{Attributs}

Les attributs\Index{attribut} d'un objet sont des �l�ments
d'information additionnels li�s � cet objet. La liste des attributs
les plus fr�quemment rencontr�s se trouve au tableau
\ref{tab:attributs}. Pour chaque attribut, il existe une fonction du
m�me nom servant � extraire l'attribut correspondant d'un objet.

\begin{table}
  \centering
  \begin{tabular}{ll}
    \toprule
    \Attribut{class}    &
    affecte le comportement d'un objet \\
    \Attribut{dim}      &
    dimensions\index{dimension} des matrices et tableaux \\
    \Attribut{dimnames} &
    �tiquettes\index{etiquette@�tiquette} des dimensions des matrices
    et tableaux \\
    \Attribut{names}    &
    �tiquettes des �l�ments d'un objet \\
    \bottomrule
  \end{tabular}
  \caption{Attributs les plus usuels d'un objet et leur effet}
  \label{tab:attributs}
\end{table}


\subsection{L'objet sp�cial \code{NA}}

\Objet{NA} est fr�quemment utilis� pour repr�senter les donn�es
manquantes.

\begin{itemize}
\item Son mode est \mode{logical}.
\item Toute op�ration impliquant une donn�e \code{NA} a comme
  r�sultat \code{NA}.
\item Certaines fonctions (\fonction{sum}, \fonction{mean}, par
  exemple), ont par cons�quent un argument \Argument{na.rm} qui,
  lorsque \code{TRUE}, �limine les donn�es manquantes avant de
  faire un calcul.
\item La fonction \Fonction{is.na} permet de tester si les �l�ments
  d'un objet sont \code{NA} ou non.
\end{itemize}

\subsection{L'objet sp�cial \code{NULL}}

\Objet{NULL} repr�sente �rien�, ou le vide.

\begin{itemize}
\item Son mode est \Mode{NULL}.
\item Sa longueur est 0.
\item Diff�rent d'un objet vide:
  \begin{itemize}
  \item un objet de longueur 0 est un contenant vide;
  \item \code{NULL} est �pas de contenant�.
  \end{itemize}
\item La fonction \Fonction{is.null} teste si un objet est \code{NULL}
  ou non.
\end{itemize}


\section{Vecteurs}
\label{bases:vecteurs}

En S, � peu de choses pr�s, \emph{tout} est un vecteur\index{vecteur}.
(Il n'y a pas de notion de scalaire.)

\begin{itemize}
\item Dans un vecteur simple, tous les �l�ments doivent �tre du
  m�me mode.
\item Il est possible (et souvent souhaitable) de donner une �tiquette
  � chacun des �l�ments d'un vecteur.
\begin{Schunk}
\begin{Sinput}
> (v <- c(a = 1, b = 2, c = 5))
\end{Sinput}
\begin{Soutput}
a b c 
1 2 5 
\end{Soutput}
\begin{Sinput}
> v <- c(1, 2, 5)
> names(v) <- c("a", "b", "c")
> v
\end{Sinput}
\begin{Soutput}
a b c 
1 2 5 
\end{Soutput}
\end{Schunk}
\item Les fonctions de base pour cr�er des vecteurs sont
  \begin{itemize}
  \item \Fonction{c} (concat�nation)
  \item \Fonction{numeric} (vecteur de mode \mode{numeric})
  \item \Fonction{logical} (vecteur de mode \mode{logical})
  \item \Fonction{character} (vecteur de mode \mode{character}).
  \end{itemize}
\item L'indi�age dans un vecteur se fait avec \fonction{[\ ]}. On
  peut extraire un �l�ment d'un vecteur par sa position ou par son
  �tiquette, si elle existe (auquel cas cette approche est beaucoup
  plus s�re).
\begin{Schunk}
\begin{Sinput}
> v[3]
\end{Sinput}
\begin{Soutput}
c 
5 
\end{Soutput}
\begin{Sinput}
> v["c"]
\end{Sinput}
\begin{Soutput}
c 
5 
\end{Soutput}
\end{Schunk}
  La section \ref{bases:indicage} traite plus en d�tail de l'indi�age de
  vecteurs et matrices.
\end{itemize}


\section{Matrices et tableaux}
\label{bases:matrices}

Une matrice\index{matrice} ou, de fa�on plus g�n�rale, un
tableau\index{tableau} (\emph{array}) n'est rien d'autre qu'un vecteur
dot� d'un attribut \attribut{dim}. � l'interne, une matrice est donc
stock�e sous forme de vecteur.

\begin{itemize}
\item La fonction de base pour cr�er des matrices est
  \Fonction{matrix}.
\item La fonction de base pour cr�er des tableaux est
  \Fonction{array}.
\item \emph{Important}: \warning les matrices et tableaux sont remplis
  en faisant d'abord varier la premi�re dimension, puis la seconde,
  etc. Pour les matrices, cela revient � remplir par colonne.
\begin{Schunk}
\begin{Sinput}
> matrix(1:6, nrow = 2, ncol = 3)
\end{Sinput}
\begin{Soutput}
     [,1] [,2] [,3]
[1,]    1    3    5
[2,]    2    4    6
\end{Soutput}
\begin{Sinput}
> matrix(1:6, nrow = 2, ncol = 3, byrow = TRUE)
\end{Sinput}
\begin{Soutput}
     [,1] [,2] [,3]
[1,]    1    2    3
[2,]    4    5    6
\end{Soutput}
\end{Schunk}
\item On extrait les �l�ments d'une matrice en pr�cisant leurs
  positions sous la forme (ligne, colonne) dans la matrice, ou encore
  leurs positions dans le vecteur sous-jacent.
\begin{Schunk}
\begin{Sinput}
> (m <- matrix(c(40, 80, 45, 21, 55, 32), 
+     nrow = 2, ncol = 3))
\end{Sinput}
\begin{Soutput}
     [,1] [,2] [,3]
[1,]   40   45   55
[2,]   80   21   32
\end{Soutput}
\begin{Sinput}
> m[1, 2]
\end{Sinput}
\begin{Soutput}
[1] 45
\end{Soutput}
\begin{Sinput}
> m[3]
\end{Sinput}
\begin{Soutput}
[1] 45
\end{Soutput}
\end{Schunk}
\item La fonction \Fonction{rbind} permet de fusionner verticalement
  deux matrices (ou plus) ayant le m�me nombre de colonnes.
\begin{Schunk}
\begin{Sinput}
> n <- matrix(1:9, nrow = 3)
> rbind(m, n)
\end{Sinput}
\begin{Soutput}
     [,1] [,2] [,3]
[1,]   40   45   55
[2,]   80   21   32
[3,]    1    4    7
[4,]    2    5    8
[5,]    3    6    9
\end{Soutput}
\end{Schunk}
\item La fonction \Fonction{cbind} permet de fusionner horizontalement
  deux matrices (ou plus) ayant le m�me nombre de lignes.
\begin{Schunk}
\begin{Sinput}
> n <- matrix(1:4, nrow = 2)
> cbind(m, n)
\end{Sinput}
\begin{Soutput}
     [,1] [,2] [,3] [,4] [,5]
[1,]   40   45   55    1    3
[2,]   80   21   32    2    4
\end{Soutput}
\end{Schunk}
\end{itemize}


\section{Listes}
\label{bases:listes}

Une liste\index{liste} est un type de vecteur sp�cial dont les
�l�ments peuvent �tre de n'importe quel mode, y compris le mode
\mode{list} (ce qui permet d'embo�ter des listes).

\begin{itemize}
\item La fonction de base pour cr�er des listes est \Fonction{list}.
\item Il est g�n�ralement pr�f�rable de nommer les �l�ments d'une
  liste. Il est en effet plus simple et s�r d'extraire les �l�ments
  par leur �tiquette.
\item L'extraction\Index{indi�age!liste} des �l�ments d'une liste peut
  se faire de deux fa�ons:
  \begin{enumerate}
  \item avec des doubles crochets \fonction{[[\ ]]};
  \item par leur �tiquette avec \code{nom.liste\$etiquette.element}.
  \end{enumerate}
\item La fonction \Fonction{unlist} convertit une liste en un vecteur
  simple. Attention, cette fonction peut �tre destructrice si la
  structure de la liste est importante.
\end{itemize}


\section{\emph{Data frames}}
\label{bases:dataframes}

Les vecteurs, matrices, tableaux (\emph{arrays}) et listes sont les
types d'objets les plus fr�quemment utilis�s en S pour la
programmation de fonctions personnelles ou la simulation. L'analyse de
donn�es --- la r�gression lin�aire, par exemple --- repose toutefois
davantage sur les \emph{data frames}\Index{data frame}.

\begin{itemize}
\item Un \emph{data frame} est une liste de classe \classe{data.frame}
  dont tous les �l�ments sont de la m�me longueur.
\item G�n�ralement repr�sent� sous forme d'un tableau � deux
  dimensions (visuellement similaire � une matrice). Chaque �l�ment de
  la liste sous-jacente correspond � une colonne.
\item On peut donc obtenir les �tiquettes des colonnes avec la fonction
  \fonction{names} (ou \fonction{colnames} \R dans \textsf{R}).  Les
  �tiquettes des lignes sont quant � elles obtenues avec
  \fonction{row.names} (ou \fonction{rownames} dans \textsf{R}).
\item Plus g�n�ral qu'une matrice puisque les colonnes peuvent �tre de
  modes diff�rents (\mode{numeric}, \mode{complex}, \mode{character}
  ou \mode{logical}).
\item Peut �tre indic� � la fois comme une liste et comme une matrice.
\item Cr�� avec la fonction \Fonction{data.frame} ou
  \Fonction{as.data.frame} (pour convertir une matrice en \emph{data
    frame}, par exemple).
\item Les fonctions \fonction{rbind} et \fonction{cbind} peuvent �tre
  utilis�es pour ajouter des lignes ou des colonnes, respectivement.
\item On peut rendre les colonnes d'un \emph{data frame} (ou d'une
  liste) visibles dans l'espace de travail avec la fonction
  \Fonction{attach}, puis les masquer avec \Fonction{detach}.
\end{itemize}

Ce type d'objet est moins important lors de l'apprentissage du langage
de programmation.




\section{Indi�age}
\label{bases:indicage}

L'indi�age des vecteurs\Index{indi�age!vecteur} et
matrices\Index{indi�age!matrice} a d�j� �t� bri�vement pr�sent� aux
sections \ref{bases:vecteurs} et \ref{bases:matrices}. Cette section
contient plus de d�tails sur cette proc�dure des plus communes lors de
l'utilisation du langage S. On se concentre toutefois sur le
traitement des vecteurs. Se r�f�rer �galement � \citet[section
2.3]{MASS} pour de plus amples d�tails.

Il existe quatre fa�ons d'indicer un vecteur dans le langage S. Dans
tous les cas, l'indi�age se fait � l'int�rieur de crochets \Fonction{[\ ]}.
\begin{enumerate}
\item Avec un vecteur d'entiers positifs. Les �l�ments se trouvant aux
  positions correspondant aux entiers sont extraits du vecteur, dans
  l'ordre. C'est la technique la plus courante.
\begin{Schunk}
\begin{Sinput}
> letters[c(1:3, 22, 5)]
\end{Sinput}
\begin{Soutput}
[1] "a" "b" "c" "v" "e"
\end{Soutput}
\end{Schunk}
\item Avec un vecteur d'entiers n�gatifs. Les �l�ments se trouvant aux
  positions correspondant aux entiers n�gatifs sont alors
  \emph{�limin�s} du vecteur.
\begin{Schunk}
\begin{Sinput}
> letters[c(-(1:3), -5, -22)]
\end{Sinput}
\begin{Soutput}
 [1] "d" "f" "g" "h" "i" "j" "k" "l" "m" "n" "o" "p"
[13] "q" "r" "s" "t" "u" "w" "x" "y" "z"
\end{Soutput}
\end{Schunk}
\item Avec un vecteur bool�en. Le vecteur d'indi�age doit alors �tre
  de la m�me longueur que le vecteur indic�. Les �l�ments
  correspondant � une valeur \code{TRUE} sont extraits du vecteur,
  alors que ceux correspondant � \code{FALSE} sont �limin�s.
\begin{Schunk}
\begin{Sinput}
> letters > "f" & letters < "q"
\end{Sinput}
\begin{Soutput}
 [1] FALSE FALSE FALSE FALSE FALSE FALSE  TRUE  TRUE
 [9]  TRUE  TRUE  TRUE  TRUE  TRUE  TRUE  TRUE  TRUE
[17] FALSE FALSE FALSE FALSE FALSE FALSE FALSE FALSE
[25] FALSE FALSE
\end{Soutput}
\begin{Sinput}
> letters[letters > "f" & letters < "q"]
\end{Sinput}
\begin{Soutput}
 [1] "g" "h" "i" "j" "k" "l" "m" "n" "o" "p"
\end{Soutput}
\end{Schunk}
\item Avec une cha�ne de caract�res. Utile pour extraire les �l�ments
  d'un vecteur � condition que ceux-ci soient nomm�s.
\begin{Schunk}
\begin{Sinput}
> x <- c(Rouge = 2, Bleu = 4, Vert = 9, Jaune = -5)
> x[c("Bleu", "Jaune")]
\end{Sinput}
\begin{Soutput}
 Bleu Jaune 
    4    -5 
\end{Soutput}
\end{Schunk}
\end{enumerate}


\section{Exemples}
\label{bases:exemples}

\lstinputlisting{bases.R}


\section{Exercices}
\label{bases:exercices}

\Opensolutionfile{reponses}[reponses-bases]
\Writetofile{reponses}{\protect\section*{Chapitre \protect\ref{bases}}}

\begin{exercice}
  \begin{enumerate}
  \item �crire une expression S pour cr�er la liste suivante:
\begin{Schunk}
\begin{Soutput}
[[1]]
[1] 1 2 3 4 5

$data
     [,1] [,2] [,3]
[1,]    1    3    5
[2,]    2    4    6

[[3]]
[1] 0 0 0

$test
[1] FALSE FALSE FALSE FALSE
\end{Soutput}
\end{Schunk}
  \item \index{etiquette@�tiquette} Extraire les �tiquettes de la
    liste.
  \item \index{mode} \index{longueur} Trouver le mode et la longueur
    du quatri�me �l�ment de la liste.
  \item \index{dimension} Extraire les dimensions du second �l�ment de
    la liste.
  \item \index{indi�age!liste} Extraire les deuxi�me et troisi�me
    �l�ments du second �l�ment de la liste.
  \item Remplacer le troisi�me �l�ment de la liste par le vecteur
    \verb|3:8|.
  \end{enumerate}
  \begin{rep}
    Soit \code{x} le nom de la liste.
    \begin{enumerate}
\item
\begin{Schunk}
\begin{Sinput}
> list(1:5, data = matrix(1:6, 2, 3), numeric(3), 
+     test = logical(4))
\end{Sinput}
\end{Schunk}
\item
\begin{Schunk}
\begin{Sinput}
> names(x)
\end{Sinput}
\end{Schunk}
\item
\begin{Schunk}
\begin{Sinput}
> mode(x$test)
> length(x$test)
\end{Sinput}
\end{Schunk}
\item
\begin{Schunk}
\begin{Sinput}
> dim(x$data)
\end{Sinput}
\end{Schunk}
\item
\begin{Schunk}
\begin{Sinput}
> x[[2]][c(2, 3)]
\end{Sinput}
\end{Schunk}
\item
\begin{Schunk}
\begin{Sinput}
> x[[3]] <- 3:8
\end{Sinput}
\end{Schunk}
    \end{enumerate}
  \end{rep}
\end{exercice}

\begin{exercice}
  \index{indi�age!vecteur}
  Soit \code{obs} un vecteur contenant les valeurs suivantes:
  \begin{center}
\begin{Schunk}
\begin{Sinput}
> obs
\end{Sinput}
\begin{Soutput}
 [1]  3 19 13  3 13 11  9  8 15  5 20 15  9 11 18  4
[17]  1 12  9 12
\end{Soutput}
\end{Schunk}
  \end{center}
  �crire une expression S permettant d'extraire les �l�ments suivants.
  \begin{enumerate}
  \item Le deuxi�me �l�ment de l'�chantillon.
  \item Les cinq premiers �l�ments de l'�chantillon.
  \item Les �l�ments strictement sup�rieurs � 14.
  \item Tous les �l�ments sauf les �l�ments en positions 6, 10 et 12.
  \end{enumerate}
  \begin{rep}
    \begin{enumerate}
\item
\begin{Schunk}
\begin{Sinput}
> obs[2]
\end{Sinput}
\end{Schunk}
\item
\begin{Schunk}
\begin{Sinput}
> obs[1:5]
\end{Sinput}
\end{Schunk}
\item
\begin{Schunk}
\begin{Sinput}
> obs[obs > 14]
\end{Sinput}
\end{Schunk}
\item
\begin{Schunk}
\begin{Sinput}
> obs[-c(6, 10, 12)]
\end{Sinput}
\end{Schunk}
    \end{enumerate}
  \end{rep}
\end{exercice}

\begin{exercice}
  \index{indi�age!matrice}
  Soit \code{mat} une matrice $10 \times 7$ obtenue al�atoirement avec
\begin{Schunk}
\begin{Sinput}
> (mat <- matrix(sample(1:100, 70), 7, 10))
\end{Sinput}
\end{Schunk}
  �crire une expression S permettant d'obtenir les �l�ments demand�s
  ci-dessous.
  \begin{enumerate}
  \item L'�l�ment $(4, 3)$ de la matrice.
  \item Le contenu de la sixi�me ligne de la matrice.
  \item Les premi�re et quatri�me colonnes de la matrice
    (simultan�ment).
  \item Les lignes de la matrice dont le premier �l�ment est
    sup�rieur � 50.
  \end{enumerate}
  \begin{rep}
    \begin{enumerate}
\item
\begin{Schunk}
\begin{Sinput}
> mat[4, 3]
\end{Sinput}
\end{Schunk}
\item
\begin{Schunk}
\begin{Sinput}
> mat[6, ]
\end{Sinput}
\end{Schunk}
\item
\begin{Schunk}
\begin{Sinput}
> mat[, c(1, 4)]
\end{Sinput}
\end{Schunk}
\item
\begin{Schunk}
\begin{Sinput}
> mat[mat[, 1] > 50, ]
\end{Sinput}
\end{Schunk}
    \end{enumerate}
  \end{rep}
\end{exercice}

\Closesolutionfile{reponses}

%%% Local Variables:
%%% mode: latex
%%% TeX-master: "introduction_programmation_S"
%%% End:

\include{donnees}
\include{pratiques}
\include{tri}
\include{debogage}
\include{import-export}
\include{paquetages}
%%% Copyright (C) 2017 Vincent Goulet
%%%
%%% Ce fichier fait partie du projet «Programmer avec R»
%%% http://github.com/vigou3/programmer-avec-r
%%%
%%% Cette création est mise à disposition selon le contrat
%%% Attribution-Partage dans les mêmes conditions 4.0
%%% International de Creative Commons.
%%% http://creativecommons.org/licenses/by-sa/4.0/

\chapter{Analyse et contrôle de texte}
\label{chap:texte}

\begin{objectifs}
\item Déterminer si une chaine de caractères correspond ou non à une
  expression régulière donnée.
\item Écrire une expression régulière correspondant à une chaine de
  caractères donnée.
\item Identifier les lignes d'un fichier de texte ayant des propriétés
  communes à l'aide de \code{grep}.
\item Rechercher et remplacer des chaines de caractères ayant des
  propriétés communes à l'aide de \code{sed}.
\item Effectuer les opérations précédentes sur des chaînes de
  caractères divisées en champs logiques avec \code{awk}.
\item Effectuer les opérations de recherche et de remplacement de
  texte à l'aide d'expressions régulières avec les divers outils de R.
\end{objectifs}


%%% Local Variables:
%%% mode: latex
%%% TeX-engine: xetex
%%% TeX-master: "programmer-avec-r"
%%% coding: utf-8
%%% End:

%%% Local Variables:
%%% mode: latex
%%% TeX-master: "programmer-avec-r"
%%% End:

\include{environnement}

\appendix
%%% Copyright (C) 2019 Vincent Goulet
%%%
%%% Ce fichier fait partie du projet
%%% «Programmer avec R»
%%% https://gitlab.com/vigou3/programmer-avec-r
%%%
%%% Cette création est mise à disposition selon le contrat
%%% Attribution-Partage dans les mêmes conditions 4.0
%%% International de Creative Commons.
%%% https://creativecommons.org/licenses/by-sa/4.0/

\chapter{RStudio: une introduction}
\index{RStudio|(}
\label{chap:rstudio}

Un environnement de développement intégré (\emph{integrated
  development environment}, IDE) est un progiciel de productivité
destiné au développement de logiciels ou, plus largement, à la
programmation informatique. Il comprend toujours un éditeur de texte
adapté au langage de programmation visé, un environnement de
compilation ou d'exécution du code et, généralement, des outils de
contrôle de versions, de gestion des projets et de navigation dans le
code source.\footnote{%
  À ce compte, GNU~Emacs constitue un environnement de développement
  intégré.} %

Offert au public depuis 2011, RStudio est un IDE convivial conçu
spécifiquement pour l'analyse de données et le développement de
packages avec R. Il est produit par RStudio~Inc.\ et est offert en
version libre ou commerciale, pour une exécution locale
(\emph{desktop}) ou pour une exécution sur un serveur via un
navigateur web.


\section{Installation}
\label{sec:rstudio:installation}

RStudio est disponible à l'identique pour les plateformes Windows,
macOS et Linux. Pour une utilisation locale sur son poste de travail,
on installera la version libre (\emph{Open Source}) de
\link{https://www.rstudio.com/products/rstudio/download/}{RStudio~Desktop}.


\section{Description sommaire}
\label{sec:rstudio:description}

La fenêtre de RStudio se divise toujours en quatre
sous-fenêtres\footnote{%
  Les sous-fenêtres sont appelées \emph{panes} (en anglais) dans
  l'application.} %
--- sauf au lancement, alors que la sous-fenêtre d'édition de code
source n'est pas visible; voir la \autoref{fig:rstudio:rstudiowindow}.
Dans le sens des aiguilles d'une montre en partant en haut à gauche,
on trouve:
\begin{enumerate}
  \index{RStudio!sous-fenêtres}
\item la sous-fenêtre d'édition de code source, avec un onglet par
  fichier de script;
\item le navigateur d'environnement de travail ou d'historique des
  commandes, selon l'onglet sélectionné;
\item le navigateur de fichiers du projet, de packages, de graphiques,
  etc., selon l'onglet sélectionné;
\item la console --- ou ligne de commande --- R.
\end{enumerate}
Au lancement de l'application, la console R occupe toute la gauche de
la fenêtre jusqu'à ce qu'un fichier de script soit ouvert.

\begin{figure}[t]
  %% Capture d'écran et espace additionnel sous l'image
  \includegraphics{images/rstudiowindow-screenshot}
  \vspace{0.5\TPVertModule}

  \begingroup
  \TPoptions{absolute=false}
  %% Identification de la console
  \begin{textblock}{0.5}(2.65,-0.62)
    \large\faLongArrowAltDown
  \end{textblock}
  \begin{textblock}{2}(2.2,-0.3)
    \footnotesize\sffamily Console R
  \end{textblock}

  %% Identification du navigateur d'environnement
  \begin{textblock}{1}(7.52,-3.65)
    \large\faLongArrowAltRight
  \end{textblock}
  \begin{textblock}{2.2}(7.95,-3.9)
    \footnotesize\sffamily\raggedright Navigateur d'environnement et d'historique
  \end{textblock}

  %% Identification de navigateur de fichiers
  \begin{textblock}{1}(7.52,-1.65)
    \large\faLongArrowAltRight
  \end{textblock}
  \begin{textblock}{2.2}(7.95,-1.9)
    \footnotesize\sffamily\raggedright Navigateur de fichiers, packages, graphiques, etc.
  \end{textblock}
  \endgroup
  \caption[Fenêtre de RStudio sous macOS]{Fenêtre de RStudio et trois
    de ses sous-fenêtres au lancement de l'application sous macOS.
    Sous Windows et Linux, la fenêtre comporte également une barre de
    menu.}
  \label{fig:rstudio:rstudiowindow}
\end{figure}

\begin{itemize}
\item Le navigateur d'environnement de travail est particulièrement
  utile pour voir le contenu, les attributs, le type et la taille de
  chaque objet sauvegardé dans la session R. Il permet également de
  visualiser le contenu des objets en cliquant sur leur nom ou sur
  l'icône de grille à droite de leur nom.
\item Il ne peut y avoir qu'un seul processus R (affiché dans la
  console) actif par fenêtre RStudio. Pour utiliser plusieurs
  processus R simultanément, il faut démarrer autant de copies de
  RStudio.
\item La position des sous-fenêtres dans la grille ne peut être
  modifiée. Par contre, chaque sous-fenêtre peut être redimensionnée.
\item On peut modifier la liste des onglets affichés dans les deux
  navigateurs dans les préférences de l'application; voir la
  \autoref{sec:rstudio:configuration}.
\end{itemize}


\section{Projets}
\index{RStudio!projets}
\label{sec:rstudio:projets}

Il est possible d'utiliser RStudio un peu comme un simple éditeur de
texte.
\begin{itemize}
\item On ouvre les fichiers de scripts un à un, soit à partir du menu
  \code{File|Open file...}, soit à partir de l'onglet \code{Files} du
  navigateur de fichiers.
\item Lorsque nécessaire, on change le répertoire de travail de R à
  partir du menu \code{Session}.
\end{itemize}

Pour faciliter l'organisation de son travail, l'ouverture des fichiers
de script et le lancement d'un processus R dans le bon répertoire de
travail, RStudio propose la notion de \emph{projet}.
\begin{itemize}
\item Un projet RStudio est associé à un répertoire de travail de R
  (\autoref{sec:presentation:workspace}).
\item On crée un nouveau projet à partir du menu \code{Project} à
  l'extrémité droite de la barre d'outils ou à partir du menu
  \code{File|New Project...} On a alors l'option de créer un nouveau
  dossier sur notre poste de travail ou de créer un projet à partir
  d'un dossier existant.
\item Lors de la création d'un projet, RStudio crée dans le dossier
  visé un fichier avec une extension \code{.Rproj} contenant diverses
  informations en lien avec le projet. De plus, le projet est
  immédiatement chargé dans RStudio.
\item L'ouverture d'un projet entraine: le lancement d'une session R
  dont le répertoire de travail est celui du projet; le
  chargement du fichier \code{.RData} (le cas échéant); l'ouverture de
  tous les fichiers de scripts qui étaient ouverts lors de la dernière
  séance de travail.
\item Chaque projet dispose de ses propres réglages. On accède à
  ceux-ci via la commande \code{Project Options...} du menu
  \code{Project} de la barre d'outils.
\item L'utilisation d'un projet permet également d'effectuer la
  gestion des versions du code avec Git ou Subversion
  directement depuis RStudio.
\end{itemize}

On trouvera plus d'information sur les projets dans l'aide en ligne
de RStudio.


\section{Commandes de base}
\label{sec:rstudio:commandes}

Comme l'interface de RStudio respecte les standards modernes, je ne
souligne ici que les commandes particulièrement utiles pour la
manipulation des fichiers de script. On accède rapidement à la liste
des commandes les plus utiles via le menu \code{Help} de
l'application.

Les raccourcis clavier sous, d'une part, Windows et Linux et sous,
d'autre part, macOS, diffèrent légèrement. Je fournis ci-dessous
les deux jeux, séparés par le symbole $\bullet$.

\begin{ttscript}{Ctrl-Shift-S $\bullet$ \cmdkey\shiftkey S}
  \index{RStudio!symbole d'affectation}
\item[\code{Alt+-} $\bullet$ \code{\optkey\,-}] insérer le symbole
  d'affectation \verb*| <- | (pour les Mac avec un clavier
  canadien-français, consulter l'encadré de la
  \autopageref{fig:rstudio:affectation})
\item[\code{Ctrl+Retour} $\bullet$ \code{\cmdkey\,\returnkey}]
  évaluer dans le processus R la ligne sous le curseur ou la région
  sélectionnée, puis déplacer le curseur à la prochaine expression
\item[\code{Ctrl+Shift+S} $\bullet$ \code{\shiftkey\,\cmdkey\,S}]
  évaluer le code du fichier courant en entier dans le processus R
\item[\code{Ctrl+Alt+B} $\bullet$ \code{\optkey\,\cmdkey\,B}]
  évaluer dans le processus R le code source du début du fichier
  jusqu'à la ligne sous le curseur
\item[\code{Ctrl+Alt+E} $\bullet$ \code{\optkey\,\cmdkey\,E}]
  évaluer dans le processus R le code source de la ligne sous le curseur
  jusqu'à la fin du fichier
\item[\code{Ctrl+Alt+F} $\bullet$ \code{\optkey\,\cmdkey\,F}]
  évaluer le code de la fonction courante dans le processus R
\end{ttscript}

À la console --- ou ligne de commande --- R, les raccourcis suivants
sont particulièrement utiles.
\begin{ttscript}{Ctrl-I $\bullet$ \cmdkey\,I}
\item[$\uparrow$ | $\downarrow$] commande
  précédente~|~suivante dans l'historique
\item[\code{Ctrl+}$\uparrow$ $\bullet$ \cmdkey\,$\uparrow$] afficher
  la fenêtre d'historique des commandes
\end{ttscript}

\begin{figure}[t]
  \refstepcounter{dummy}
  \label{fig:rstudio:affectation}
  \begin{titled-frame}{Symbole d'affectation et
      Mac munis d'un clavier canadien-français}
    Le très pratique raccourci clavier {\optkey\,-} servant à insérer
    le symbole d'affectation ne fonctionne pas avec le clavier
    canadien-français sur les Mac. Cette combinaison de touches sert
    déjà à insérer le symbole «\textbar». C'est en fait un doublon
    puisque ce symbole est principalement associé à la touche
    {\optkey\,/}.

    Il est possible de régler le problème en redéfinissant le
    raccourci clavier associé au symbole d'affectation. Je recommande
    d'utiliser la combinaison {\optkey\,=} puisque le caractère lié
    normalement à cette touche («\code{¬}») n'est d'aucune utilité dans
    R.

    \begin{enumerate}
    \item Accéder à la fenêtre d'édition des raccourcis clavier via le
      menu \code{Tools|Modify Keyboard Shortcuts\dots}.
    \item Dans le champ de recherche en haut, taper les lettres de
      «assignment» jusqu'à ce que la commande \code{Insert Assignment
        Operator} apparaisse.
    \item Cliquer sur la case du raccouci clavier (qui devrait être
      \code{Alt+-} par défaut) pour l'activer. Appuyer sur \optkey\,=
      pour utiliser cette combinaison de touches comme raccourci.
    \item Appuyer sur le bouton \code{Apply} pour activer le raccourci
      et pour fermer la fenêtre de configuration.
    \end{enumerate}

    \videobox{\link{https://youtu.be/alGs2KFpBJw}{Configuration du
        raccourci pour \code{<-} dans RStudio pour les Mac avec clavier
        canadien-français}}{%
      La description de cette procédure de configuration est également
      disponible en \link{https://youtu.be/alGs2KFpBJw}{vidéo}.}
  \end{titled-frame}
\end{figure}


\section{Anatomie d'une session de travail (bis)}
\label{sec:rstudio:session}

On reprend ici la description de la session de travail type présentée
à la \autoref{sec:presentation:session}, mais en expliquant comment
compléter chaque étape dans RStudio. Sont intercalés dans les
instructions les raccourcis clavier des commandes RStudio et les accès
par les menus.

\begin{enumerate}
\item Lancer RStudio et ouvrir un nouveau fichier de script ou un
  fichier de script existant.
  \begin{trivlist}
  \item
    \makebox[0.38\linewidth][l]{%
      \colorbox{codebg}{\code{Ctrl+Shift+N} $\bullet$ \code{\shiftkey\,\cmdkey\,N}}}
    \hfill
    \makebox[0.58\linewidth][l]{%
      \colorbox{codebg}{\code{File|New File|R Script...}}}
  \item
    \makebox[0.38\linewidth][l]{%
      \colorbox{codebg}{\code{Ctrl+O} $\bullet$ \code{\cmdkey\,O}}}
    \hfill
    \makebox[0.58\linewidth][l]{%
      \colorbox{codebg}{\code{File|Open File...}}}
  \end{trivlist}
\item C'est une bonne pratique de faire du dossier où se trouve son ou
  ses fichiers de scripts le répertoire de travail de R.
  \begin{trivlist}
  \item
    \colorbox{codebg}{\code{Session|Set Working Directory|To Source File Location}}
  \end{trivlist}
\item Composer le code. Lors de cette étape de programmation, on se
  déplacera souvent du fichier de script à la console R afin
  d'essayer diverses expressions. On exécutera également des parties
  seulement du code se trouvant dans le fichier de script.
  \begin{trivlist}
  \item
    \makebox[0.38\linewidth][l]{%
      \colorbox{codebg}{\code{Ctrl+Retour} $\bullet$ \code{\cmdkey\,\returnkey}}}
    \hfill
    \makebox[0.58\linewidth][l]{%
      \colorbox{codebg}{\code{Code|Run Selected Line(s)}}}
  \item
    \makebox[0.38\linewidth][l]{%
      \colorbox{codebg}{\code{Ctrl+1} $\bullet$ \code{\cmdkey\,1}}}
    \hfill
    \makebox[0.58\linewidth][l]{%
      \colorbox{codebg}{\code{View|Move Focus to Source}}}
  \item
    \makebox[0.38\linewidth][l]{%
      \colorbox{codebg}{\code{Ctrl+2} $\bullet$ \code{\cmdkey\,2}}}
    \hfill
    \makebox[0.58\linewidth][l]{%
      \colorbox{codebg}{\code{View|Move Focus to Console}}}
  \end{trivlist}
\item Sauvegarder le fichier de script.
  \begin{trivlist}
  \item
    \makebox[0.38\linewidth][l]{%
      \colorbox{codebg}{\code{Ctrl+S} $\bullet$ \code{\cmdkey\,S}}}
    \hfill
    \makebox[0.58\linewidth][l]{%
      \colorbox{codebg}{\code{File|Save}}}
  \end{trivlist}
  (S'il s'agit d'un nouveau fichier, s'assurer de terminer son nom par
  \code{.R}.) Le nom du fichier dans l'onglet de la sous-fenêtre passe
  du rouge au noir.
\item Sauvegarder, si désiré, l'espace de travail de R avec
  \code{save.image()}\index{save.image@\code{save.image}}. Cela n'est
  habituellement pas nécessaire à moins que l'espace de travail ne
  contienne des objets importants ou longs à recréer.
  \begin{trivlist}
  \item
    \colorbox{codebg}{\code{Session|Save Workspace As...}}
  \end{trivlist}
\item Quitter RStudio de la manière usuelle. Par défaut, RStudio
  devrait demander si l'on souhaite sauvegarder l'espace de travail de
  R. Nous suggérons de ne pas le faire.

  La \autoref{sec:rstudio:configuration} explique comment configurer
  RStudio afin d'éviter de se faire poser la question à chaque
  fermeture de l'application.
\end{enumerate}


\section{Configuration de l'éditeur}
\index{RStudio!configuration}
\label{sec:rstudio:configuration}

Il est possible de configurer plusieurs facettes de RStudio à partir
d'une interface familière. On accède aux options de configuration par
le menu \code{Tools|Global Options...} sous Windows et Linux, et par le
menu standard \code{RStudio|Preferences} (\code{\cmdkey\,,})
sous macOS.

Je recommande d'effectuer les configurations suivantes.
\begin{enumerate}
\item Dans la catégorie des options de configuration générales, régler
  l'option \code{Save workspace to .RData on exit} à \code{Never}, tel
  qu'illustré à la \autoref{fig:rstudio:rstudio-on-exit}. Avec ce
  réglage, l'espace de travail de R ne sera pas sauvegardé à la
  fermeture de RStudio.
\item Dans la catégorie \code{Code} et la sous-catégorie
  \code{Editing}, régler l'option \code{Tab width} à \code{4} afin
  d'indenter le code de quatre caractères par défaut. (Ce réglage fait
  aussi l'objet de l'\autoref{ex:pratiques:rstudio-tab-width}.)
\item Toujours dans la catégorie \code{Code}, mais dans la
  sous-catégorie \code{Saving}, régler l'option \code{Default text
    encoding} à \code{UTF-8}, tel qu'illustré à la
  \autoref{fig:rstudio:rstudio-utf-8}. Les fichiers de script
  contenant des accents et autres symboles seront ensuite sauvegardés
  dans un format d'encodage standard sur toutes les plateformes
  informatiques.
\end{enumerate}

\begin{figure}
  \centering
  \includegraphics{images/rstudio-on-exit}
  \caption{Réglage de RStudio pour ne pas sauvegarder l'espace
    de travail de R au moment de quitter l'application}
  \label{fig:rstudio:rstudio-on-exit}
\end{figure}

\begin{figure}
  \centering
  \includegraphics{images/rstudio-utf-8}
  \caption{Réglage de RStudio pour enregistrer les fichiers dans
    l'encodage UTF-8}
  \label{fig:rstudio:rstudio-utf-8}
\end{figure}


\section{Aide et documentation}
\label{sec:rstudio:aide}

La documentation de RStudio se trouve entièrement en ligne. On y
accède par le menu \code{Help}. L'onglet \code{Help} du navigateur de
fichiers (sous-fenêtre en bas à droite) offre également une interface
unifiée pour accéder à l'aide de R et à celle de RStudio.

\index{RStudio|)}

%%% Local Variables:
%%% mode: latex
%%% TeX-engine: xetex
%%% TeX-master: "programmer-avec-r"
%%% coding: utf-8
%%% End:

\chapter{GNU Emacs et ESS: la base}
\label{emacs+ess}

Emacs est l'Éditeur de texte des éditeurs de texte. Conçu à l'origine
comme un éditeur pour les programmeurs (avec des modes spéciaux pour
une multitude de langages différents), Emacs est devenu au fil du
temps un environnement logiciel en soi dans lequel on peut réaliser
une foule de tâches différentes: rédiger des documents \LaTeX,
interagir avec R, SAS ou un logiciel de base de données, consulter son
courrier électronique, gérer son calendrier ou même jouer à Tetris!

Cette annexe passe en revue les quelques commandes essentielles à
connaître pour commencer à travailler avec GNU Emacs et le mode ESS.
L'ouvrage de \cite{Cameron:Emacs:2004} constitue une excellente
référence pour l'apprentissage plus poussé de l'éditeur.


\section{Mise en contexte}
\label{emacs+ess:contexte}

Emacs est le logiciel étendard du projet GNU («\emph{GNU is not
  Unix}»), dont le principal commanditaire est la \emph{Free Software
  Foundation} (FSF) à l'origine de tout le mouvement du logiciel
libre.
\begin{itemize}
\item Richard M.\ Stallman, président de la FSF et grand apôtre du
  libre, a écrit la première version de Emacs et il continue à ce jour
  à contribuer au projet.
\item Les origines de Emacs remontent au début des années 1980, une
  époque où les interfaces graphiques n'existaient pas, le parc
  informatique était beaucoup plus hétérogène qu'aujourd'hui (les
  claviers n'étaient pas les mêmes d'une marque d'ordinateur à une
  autre) et les modes de communication entre les ordinateurs
  demeuraient rudimentaires.
\item L'âge vénérable de Emacs transparaît à plusieurs endroits,
  notamment dans la terminologie inhabituelle, les raccourcis clavier
  non conformes aux standards d'aujourd'hui ou la manipulation des
  fenêtres qui ne se fait pas avec une souris.
\end{itemize}

Emacs s'adapte à différentes tâches par l'entremise de \emph{modes}
qui modifient son comportement ou lui ajoutent des fonctionnalités.
L'un de ces modes est ESS (\emph{Emacs Speaks Statistics}).
\begin{itemize}
\item ESS permet d'interagir avec des logiciels statistiques (en
  particulier R, S+ et SAS) directement depuis Emacs.
\item Quelques-uns des développeurs de ESS sont aussi des développeurs
  de R, d'où la grande compatibilité entre les deux logiciels.
\item Lorsque ESS est installé, le mode est activé automatiquement en
  ouvrant dans Emacs un fichier dont le nom se termine par l'extension
  \code{.R}.
\end{itemize}


\section{Installation}
\label{emacs+ess:installation}

GNU Emacs et le mode ESS sont normalement livrés d'office avec toutes
les distributions Linux. Pour les environnements Windows et Mac OS~X,
le plus simple consiste à télécharger et installer les distributions
préparées par le présent auteur. Consulter le site
\begin{quote}
  \url{http://vgoulet.act.ulaval.ca/emacs/}
\end{quote}


\section{Description sommaire}
\label{emacs+ess:description}

Au lancement, Emacs affiche un écran d'information contenant des liens
vers différentes ressources. Cet écran disparaît dès que l'on appuie
sur une touche. La fenêtre Emacs se divise en quatre zone principales
(voir la \autoref{fig:ess:emacswindow}):
\begin{enumerate}
\item tout au haut de la fenêtre (ou de l'écran sous OS~X), on trouve
  l'habituelle barre de menu dont le contenu change selon le mode dans
  lequel se trouve Emacs;
\item l'essentiel de la fenêtre sert à afficher un \emph{buffer}, soit
  le contenu d'un fichier ouvert ou l'invite de commande d'un
  programme externe;
\item la ligne de mode est le séparateur horizontal contenant diverses
  informations sur le fichier ouvert et l'état de Emacs;
\item le \emph{minibuffer} est la région au bas de la fenêtre où l'on
  entre des commandes et reçoit de l'information de Emacs.
\end{enumerate}
Il est possible de séparer la fenêtre Emacs en sous-fenêtres pour
afficher plusieurs \emph{buffers} à la fois. Il y a alors une ligne de
mode pour chaque \emph{buffer}.

\begin{figure}[t]
  %% Capture d'écran
  \includegraphics{emacswindow-screenshot}

  %% Identification du minibuffer
  \begin{textblock}{0.35}(7.45,-0.18)
    \LARGE$\leadsto$
  \end{textblock}
  \begin{textblock}{1.75}(7.95,-0.2)
    \small \emph{Minibuffer}
  \end{textblock}

  %% Identification de la barre de menu
  \begin{textblock}{0.35}(7.45,-6.68)
    \LARGE$\leadsto$
  \end{textblock}
  \begin{textblock}{1.75}(7.95,-6.7)
    \small Barre de menu
  \end{textblock}

  %% Identification du buffer
  \begin{textblock}{0.35}(7.45,-3.48)
    \LARGE$\leadsto$
  \end{textblock}
  \begin{textblock}{1.75}(7.95,-3.5)
    \small \emph{Buffer}
  \end{textblock}

  %% Identification de la mode line
  \begin{textblock}{0.35}(7.45,-0.38)
    \LARGE$\leadsto$
  \end{textblock}
  \begin{textblock}{1.75}(7.95,-0.4)
    \small Ligne de mode
  \end{textblock}
  \caption{Fenêtre GNU~Emacs et ses différentes parties au lancement
    de l'application sous Mac OS~X. Sous Windows et Linux, la barre de
    menu se trouve à l'intérieur de la fenêtre.}
  \label{fig:ess:emacswindow}
\end{figure}


\section{\emph{Emacs-ismes} et \emph{Unix-ismes}}
\label{emacs+ess:ismes}

Emacs possède sa propre terminologie qu'il vaut mieux connaître
lorsque l'on consulte la documentation. De plus, l'éditeur fait
siennes certaines conventions du monde Unix et qui sont moins usitées
sur les plateformes Windows et OS~X.

\begin{itemize}
\item Dans les définitions de raccourcis claviers:
  \begin{itemize}
  \item \code{C} est la touche \code{Contrôle} (\ctlkey);
  \item \code{M} est la touche \code{Meta}, qui correspond à la touche
    \code{Alt} de gauche sur un PC ou la touche \code{Option}
    (\optkey) sur un Mac (toutefois, voir l'encadré ci-dessus);
  \item \code{ESC} est la touche \code{Échap} (\esckey) et
    est équivalente à \code{Meta};
  \item \code{SPC} est la barre d'espacement;
  \item \code{DEL} est la touche \code{Retour arrière} (\delkey) ---
    \emph{et non la touche} \code{Supprimer}.
  \item \code{RET} est la touche \code{Entrée} (\returnkey);
  \end{itemize}
\item Toutes les fonctionnalités de Emacs correspondent à une commande
  pouvant être tapée dans le \emph{minibuffer}. \code{M-x} démarre
  l'invite de commande.
\item Le caractère \,\verb=~=\, représente le dossier vers lequel
  pointe la variable d'environnement \code{\$HOME} (Linux, OS~X) ou
  \code{\%HOME\%} (Windows). C'est le dossier par défaut de Emacs.
\item La barre oblique (\code{/}) est utilisée pour séparer les
  dossiers dans les chemins d'accès aux fichiers, même sous Windows.
\item En général, il est possible d'appuyer sur \code{TAB} dans le
  \emph{minibuffer} pour compléter les noms de fichiers ou de
  commandes.
\end{itemize}

\begin{figure}[t]
  \begin{osx}
    Par défaut sous Mac OS~X, la touche \code{Meta} est assignée à
    \code{Option} (\optkey). Sur les claviers français, cela empêche
    d'accéder à certains caractères spéciaux tels que [, ], \{ ou \}.

    Une solution consiste à plutôt assigner la touche \code{Meta} à
    \code{Commande} (\cmdkey). Cela bloque alors l'accès à certains
    raccourcis Mac, mais la situation est moins critique ainsi.

    Pour assigner la touche \code{Meta} à \code{Commande} (\cmdkey) et
    laisser la touche \code{Option} (\optkey) jouer son rôle usuel, il
    suffit d'insérer les lignes suivantes dans son fichier de
    configuration \code{.emacs} (voir la
    \autoref{emacs+ess:configuration}):
\begin{verbatim}
;;; ====================================
;;;  Assigner la touche Meta à Commande
;;;  et laisser Option être Option
;;; ====================================
(setq-default ns-command-modifier 'meta)
(setq-default ns-option-modifier 'none)
\end{verbatim}
  \end{osx}
\end{figure}



\section{Commandes de base}
\label{emacs+ess:commandes}

Emacs comporte une pléthore de commandes, il serait donc futile de
tenter d'en faire une liste exhaustive ici. Nous nous contenterons de
mentionner les commandes les plus importantes regroupées par tâche.

Pour débuter, il est utile de suivre le Tour guidé de Emacs\footnote{%
  \url{http://www.gnu.org/software/emacs/tour/} ou cliquer sur le lien
  dans l'écran d'accueil.} %
et de lire le tutoriel de Emacs, que l'on démarre avec \code{C-h t}.

\subsection{Les essentielles}
\label{emacs+ess:commandes:essentielles}

\begin{ttscript}{M-x}
\item[\emacs{M-x}] démarrer l'invite de commande
\item[\emacs{C-g}] bouton de panique: annuler, quitter! Presser plus
  d'une fois au besoin.
\end{ttscript}

\subsection{Manipulation de fichiers}
\label{emacs+ess:commandes:fichiers}

Entre parenthèses, le nom de la commande Emacs correspondante. On peut
entrer cette commande dans le \emph{minibuffer} au lieu d'utiliser le
raccourci clavier.

\begin{important}
  On remarquera qu'il n'existe pas de commande «nouveau fichier» dans
  Emacs. Pour créer un nouveau fichier, il suffit d'ouvrir un fichier
  n'existant pas.\index{Emacs!nouveau fichier}
\end{important}

\begin{ttscript}{C-x C-w}
  \raggedright
\item[\emacs{C-x C-f}] ouvrir un fichier (\code{find-file})
\item[\emacs{C-x C-s}] sauvegarder
  (\code{save-buffer})\index{Emacs!sauvegarder}
\item[\emacs{C-x C-w}] sauvegarder sous
  (\code{write-file})\index{Emacs!sauvegarder sous}
\item[\emacs{C-x k}] fermer un fichier (\code{kill-buffer})
  \\[\baselineskip]
\item[\emacs{C-\_}] annuler (pratiquement illimité); aussi
  \emacs{C-x u} (\code{undo})
  \\[\baselineskip]
\item[\emacs{C-s}] recherche incrémentale avant
  (\code{isearch-forward})
\item[\emacs{C-r}] Recherche incrémentale arrière
  (\code{isearch-backward})
\item[\emacs{M-\%}] rechercher et remplacer
  (\code{query-replace})\index{Emacs!rechercher et remplacer}
\end{ttscript}


\subsection{Déplacements simples du curseur}
\index{Emacs!déplacement}
\label{emacs+ess:commandes:deplacement}

\begin{ttscript}{C-DEL | C-w}
  \raggedright
\item[\emacs{C-b} | \emacs{C-f}] déplacer d'un caractère vers l'arrière~|~l'avant
  (\code{backward-char} | \code{forward-char})
\item[\emacs{C-a} | \emacs{C-e}] aller au début~|~fin de la ligne
  (\code{move-beginning-of-line} | \code{move-end-of-line})
\item[\emacs{C-p} | \emacs{C-n}] aller à la ligne précédente~|~suivante
  (\code{previous-line} | \code{next-line})
\item[\emacs{M-<} | \emacs{M->}] aller au début~|~fin du fichier
  (\code{beginning-of-buffer} | \code{end-of-buffer}) \\[\baselineskip]
\item[\emacs{DEL} | \emacs{C-d}] effacer le caractère à
  gauche~|~droite du curseur (\code{delete-backward-char} | \code{delete-char})
\item[\emacs{M-DEL} | \emacs{M-d}] effacer le mot à gauche~|~droite
  du curseur (\code{backward-kill-word} | \code{kill-word})
\item[\emacs{C-k}] supprimer jusqu'à la fin de la ligne (\code{kill-line})
\end{ttscript}

\begin{figure}[h]
  \begin{osx}
    Plusieurs des raccourcis clavier de Emacs composés avec la touche
    \code{Contrôle} (\ctlkey) sont valides sous Mac OS~X. Par exemple,
    \ctlkey\,\code{A} et \ctlkey\,\code{E} déplacent le curseur au
    début et à la fin de la ligne dans les champs texte.
  \end{osx}
\end{figure}


\subsection{Sélection de texte, copier, coller, couper}
\index{Emacs!sélection}
\label{emacs+ess:commandes:selection}

\begin{ttscript}{C-x C-w}
  \raggedright
\item[\emacs{C-SPC}] débute la sélection (\code{set-mark-command})
\item[\emacs{C-w}] couper la sélection (\code{kill-region})
\item[\emacs{M-w}] copier la sélection (\code{kill-ring-save})
\item[\emacs{C-y}] coller (\code{yank})
\item[\emacs{M-y}] remplacer le dernier texte collé par la
  sélection précédente (\code{yank-pop})
\end{ttscript}

\begin{itemize}
\item Il est possible d'utiliser les raccourcis clavier usuels de
  Windows (\code{C-c}, \code{C-x}, \code{C-v}) et OS~X
  (\cmdkey\,\code{C}, \cmdkey\,\code{X}, \cmdkey\,\code{V}) en
  activant le mode CUA dans le menu \code{Options}.
\item On peut copier-coller directement avec la souris dans Windows en
  sélectionnant du texte puis en appuyant sur le bouton central (ou la
  molette) à l'endroit souhaité pour y copier le texte.
\end{itemize}


\subsection{Manipulation de fenêtres}
\label{emacs+ess:commandes:fenetres}

\begin{ttscript}{C-x C-w}
\item[\emacs{C-x b}] changer de \emph{buffer}
  (\code{switch-buffer})
\item[\emacs{C-x 2}] séparer l'écran en deux fenêtres
  (\code{split-window-vertically})
\item[\emacs{C-x 1}] conserver uniquement la fenêtre courante
  (\code{delete-other-windows})
\item[\emacs{C-x 0}] fermer la fenêtre courante
  (\code{delete-window})
\item[\emacs{C-x o}] aller vers une autre fenêtre lorsqu'il y en a
  plus d'une (\code{other-window})
\end{ttscript}

\subsection{Manipulation de fichiers de script dans le mode ESS}
\label{emacs+ess:commandes:script}

Dans les versions récentes de ESS (à partir de 12.09 environ), la
manipulation des fichiers de script a été passablement simplifiée par
l'introduction de fonctions «intelligentes» qui s'adaptent à la
situation. Les deux principales commandes à connaître sont les
suivantes:

\begin{ttscript}{C-x C-w}
  \raggedright
\item[\ess{C-RET}] évaluer dans le processus R la ligne sous le
  curseur ou la région sélectionnée, puis déplacer le curseur à la
  prochaine expression (\code{ess-eval-region-or-line-and-step})
\item[\ess{C-c C-c}] évaluer dans le processus R la région sélectionnée, la
  fonction ou le paragraphe dans lequel se trouve le curseur, puis
  déplacer le curseur à la prochaine expression
  (\code{ess-eval-region-or-function-or-paragraph-and-step})
\end{ttscript}
Quelques autres fonctions utiles:
\begin{ttscript}{C-x C-w}
  \raggedright
\item[\ess{C-c C-v}] aide sur une commande R
  (\code{ess-display-help-on-object})
\item[\ess{C-c C-l}] évaluer le code du fichier courant en entier dans
  le processus R (\code{ess-load-file})
\item[\ess{C-c C-n}] évaluer la ligne sous le curseur dans le
  processus R, puis déplacer le curseur à la prochaine expression
  (\code{ess-eval-line-and-step})
\item[\ess{C-c C-r}] évaluer la région sélectionnée dans le processus
  R (\code{ess-eval-region})
\item[\ess{C-c C-f}] évaluer le code de la fonction courante dans
  le processus R (\code{ess-eval-function})
\end{ttscript}

\subsection{Interaction avec l'invite de commande R}
\label{emacs+ess:commandes:invite}

\begin{ttscript}{C-x | C-w}
  \raggedright
\item[\ess{M-p} | \ess{M-n}] commande précédente~|~suivante
  dans l'historique (\code{previous-matching-history-from-input},
  \code{next-matching-history-from-input})
\item[\ess{C-c C-e}] replacer la dernière ligne au bas de la fenêtre
  (\code{comint-show-maximum-output})
\item[\ess{M-h}] sélectionner le résultat de la dernière commande
  (\code{mark-paragraph})
\item[\ess{C-c C-o}] effacer le résultat de la dernière commande
  (\code{comint-delete-output})
\item[\ess{C-c C-v}] aide sur une commande R
  (\code{ess-display-help-on-object})
\item[\ess{C-c C-q}] terminer le processus R (\code{ess-quit})
\end{ttscript}

\subsection{Consultation des rubriques d'aide de R}
\label{emacs+ess:commandes:aide}

\begin{ttscript}{m, m}
  \raggedright
\item[\ess{p} | \ess{n}] aller à la section précédente~|~suivante de
  la rubrique (\code{ess-skip-to-previous-section} |
  \code{ess-skip-to-next-section})
\item[\ess{s a}] aller à la section de la liste des arguments \emph{(Arguments)}
\item[\ess{s D}] aller à la section des détails sur la fonction \emph{(Details)}
\item[\ess{s v}] aller à la section sur la valeur retournée par la
  fonction \emph{(Value)}
\item[\ess{s s}] aller à la section des fonctions apparentée \emph{(See Also)}
\item[\ess{s e}] aller à la section des exemples \emph{(Examples)}
\item[\ess{l}] évaluer la ligne sous le curseur; pratique pour
  exécuter les exemples (\code{ess-eval-line-and-step})
\item[\ess{r}] évaluer la région sélectionnée (\code{ess-eval-region})
\item[\ess{h}] ouvrir une nouvelle rubrique d'aide, par défaut pour le
  mot se trouvant sous le curseur (\code{ess-display-help-on-object})
\item[\ess{q}] retourner au processus ESS en laissant la rubrique
  d'aide visible (\code{ess-switch-to-end-of-ESS})
\item[\ess{x}] fermer la rubrique d'aide et retourner au processus ESS
  (\code{ess-kill-buffer-and-go})
\end{ttscript}



\section{Anatomie d'une session de travail (bis)}
\label{emacs+ess:session}

On reprend ici les étapes d'une \capsule{session de travail} type
présentées à la \autoref{presentation:session}, mais en expliquant
comment compléter chacune dans Emacs avec le mode ESS.

\begin{enumerate}
\item Lancer Emacs et ouvrir un fichier de script avec
  \begin{quote}
    \code{C-x C-f}
  \end{quote}
  ou avec le menu
  \begin{quote}
    \code{File|Open file...}
  \end{quote}
  En spécifiant un nom de fichier qui n'existe pas déjà, on se trouve
  à créer un nouveau fichier de script. S'assurer de terminer le nom
  des nouveaux fichiers par \code{.R} pour que Emacs reconnaisse
  automatiquement qu'il s'agit de fichiers de script R.
\item Démarrer un processus R à l'intérieur même de Emacs avec
  \begin{quote}
    \code{M-x R }\returnkey
  \end{quote}
  Emacs demandera alors de spécifier de répertoire de travail
  (\emph{starting data directory}). Accepter la valeur par défaut, par
  exemple
  \begin{quote}
    \verb=~/ =\returnkey
  \end{quote}
  ou indiquer un autre dossier. Un éventuel message de Emacs à l'effet
  que le fichier \code{.Rhistory} n'a pas été trouvé est sans
  conséquence et peut être ignoré.
\item Composer le code. Lors de cette étape, on se déplacera souvent
  du fichier de script à la ligne de commande afin d'essayer diverses
  expressions. On exécutera également des parties seulement du code se
  trouvant dans le fichier de script. Les commandes les plus utilisées
  sont alors
  \begin{quote}
    \ess{C-RET}\ pour exécuter une ligne du fichier de script; \\
    \ess{C-c C-c}\ pour exécuter un paragraphe du fichier de script; \\
    \emacs{C-x o}\ pour se déplacer d'une fenêtre à l'autre; \\
    \ess{C-c C-e}\ pour replacer la ligne de commande au bas de la
    fenêtre.
  \end{quote}
\item Sauvegarder le fichier de script:
  \begin{quote}
    \emacs{C-x C-s}
  \end{quote}
  Les quatrième et cinquième caractères de la ligne de mode changent
  de \,\verb|**|\, à \,\verb|--|.
\item Sauvegarder si désiré l'espace de travail de R avec
  \code{save.image()}\indexfonction{save.image}. On le répète, cela
  n'est habituellement pas nécessaire à moins que l'espace de travail
  ne contienne des objets importants ou longs à recréer.
\item Quitter le processus R avec
  \begin{quote}
    \ess{C-c C-q}
  \end{quote}
  Cette commande ESS se chargera de fermer tous les fichiers associés
  au processus R. On peut ensuite quitter Emacs en fermant
  l'application de la manière usuelle.
\end{enumerate}



\section{Configuration de l'éditeur}
\index{Emacs!configuration}
\label{emacs+ess:configuration}

Une des grandes forces de Emacs est qu'à peu près chacune de ses
facettes est configurable: couleurs, polices de caractère, raccourcis
clavier, etc.

\begin{itemize}
\item Un utilisateur place ses commandes de configuration dans un
  fichier nommé \code{.emacs} (le point est important!) que Emacs lit
  au démarrage.
\item Le fichier \code{.emacs} doit se trouver dans le dossier
  \,\verb=~/=, c'est-à-dire dans le dossier de départ de l'utilisateur
  sous Linux et OS~X, et dans le dossier référencé par la variable
  d'environnement \code{\%HOME\%} sous Windows.
\end{itemize}


\section{Aide et documentation}
\label{emacs+ess:aide}

Emacs possède son propre système d'aide très exhaustif, mais dont la
navigation est peu intuitive selon les standards d'aujourd'hui.
Consulter le menu \code{Help}.

Autrement, on trouvera les manuels de Emacs et de ESS en divers
formats dans les sites respectifs des deux projets:
\begin{quote}
  \url{http://www.gnu.org/software/emacs} \\
  \url{http://ess.r-project.org}
\end{quote}

Enfin, si le désespoir vous prend au cours d'une séance de codage
intensive, vous pouvez toujours consulter le psychothérapeute Emacs.
On le trouve, bien entendu, dans le menu \code{Help}!

%%% Local Variables:
%%% mode: latex
%%% TeX-master: "introduction_programmation_r"
%%% coding: utf-8
%%% End:

%%% Copyright (C) 2018 Vincent Goulet
%%%
%%% Ce fichier fait partie du projet
%%% «Programmer avec R»
%%% https://gitlab.com/vigou3/programmer-avec-r
%%%
%%% Cette création est mise à disposition selon le contrat
%%% Attribution-Partage dans les mêmes conditions 4.0
%%% International de Creative Commons.
%%% https://creativecommons.org/licenses/by-sa/4.0/

\chapter{Solutions des exercices}
\label{chap:solutions}

\begingroup

%% Environnement Schunk simplifié pour l'affichage des réponses
\renewenvironment{Schunk}{%
  \setlength{\topsep}{0pt}
  \colorlet{shadecolor}{codebg}
  \begin{snugshade*}}%
  {\end{snugshade*}}

\input{solutions-informatique}
\input{solutions-premiers}
\input{solutions-fonctions}
\input{solutions-implementation}
\input{solutions-algorithmique}
\input{solutions-donnees}
\input{solutions-application}
\input{solutions-cli}
\input{solutions-texte}
\input{solutions-git}
\input{solutions-debogage}

\endgroup

%%% Local Variables:
%%% mode: latex
%%% TeX-engine: xetex
%%% TeX-master: "programmer-avec-r"
%%% coding: utf-8
%%% End:


\bibliography{r,stat,informatique,vg}

\cleardoublepage
\chaptermark{Index}
\printindex

\cleartoverso
\pagestyle{empty}
%%% Copyright (C) 2019 Vincent Goulet
%%%
%%% Ce fichier fait partie du projet
%%% «Programmer avec R»
%%% https://gitlab.com/vigou3/programmer-avec-r
%%%
%%% Cette création est mise à disposition sous licence
%%% Attribution-Partage dans les mêmes conditions 4.0
%%% International de Creative Commons.
%%% https://creativecommons.org/licenses/by-sa/4.0/

\vspace*{\fill}

\begingroup
\calccentering{\unitlength}
\begin{adjustwidth*}{\unitlength}{-\unitlength}
  \begin{flushleft}
    \small %
    Ce document a été produit avec le système de mise en page
    {\XeLaTeX}. Le texte principal est composé en Lucida Bright~OT
    11~points, les mathématiques en Lucida Bright Math~OT, le code
    informatique en Lucida Grande Mono~DK et les titres en Fira~Sans.
    Des icônes proviennent de la police Font~Awesome. Les graphiques
    ont été réalisés avec R.
  \end{flushleft}
\end{adjustwidth*}
\endgroup
\vfill


\cleartoverso
%%% Copyright (C) 2019 Vincent Goulet
%%%
%%% Ce fichier fait partie du projet
%%% «Programmer avec R»
%%% https://gitlab.com/vigou3/programmer-avec-r
%%%
%%% Cette création est mise à disposition sous licence
%%% Attribution-Partage dans les mêmes conditions 4.0
%%% International de Creative Commons.
%%% https://creativecommons.org/licenses/by-sa/4.0/

\begingroup
\TPGrid{8}{64}
\textblockorigin{0mm}{0mm}
\setlength{\parindent}{0mm}
\setlength{\bandeorwidth}{12.8125\TPVertModule} % == 4\logoheight
\setlength{\banderougewidth}{\paperwidth}
\addtolength{\banderougewidth}{-\bandeorwidth}
\setlength{\bandeorheight}{\TPVertModule}
\setlength{\banderougeheight}{\TPVertModule}

%% bandeau identitaire arrière
\begin{textblock*}{\paperwidth}[0,1](0mm,56\TPVertModule)
  \textcolor{or}{\rule{\bandeorwidth}{\TPVertModule}}%      % bande or
  \textcolor{rouge}{\rule{\banderougewidth}{\TPVertModule}} % bande rouge
\end{textblock*}
\endgroup

%%% Local Variables:
%%% mode: latex
%%% TeX-master: "programmer-avec-r"
%%% TeX-engine: xetex
%%% End:


\end{document}

%%% Local Variables:
%%% mode: latex
%%% TeX-engine: xetex
%%% TeX-master: t
%%% coding: utf-8
%%% End:
